% Environment well-being and health — Anglais Tle Tech — Cours signé OnBuch+
% Programme officiel MINESEC. Compiler avec : tectonic N03-environment-well-being-and-health.tex
\newcommand{\DOCMATIERE}{Anglais}
\newcommand{\DOCNIVEAU}{Tle Tech}
\newcommand{\DOCMODULE}{Anglais – classe de Terminale technique}
\newcommand{\DOCLECON}{N03}
\newcommand{\DOCTITRE}{Environment well-being and health}
% OnBuch+ — préambule « cours signés » ENRICHI (compilé avec tectonic / XeLaTeX)
% Système visuel « Pop » d'OnBuch+ : fond crème (jamais blanc), bordures encre
% épaisses, ombres « dures » décalées, titres Archivo Black en capitales.
% Version MATHÉMATIQUES : graphes (pgfplots/tikz), boîtes pédagogiques
% (définition, propriété, méthode, attention, le savais-tu, exercice résolu,
% exercice, TP, bilan), page de garde et signature OnBuch+ ancrée en bas.
%
% Macros attendues dans main.tex AVANT \input{preamble} :
%   \DOCMATIERE  \DOCNIVEAU  \DOCMODULE  \DOCLECON (numéro)  \DOCTITRE
\documentclass[11pt]{article}

\usepackage{fontspec}
\usepackage[french]{babel}
\usepackage[a4paper,margin=2cm,top=3.4cm,bottom=2.8cm,headheight=1.4cm,footskip=1.3cm]{geometry}
\usepackage{amsmath,amssymb,amsfonts}
\usepackage{mathtools} % \xmapsto, \coloneqq... souvent utilisés par les modèles
\usepackage{cancel} % simplifications barrées (bilans de forces)
% \abs{z} (module, valeur absolue) et \norm{u} (norme), utilisés par les modèles
\providecommand{\abs}{}\let\abs\relax\DeclarePairedDelimiter{\abs}{\lvert}{\rvert}
\providecommand{\norm}{}\let\norm\relax\DeclarePairedDelimiter{\norm}{\lVert}{\rVert}
\usepackage[table]{xcolor} % [table] : fournit aussi \rowcolors (alternance de lignes)
\usepackage{pagecolor}
\usepackage{tikz}
\usetikzlibrary{arrows.meta,positioning,calc,shapes.geometric,shapes.misc,shapes.arrows,decorations.pathmorphing,decorations.pathreplacing,decorations.markings,patterns,shadows,mindmap,trees,fit,backgrounds,matrix,intersections,angles,quotes}
\usepackage{pgfplots}
\usepackage[version=4]{mhchem} % équations nucléaires \ce{^{235}_{92}U}
\usepackage[european,siunitx]{circuitikz} % schémas électriques
\pgfplotsset{compat=1.18}
\usepgfplotslibrary{fillbetween,groupplots} % aires entre courbes (calcul intégral)
\usepackage{enumitem}
\usepackage{fancyhdr}
% [most] charge toutes les bibliothèques tcolorbox (listings, minted, raster,
% documentation...) et gonfle inutilement la mémoire TeX sur un document long
% (~300 boîtes) : on ne charge que ce qui est réellement utilisé.
\usepackage[skins,breakable]{tcolorbox}
\usepackage{graphicx}
\usepackage{colortbl}
\usepackage{booktabs}
\usepackage{tabularx}
\usepackage{multirow}
\usepackage{diagbox} % case d'en-tête diagonale des tableaux à double entrée
% Colonnes extensibles centrée / gauche / droite (C, L, R), souvent utilisées par les modèles
\newcolumntype{C}{>{\centering\arraybackslash}X}
\newcolumntype{L}{>{\raggedright\arraybackslash}X}
\newcolumntype{R}{>{\raggedleft\arraybackslash}X}
\usepackage{array}
\usepackage{multicol}
\usepackage{siunitx}
% parse-numbers=false : accepte aussi \SI{2,5 \times 10^{-3}}{...}
\sisetup{locale=FR,output-decimal-marker={,},group-minimum-digits=4,parse-numbers=false}
\DeclareSIUnit\ohm{\ensuremath{\symup{\Omega}}} % U+2126 absent de la police
\DeclareSIUnit\division{div} % réglages d'oscilloscope (ms/div, V/div)
\DeclareSIUnit\tr{tr} % tours (vitesse de rotation, ex. \SI{1500}{\tr\per\minute})
\DeclareSIUnit\molar{M} % concentration molaire (ex. \SI{3}{\molar}, \SI{1}{\micro\molar})
\DeclareSIUnit\an{an} % année (le modèle confond parfois avec \year, un compteur TeX) — ex. \SI{5}{\kilo\meter\per\an}
\DeclareSIUnit\FCFA{FCFA} % franc CFA (ex. \SI{500}{\FCFA\per\kilo\gram})
\DeclareSIUnit\degreeBrix{\textdegree Brix} % teneur en sucre d'un sirop/jus (réfractométrie)
\DeclareSIUnit\month{mois} % le modèle utilise parfois \month, un compteur TeX, comme unité de durée
\DeclareSIUnit\microliter{\micro\liter} % le modèle écrit parfois \microliter en un seul token
\DeclareSIUnit\million{millions} % dénombrement (ex. \SI{15}{\million\per\mL} spermatozoïdes)
\usepackage{qrcode}
% \degree : commande fréquemment utilisée par les modèles pour un angle ou une
% température en dehors de \SI{}{} (non fournie par les paquets chargés).
\providecommand{\degree}{^{\circ}}
% \qed (fin de démonstration), utilisé par les modèles sans amsthm
\providecommand{\qed}{\hfill$\square$}
% \barre : double barre des valeurs interdites dans un tableau de variations (tkz-tab)
\providecommand{\barre}{\ensuremath{\|}}
% environnement proof (amsthm non chargé) utilisé par les modèles
\ifdefined\proof\else\newenvironment{proof}[1][Démonstration]{\par\noindent\textit{#1.}\ }{\qed\par}\fi
\usepackage{lastpage}
\usepackage{hyperref}
\hypersetup{hidelinks}

% ============================================================
% Palette « Pop » OnBuch+ — mêmes hex que la classe P (plus_ui.dart)
% ============================================================
\definecolor{poporange}{HTML}{F59321}
\definecolor{popdark}{HTML}{E07A0C}
\definecolor{popcream}{HTML}{FFF6E8}
\definecolor{popink}{HTML}{1C1714}
\definecolor{poppurple}{HTML}{7A5AE0}
\definecolor{popgreen}{HTML}{2E9E6A}
\definecolor{poppink}{HTML}{E0457B}
\definecolor{popblue}{HTML}{2F7DE1}
\definecolor{popgold}{HTML}{C98A12}
\definecolor{popmuted}{HTML}{8A7F76}
\definecolor{popline}{HTML}{EADFD2}
\definecolor{popgray}{HTML}{8A7F76}
\colorlet{popgrey}{popgray}
% Alias tolérants (le modèle nomme parfois les couleurs différemment)
\colorlet{popwhite}{white}
\colorlet{popblack}{popink}
\colorlet{popred}{poppink}
\colorlet{popyellow}{popgold}
% Teintes claires (fonds de tableaux/encadrés) — le modèle nomme parfois
% les couleurs avec un suffixe L (ex. popblueL) sans les avoir définies.
\colorlet{poporangeL}{poporange!20}
\colorlet{popdarkL}{popdark!20}
\colorlet{poppurpleL}{poppurple!20}
\colorlet{popgreenL}{popgreen!20}
\colorlet{poppinkL}{poppink!20}
\colorlet{popblueL}{popblue!20}
\colorlet{popgoldL}{popgold!20}
\colorlet{popredL}{popred!20}
\colorlet{popgrayL}{popgray!20}
% Teintes claires pour fonds de tableaux / zones
\colorlet{popblueL}{popblue!10!white}
\colorlet{poporangeL}{poporange!14!white}
\colorlet{popgreenL}{popgreen!12!white}
\colorlet{poppurpleL}{poppurple!12!white}
\colorlet{poppinkL}{poppink!12!white}
\colorlet{popgoldL}{popgold!14!white}

\pagecolor{popcream}

% ============================================================
% Typographie
% ============================================================
\defaultfontfeatures{Ligatures=TeX}
\setmainfont{PlusJakartaSans-Medium.ttf}[Path=../../../fonts/,BoldFont=PlusJakartaSans-Bold.ttf]
% Maths en sans-serif (Fira Math) avec chiffres et lettres droites de Plus
% Jakarta, pour que les formules se fondent dans le texte.
\usepackage[math-style=ISO,bold-style=ISO]{unicode-math}
\setmathfont{FiraMath-Regular.otf}[Path=../../../fonts/]
\setmathfont{PlusJakartaSans-Medium.ttf}[Path=../../../fonts/,range={up/num,up/{Latin,latin},\mathup}]
\usepackage{icomma} % virgule décimale sans espace en mode math
% Caractères mathématiques Unicode absents des polices de texte (Plus Jakarta,
% Archivo Black) : rendus par leur équivalent mathématique, en texte comme en
% formule — sinon ils s'affichent en carré vide (ex. « Arithmétique dans ℤ »).
\usepackage{newunicodechar}
\newunicodechar{ℝ}{\ensuremath{\mathbb{R}}}
\newunicodechar{ℤ}{\ensuremath{\mathbb{Z}}}
\newunicodechar{ℂ}{\ensuremath{\mathbb{C}}}
\newunicodechar{ℕ}{\ensuremath{\mathbb{N}}}
\newunicodechar{ℚ}{\ensuremath{\mathbb{Q}}}
\newunicodechar{𝔻}{\ensuremath{\mathbb{D}}}
\newunicodechar{α}{\ensuremath{\alpha}}
\newunicodechar{β}{\ensuremath{\beta}}
\newunicodechar{γ}{\ensuremath{\gamma}}
\newunicodechar{δ}{\ensuremath{\delta}}
\newunicodechar{ε}{\ensuremath{\varepsilon}}
\newunicodechar{θ}{\ensuremath{\theta}}
\newunicodechar{λ}{\ensuremath{\lambda}}
\newunicodechar{μ}{\ensuremath{\mu}}
\newunicodechar{σ}{\ensuremath{\sigma}}
\newunicodechar{τ}{\ensuremath{\tau}}
\newunicodechar{φ}{\ensuremath{\varphi}}
\newunicodechar{ψ}{\ensuremath{\psi}}
\newunicodechar{ω}{\ensuremath{\omega}}
\newunicodechar{Σ}{\ensuremath{\Sigma}}
% majuscules produites par \MakeUppercase dans les titres (α-aminés -> Α)
\newunicodechar{Α}{\ensuremath{\alpha}}
\newunicodechar{Β}{\ensuremath{\beta}}
\newunicodechar{Γ}{\ensuremath{\Gamma}}
\newunicodechar{Δ}{\ensuremath{\Delta}}
\newunicodechar{Ε}{\ensuremath{\varepsilon}}
\newunicodechar{Θ}{\ensuremath{\Theta}}
\newunicodechar{Λ}{\ensuremath{\Lambda}}
\newunicodechar{Μ}{\ensuremath{\mu}}
\newunicodechar{Τ}{\ensuremath{\tau}}
\newunicodechar{Φ}{\ensuremath{\Phi}}
\newunicodechar{Ψ}{\ensuremath{\Psi}}
\newunicodechar{Ω}{\ensuremath{\Omega}}
\newunicodechar{↦}{\ensuremath{\mapsto}}
\newunicodechar{∈}{\ensuremath{\in}}
\newunicodechar{∉}{\ensuremath{\notin}}
\newunicodechar{∧}{\ensuremath{\wedge}}
\newunicodechar{⇔}{\ensuremath{\Leftrightarrow}}
\newunicodechar{⟺}{\ensuremath{\Longleftrightarrow}}
\newunicodechar{⇒}{\ensuremath{\Rightarrow}}
\newunicodechar{⟹}{\ensuremath{\Longrightarrow}}
\newunicodechar{∘}{\ensuremath{\circ}}
\newunicodechar{∪}{\ensuremath{\cup}}
\newunicodechar{∩}{\ensuremath{\cap}}
\newunicodechar{≡}{\ensuremath{\equiv}}
\newunicodechar{⊂}{\ensuremath{\subset}}
\newunicodechar{⊥}{\ensuremath{\perp}}
\newunicodechar{∃}{\ensuremath{\exists}}
\newunicodechar{∀}{\ensuremath{\forall}}
\newunicodechar{∛}{\ensuremath{\sqrt[3]{\,}}}
\newunicodechar{∜}{\ensuremath{\sqrt[4]{\,}}}
\newunicodechar{⃗}{\ensuremath{{}^{\to}}}
\newunicodechar{ⁿ}{\ifmmode^{n}\else\textsuperscript{\ensuremath{n}}\fi}
\newunicodechar{ˣ}{\ifmmode^{x}\else\textsuperscript{\ensuremath{x}}\fi}
\newunicodechar{ᵇ}{\ifmmode^{b}\else\textsuperscript{\ensuremath{b}}\fi}
\newunicodechar{ʳ}{\ifmmode^{r}\else\textsuperscript{\ensuremath{r}}\fi}
\newunicodechar{ᵘ}{\ifmmode^{u}\else\textsuperscript{\ensuremath{u}}\fi}
\newunicodechar{ʸ}{\ifmmode^{y}\else\textsuperscript{\ensuremath{y}}\fi}
\newunicodechar{ᵃ}{\ifmmode^{a}\else\textsuperscript{\ensuremath{a}}\fi}
\newunicodechar{ᵉ}{\ifmmode^{e}\else\textsuperscript{\ensuremath{e}}\fi}
\newunicodechar{ᵏ}{\ifmmode^{k}\else\textsuperscript{\ensuremath{k}}\fi}
\newunicodechar{ᵖ}{\ifmmode^{p}\else\textsuperscript{\ensuremath{p}}\fi}
\newunicodechar{ᴺ}{\ifmmode^{N}\else\textsuperscript{\ensuremath{N}}\fi}
\newunicodechar{ᶜ}{\ifmmode^{c}\else\textsuperscript{\ensuremath{c}}\fi}
\newunicodechar{ʰ}{\ifmmode^{h}\else\textsuperscript{\ensuremath{h}}\fi}
\newunicodechar{ᵠ}{\ifmmode^{q}\else\textsuperscript{\ensuremath{q}}\fi}
\newunicodechar{ₙ}{\ifmmode_{n}\else\textsubscript{\ensuremath{n}}\fi}
\newunicodechar{ₐ}{\ifmmode_{a}\else\textsubscript{\ensuremath{a}}\fi}
\newunicodechar{ᵢ}{\ifmmode_{i}\else\textsubscript{\ensuremath{i}}\fi}
\newunicodechar{ᵣ}{\ifmmode_{r}\else\textsubscript{\ensuremath{r}}\fi}
\newunicodechar{ₖ}{\ifmmode_{k}\else\textsubscript{\ensuremath{k}}\fi}
\newunicodechar{ᵦ}{\ifmmode_{\beta}\else\textsubscript{\ensuremath{\beta}}\fi}
\newunicodechar{ₓ}{\ifmmode_{x}\else\textsubscript{\ensuremath{x}}\fi}
\newfontfamily\popdisplay{ArchivoBlack-Regular.ttf}[Path=../../../fonts/]
\newfontfamily\poplogo{SpaceGrotesk-Bold.ttf}[Path=../../../fonts/]
\newfontfamily\popsemi{PlusJakartaSans-SemiBold.ttf}[Path=../../../fonts/]
\newfontfamily\popxbold{PlusJakartaSans-ExtraBold.ttf}[Path=../../../fonts/]
\newfontfamily\popxboldls{PlusJakartaSans-ExtraBold.ttf}[Path=../../../fonts/,LetterSpace=3.0]

\setlist[enumerate]{leftmargin=1.7em,itemsep=5pt,topsep=4pt}
\setlist[itemize]{leftmargin=1.7em,itemsep=4pt,topsep=4pt,label={\color{poporange}\textbullet}}
\setlength{\parindent}{0pt}
\setlength{\parskip}{5pt}
\renewcommand{\arraystretch}{1.25}

% pgfplots : style Pop par défaut
\pgfplotsset{
  every axis/.append style={axis line style={popink,line width=1pt},tick style={popink},
    grid style={popline,line width=0.5pt},label style={font=\small},tick label style={font=\footnotesize}},
  popcurve/.style={poporange,line width=1.8pt,no markers,smooth},
}
% Même style disponible dans un \draw TikZ simple (hors pgfplots)
\tikzset{popcurve/.style={poporange,line width=1.8pt}}
\tikzset{popgrid/.style={popline,very thin}}
\colorlet{popcurve}{poporange} % le nom du style est parfois utilisé comme couleur

% ============================================================
% Pill / squiggle
% ============================================================
\newcommand{\poppill}[3]{%
  \tikz[baseline=-0.55ex]\node[draw=popink,line width=1.2pt,rounded corners=2.6pt,%
    fill=#2,inner xsep=6.5pt,inner ysep=3pt]{%
    \popxboldls\fontsize{7}{7}\selectfont\color{#3}\MakeUppercase{#1}};%
}
\newcommand{\popsquiggle}[1]{%
  \tikz[baseline=-0.4ex]{
    \draw[#1,line width=2.6pt,line cap=round]
      plot [smooth] coordinates {(0,0.09) (0.34,0.01) (0.68,0.21) (1.02,0.01) (1.36,0.10)};
  }%
}
% Mot-clé surligné (marqueur orange sous le texte)
\newcommand{\cle}[1]{{\popxbold\color{popdark}#1}}

% ============================================================
% En-tête / pied
% ============================================================
\pagestyle{fancy}
\fancyhf{}
\renewcommand{\headrulewidth}{0pt}
\renewcommand{\footrulewidth}{0pt}
\lhead{\raisebox{-0.15ex}{\poplogo\fontsize{13}{13}\selectfont\color{popink}OnBuch\color{poporange}+}}
\rhead{\poppill{\DOCMATIERE\ \textbullet\ \DOCNIVEAU\ \textbullet\ Leçon \DOCLECON}{white}{popink}}
\lfoot{\popxbold\fontsize{7.6}{7.6}\selectfont\color{popmuted}\MakeUppercase{Cours signé OnBuch+}\ \textbullet\ {\popsemi\DOCTITRE}}
\rfoot{\tikz[baseline=-0.6ex]\node[rounded corners=8.5pt,fill=popink,minimum height=17pt,minimum width=17pt,inner xsep=5pt,inner sep=0]{\popxbold\fontsize{8}{8}\selectfont\color{popcream}\thepage\,/\,\pageref*{LastPage}};}

% ============================================================
% Titres
% ============================================================
\newcommand{\chaptitre}[2]{%
  \begin{tcolorbox}[enhanced,colback=poporange,colframe=popink,boxrule=2.6pt,arc=11pt,
      drop shadow=popink,left=15pt,right=15pt,top=12pt,bottom=11pt,before skip=3pt,after skip=18pt]
    \poppill{#2}{popink}{popcream}\par\vspace{9pt}
    {\raggedright\hyphenpenalty=10000\exhyphenpenalty=10000\popdisplay\fontsize{22}{24}\selectfont\color{popcream}\MakeUppercase{#1}\par}
  \end{tcolorbox}
}
% Section numérotée : pastille orange avec le numéro + titre Archivo
\newcounter{popsec}
\newcommand{\coursec}[1]{%
  \stepcounter{popsec}\setcounter{popsub}{0}%
  \par\vspace{16pt}\noindent
  \begin{minipage}{\linewidth}
  \tikz[baseline=-0.7ex]\node[draw=popink,line width=1.6pt,fill=poporange,rounded corners=4pt,minimum size=22pt,inner sep=0]{\popdisplay\fontsize{12}{12}\selectfont\color{popcream}\thepopsec};%
  \hspace{8pt}{\popdisplay\fontsize{15}{17}\selectfont\color{popink}\MakeUppercase{#1}}\par
  \vspace{4pt}\noindent{\color{poporange}\rule{36pt}{3.6pt}}
  \end{minipage}\par\nopagebreak\vspace{8pt}
}
\newcounter{popsub}
\newcommand{\courssub}[1]{%
  \stepcounter{popsub}%
  \par\vspace{9pt}\noindent{\popxbold\fontsize{11.5}{13}\selectfont\color{popink}{\color{poporange}\thepopsec.\thepopsub}\hspace{6pt}#1}\par\nopagebreak\vspace{4pt}
}

% ============================================================
% Boîtes pédagogiques — même recette : fond blanc, bordure épaisse colorée,
% ombre dure encre, badge plein. Titre optionnel [#1].
% ============================================================
\tcbset{popbase/.style={breakable,enhanced,colback=white,boxrule=2.3pt,arc=9pt,
  shadow={3pt}{-3pt}{0pt}{popink},left=14pt,right=14pt,top=11pt,bottom=11pt,before skip=11pt,after skip=15pt,
  fontupper=\popsemi\fontsize{10.6}{14.5}\selectfont\color{popink},
  fontlower=\popsemi\fontsize{10.6}{14.5}\selectfont\color{popink}}}
% Coche dessinée (le glyphe ✓ n'existe pas dans les polices du document)
\newcommand{\popcheck}[1]{\tikz[baseline=-0.5ex]\draw[#1,line width=1.6pt,line cap=round,line join=round](-0.13,0.0)--(-0.04,-0.09)--(0.13,0.1);}
\providecommand{\coche}{\popcheck{popgreen}}
\providecommand{\resultat}[1]{\par\smallskip\noindent\fcolorbox{popgreen}{popgreenL}{\parbox{\dimexpr\linewidth-2\fboxsep-2\fboxrule\relax}{\textbf{Résultat.} #1}}\par\smallskip}
\newcommand{\popboxhead}[3]{\poppill{#1}{#2}{white}\ifx\relax#3\relax\else\hspace{6pt}{\popxbold\fontsize{10.6}{12}\selectfont\color{popink}#3}\fi\par\vspace{7pt}}

\newtcolorbox{aretenir}[1][]{popbase,colframe=popgreen,before upper={\popboxhead{À retenir}{popgreen}{#1}}}
\newtcolorbox{exemplebox}[1][]{popbase,colframe=poporange,before upper={\popboxhead{Exemple}{poporange}{#1}}}
\newtcolorbox{definition}[1][]{popbase,colframe=popblue,before upper={\popboxhead{Définition}{popblue}{#1}}}
\newtcolorbox{propriete}[1][]{popbase,colframe=poppurple,before upper={\popboxhead{Propriété}{poppurple}{#1}}}
\newtcolorbox{methode}[1][]{popbase,colframe=popgold,before upper={\popboxhead{Méthode}{popgold}{#1}}}
\newtcolorbox{attention}[1][]{popbase,colframe=poppink,before upper={\popboxhead{Attention}{poppink}{#1}}}
\newtcolorbox{experience}[1][]{popbase,colframe=popblue,colback=popblueL,before upper={\popboxhead{Expérience}{popblue}{#1}}}
% « Le savais-tu ? » : carte violette pleine (contexte, culture, Cameroun)
\newtcolorbox{savaistu}[1][]{popbase,colback=poppurple,colframe=popink,
  fontupper=\popsemi\fontsize{10.4}{14.2}\selectfont\color{white},
  before upper={\poppill{Le savais-tu\,?}{white}{poppurple}\ifx\relax#1\relax\else\hspace{6pt}{\popxbold\color{white}#1}\fi\par\vspace{7pt}}}
% Exercice résolu : énoncé en haut, \tcblower puis la solution sur fond crème
\newcounter{popexo}\newcounter{popexr}
\newtcolorbox{exoresolu}[1][]{popbase,colframe=poporange,colbacklower=popcream,
  segmentation style={popink,line width=1.2pt,dashed},
  before upper={\stepcounter{popexr}\popboxhead{Exercice résolu \thepopexr}{poporange}{#1}},
  before lower={\poppill{Solution}{popink}{popcream}\par\vspace{6pt}}}
% Exercice (non résolu en place) — la correction est dans la partie Corrigés
\newtcolorbox{exercice}[1][]{popbase,colframe=popink,
  before upper={\stepcounter{popexo}\popboxhead{Exercice \thepopexo}{popink}{#1}}}
% Corrigé
\newtcolorbox{corrige}[1][]{popbase,colframe=popgreen,colback=popgreenL,
  before upper={\popboxhead{Corrigé}{popgreen}{#1}}}
% Objectifs / prérequis / bilan
\newtcolorbox{objectifs}{popbase,colframe=popink,colback=poporangeL,
  before upper={\popboxhead{Objectifs de la leçon}{popink}{}}}
\newtcolorbox{prerequis}{popbase,colframe=popink,colback=popblueL,
  before upper={\popboxhead{Prérequis}{popblue}{}}}
\newtcolorbox{bilan}{popbase,colframe=popink,colback=white,boxrule=2.8pt,
  before upper={\popboxhead{Fiche bilan}{poporange}{L'essentiel à connaître pour l'examen}}}
% Figure encadrée avec légende
\newtcolorbox{popfigure}[1][]{enhanced,breakable=false,colback=white,colframe=popline,boxrule=1.4pt,arc=8pt,
  left=10pt,right=10pt,top=10pt,bottom=8pt,before skip=10pt,after skip=12pt,halign=center,
  lower separated=false,halign lower=center,
  fontlower=\popsemi\fontsize{9}{11.5}\selectfont\color{popmuted},
  #1}
\newcommand{\legende}[1]{\par\vspace{6pt}{\popsemi\fontsize{9}{11.5}\selectfont\color{popmuted}#1}}

% ============================================================
% Signature OnBuch+
% ============================================================
\newcommand{\poplogoinline}[1]{{\poplogo\fontsize{#1}{#1}\selectfont\color{popink}OnBuch\color{poporange}+}}
\newcommand{\signatureonbuch}{%
  \begin{tcolorbox}[enhanced,colback=popink,colframe=popink,boxrule=2.2pt,arc=10pt,
      drop shadow=poporange,left=14pt,right=14pt,top=10pt,bottom=10pt,before skip=0pt,after skip=0pt]
    \tikz[baseline=-0.7ex]\node[circle,fill=popgreen,draw=popcream,line width=1.3pt,minimum size=22pt,inner sep=0]{\popcheck{white}};%
    \hspace{9pt}\begin{minipage}[c]{0.62\linewidth}
      {\popxbold\fontsize{12}{13}\selectfont\color{popcream}Signé\ \textbullet\ {\poplogo OnBuch\color{poporange}+}}\par\vspace{2pt}
      {\raggedright\popsemi\fontsize{8}{10.5}\selectfont\color{popline}\MakeUppercase{Équipe pédagogique OnBuch+}\\\MakeUppercase{Conforme au programme officiel MINESEC}\par}
    \end{minipage}\hfill
    {\poplogo\fontsize{22}{22}\selectfont\color{popcream}OnBuch\color{poporange}+}
  \end{tcolorbox}%
}

% ============================================================
% Page de garde
% #1 titre · #2 matière · #3 niveau · #4 module · #5 n° leçon · #6 durée ·
% #7 description / objectifs (texte ou itemize)
% ============================================================
\newcommand{\pagegarde}[7]{%
  \thispagestyle{empty}
  \newgeometry{a4paper,margin=1.7cm,top=1.6cm,bottom=1.4cm}%
  \begin{tikzpicture}[remember picture,overlay]
    \fill[poporange!18] ([xshift=-3.2cm,yshift=-3.6cm]current page.north east) circle (4.6cm);
    \fill[poppurple!14] ([xshift=2.4cm,yshift=7.6cm]current page.south west) circle (3.8cm);
  \end{tikzpicture}%
  {\poplogo\fontsize{30}{30}\selectfont\color{popink}OnBuch\color{poporange}+}\hfill
  \raisebox{0.6ex}{\poppill{Cours signé \textbullet\ Premium}{popink}{popcream}}\par
  \vspace{1pt}\popsquiggle{poppurple}\par
  \vspace{8pt}
  \begin{tcolorbox}[enhanced,colback=poporange,colframe=popink,boxrule=2.8pt,arc=13pt,
      drop shadow=popink,left=18pt,right=18pt,top=12pt,bottom=13pt,after skip=10pt]
    \poppill{#2 \textbullet\ #3}{popink}{popcream}\hspace{5pt}\poppill{#4}{white}{popink}\par\vspace{8pt}
    {\popdisplay\fontsize{14}{14}\selectfont\color{popink}LEÇON #5}\par\vspace{4pt}
    {\raggedright\hyphenpenalty=10000\exhyphenpenalty=10000\popdisplay\fontsize{25}{27}\selectfont\color{popcream}\MakeUppercase{#1}\par}
    \vspace{8pt}
    {\popxbold\fontsize{9.5}{9.5}\selectfont\color{popink}\MakeUppercase{Durée conseillée : #6}}
  \end{tcolorbox}
  \begin{tcolorbox}[enhanced,colback=white,colframe=popink,boxrule=2.2pt,arc=10pt,
      drop shadow=popink,left=16pt,right=16pt,top=10pt,bottom=10pt,after skip=8pt]
    \poppill{Dans cette leçon}{poppurple}{white}\par\vspace{5pt}
    \popsemi\fontsize{9.3}{12.6}\selectfont\color{popink}#7
  \end{tcolorbox}
  \noindent\begin{minipage}[t]{0.487\linewidth}
    \begin{tcolorbox}[enhanced,colback=popgreen,colframe=popink,boxrule=2.2pt,arc=10pt,
        drop shadow=popink,left=11pt,right=11pt,top=10pt,bottom=10pt,height=3.1cm,valign=center]
      \tcbox[colback=white,colframe=popink,boxrule=1.3pt,arc=5pt,boxsep=0pt,nobeforeafter,
        left=3pt,right=3pt,top=3pt,bottom=3pt,valign=center]{\qrcode[height=1.9cm]{https://play.google.com/store/apps/details?id=com.luvvix.onbuch&hl=fr}}%
      \hspace{8pt}\begin{minipage}[c]{0.5\linewidth}
        {\popdisplay\fontsize{11}{12.5}\selectfont\color{white}\MakeUppercase{Télécharge OnBuch}}\par\vspace{4pt}
        {\raggedright\popsemi\fontsize{8.4}{11}\selectfont\color{white}Gratuit \textbullet\ Google Play\\Tuteur IA, annales, concours, résultats\par}
      \end{minipage}
    \end{tcolorbox}
  \end{minipage}\hfill
  \begin{minipage}[t]{0.487\linewidth}
    \begin{tcolorbox}[enhanced,colback=poppurple,colframe=popink,boxrule=2.2pt,arc=10pt,
        drop shadow=popink,left=11pt,right=11pt,top=10pt,bottom=10pt,height=3.1cm,valign=center]
      \tikz[baseline=-0.7ex]\node[circle,fill=white,draw=popink,line width=1.3pt,minimum size=1.6cm,inner sep=0]{\popdisplay\fontsize{24}{24}\selectfont\color{poppurple}+};%
      \hspace{8pt}\begin{minipage}[c]{0.6\linewidth}
        {\popdisplay\fontsize{11}{12.5}\selectfont\color{white}\MakeUppercase{Passe à OnBuch+}}\par\vspace{4pt}
        {\raggedright\popsemi\fontsize{8.4}{11}\selectfont\color{white}Tous les cours signés, Tuteur IA illimité, fiches PDF — onglet OnBuch+ de l'app\par}
      \end{minipage}
    \end{tcolorbox}
  \end{minipage}
  \vspace{8pt}
  \popcontenu
  \vfill
  \signatureonbuch
  \restoregeometry
  \newpage
}

% Rangée « Ce cours contient » (page de garde)
\newcommand{\poptile}[3]{% #1 couleur #2 gros texte #3 légende
  \begin{tcolorbox}[enhanced,colback=#1,colframe=popink,boxrule=2pt,arc=9pt,shadow={2.5pt}{-2.5pt}{0pt}{popink},
      left=6pt,right=6pt,top=9pt,bottom=9pt,halign=center,valign=center,height=1.75cm,width=\linewidth,nobeforeafter]
    {\popdisplay\fontsize{12.5}{13}\selectfont\color{white}\MakeUppercase{#2}}\par\vspace{3pt}
    {\centering\popsemi\fontsize{7.8}{9.5}\selectfont\color{white}#3\par}
  \end{tcolorbox}}
\newcommand{\popcontenu}{%
  \par\noindent{\popxboldls\fontsize{7.5}{7.5}\selectfont\color{popmuted}CE COURS SIGNÉ CONTIENT}\par\vspace{7pt}
  \noindent\begin{minipage}[t]{0.235\linewidth}\poptile{popblue}{Cours}{complet, illustré, pas à pas}\end{minipage}\hfill
  \begin{minipage}[t]{0.235\linewidth}\poptile{poporange}{Méthodes}{et exercices résolus}\end{minipage}\hfill
  \begin{minipage}[t]{0.235\linewidth}\poptile{poppink}{Exercices}{type examen + corrigés}\end{minipage}\hfill
  \begin{minipage}[t]{0.235\linewidth}\poptile{popgreen}{Fiche bilan}{l'essentiel à retenir}\end{minipage}\par}

% Sommaire
\newtcolorbox{sommaire}{enhanced,colback=white,colframe=popink,boxrule=2.2pt,arc=10pt,
  drop shadow=popink,left=18pt,right=18pt,top=13pt,bottom=13pt,before skip=6pt,after skip=20pt,
  before upper={\poppill{Sommaire}{poppurple}{white}\par\vspace{9pt}\popsemi\fontsize{10.6}{14.5}\selectfont\color{popink}}}

% Signature de fin de cours — poussée tout en bas de la dernière page
\newcommand{\findecours}{%
  \par\vfill\nopagebreak
  \begin{center}\popsquiggle{poporange}\end{center}\vspace{4pt}
  \signatureonbuch
}
\AtBeginDocument{\def\llbracket{\mathopen{[\mkern-3.5mu[}}\def\rrbracket{\mathclose{]\mkern-3.5mu]}}} % intervalles d'entiers (glyphe ⟦ absent de la police)
% Glyphes absents de Fira Math : redéfinis à partir de glyphes disponibles
\AtBeginDocument{%
  \def\setminus{\mathbin{\backslash}}\let\smallsetminus\setminus
  \def\vdots{\mathord{\vbox{\baselineskip3.2pt\lineskiplimit0pt\kern4pt\hbox{.}\hbox{.}\hbox{.}}}}%
  \def\ddots{\mathinner{\mkern1mu\raise6.5pt\hbox{.}\mkern2mu\raise3.6pt\hbox{.}\mkern2mu\raise0.7pt\hbox{.}\mkern1mu}}%
  \def\triangle{\mathord{\scalebox{0.9}{$\symup{\Delta}$}}}%
  \def\nabla{\mathord{\raisebox{\depth}{\scalebox{1}[-1]{$\symup{\Delta}$}}}}%
  \def\checkmark{\popcheck{popdark}}%
  \def\star{\mathbin{\ast}}\def\bigstar{\mathord{\ast}}%
  \def\diamond{\mathbin{\scalebox{0.6}{$\Diamond$}}}%
  \def\bigcup{\mathop{\mathchoice{\raisebox{-0.3em}{\scalebox{1.7}{$\cup$}}}{\scalebox{1.25}{$\cup$}}{\cup}{\cup}}\displaylimits}%
  \def\bigcap{\mathop{\mathchoice{\raisebox{-0.3em}{\scalebox{1.7}{$\cap$}}}{\scalebox{1.25}{$\cap$}}{\cap}{\cap}}\displaylimits}%
  \def\lesseqqgtr{\mathrel{\lessgtr}}\def\bigoplus{\mathop{\oplus}\displaylimits}%
}
\providecommand{\ssi}{si et seulement si} % abréviation fréquente
\AtBeginDocument{\providecommand{\wideparen}{\overparen}} % arc de cercle

\begin{document}

\pagegarde{Environment well-being and health}{Anglais}{Tle Tech}{Anglais – classe de Terminale technique}{N03}{13 séances (fiche de progression harmonisée)}{Cette leçon développe les compétences langagières pour analyser et communiquer sur les problèmes de santé publique (drogues, maladies transmissibles) et les enjeux environnementaux (changement climatique, gestion des déchets) au Cameroun. Elle consolide la grammaire complexe (relatives, discours rapporté, conditionnels, structures adjectivales) pour une expression précise et argumentée.
\vspace{4pt}
\begin{multicols}{2}\small
\begin{itemize}
\item À la fin de cette leçon, je sais utiliser les pronoms relatifs (who, which, whose, where) pour décrire avec précision les acteurs et les lieux liés aux problèmes de toxicomanie.
\item Je sais transformer des témoignages directs en discours rapporté pour rédiger un rapport d'alerte sur une épidémie dans mon établissement.
\item Je sais employer le vocabulaire spécifique (addiction, withdrawal, recycling, landfill, greenhouse effect) pour expliquer les causes et conséquences du réchauffement climatique au Cameroun.
\item Je sais construire des phrases au premier et second conditionnel pour proposer des solutions réalistes (If we sort waste...) ou hypothétiques (If I were the Minister...).
\item Je sais manipuler les structures 'how + adjective' et 'too... to' pour évaluer la gravité des situations sanitaires et environnementales.
\item Je sais rédiger une composition structurée (introduction, développement, conclusion) sur la gestion des ordures dans ma ville, en justifiant mes arguments.
\end{itemize}\end{multicols}}

\begin{sommaire}
\begin{multicols}{2}
\begin{enumerate}
\item 1. Substance Abuse : Description \& Analyse Grammaticale (Relative Clauses)
\item 2. Preventing Communicable Diseases : Listening \& Reported Speech
\item 3. Environmental Issues : Climate Change \& Waste Management (Reading \& Conditionals)
\item 4. Advanced Adjective Structures : Degree \& Comparison (How + Adj, Too...To, Superlatives)
\item 5. Writing \& Integration : From Analysis to Action (Composition \& Project)
\item Activité d'intégration
\item Méthodes et exercices résolus
\item Exercices
\item Corrigés des exercices
\item Fiche bilan
\end{enumerate}
\end{multicols}
\end{sommaire}

\begin{objectifs}
\begin{itemize}
\item À la fin de cette leçon, je sais utiliser les pronoms relatifs (who, which, whose, where) pour décrire avec précision les acteurs et les lieux liés aux problèmes de toxicomanie.
\item Je sais transformer des témoignages directs en discours rapporté pour rédiger un rapport d'alerte sur une épidémie dans mon établissement.
\item Je sais employer le vocabulaire spécifique (addiction, withdrawal, recycling, landfill, greenhouse effect) pour expliquer les causes et conséquences du réchauffement climatique au Cameroun.
\item Je sais construire des phrases au premier et second conditionnel pour proposer des solutions réalistes (If we sort waste...) ou hypothétiques (If I were the Minister...).
\item Je sais manipuler les structures 'how + adjective' et 'too... to' pour évaluer la gravité des situations sanitaires et environnementales.
\item Je sais rédiger une composition structurée (introduction, développement, conclusion) sur la gestion des ordures dans ma ville, en justifiant mes arguments.
\item Je sais présenter oralement une campagne de prévention contre l'alcoolisme ou le tabagisme en milieu scolaire en utilisant les connecteurs logiques appropriés.
\end{itemize}
\end{objectifs}

\begin{prerequis}
\begin{itemize}
\item Maîtrise des temps simples (Present Simple, Past Simple) et du Present Perfect pour relier passé et présent.
\item Connaissance de base des pronoms relatifs simples (who, which, that) vus en Première.
\item Structure du conditionnel de base (If + present -> will + V) et notions de modal verbs (must, should, can).
\item Vocabulaire de base sur le corps humain, les maladies courantes et l'environnement immédiat (déforestation, pollution).
\item Capacité à rédiger un paragraphe cohérent avec une phrase sujet et des arguments simples.
\end{itemize}
\end{prerequis}

\newpage

\coursec{Substance Abuse: Description \& Grammatical Analysis (Relative Clauses)}

\courssub{Hook Situation \& Problem Statement}

\begin{popfigure}
\begin{tikzpicture}[scale=0.85, every node/.style={font=\small}]
% Central node
\node[draw, rounded corners=8pt, fill=popblue!15, thick, align=center, minimum width=4.5cm, minimum height=1.8cm] (main) at (0,0) {\textbf{Substance Abuse in Cameroon}\\(Tramadol, Kobolo, Alcohol)};

% Causes
\node[draw, rounded corners=8pt, fill=poporange!15, thick, align=center, minimum width=3.2cm, minimum height=1.5cm] (cause1) at (-5,2.5) {Peer Pressure\\(Influence du groupe)};
\node[draw, rounded corners=8pt, fill=poporange!15, thick, align=center, minimum width=3.2cm, minimum height=1.5cm] (cause2) at (-5,0) {Unemployment\\(Ch\^omage des jeunes)};
\node[draw, rounded corners=8pt, fill=poporange!15, thick, align=center, minimum width=3.2cm, minimum height=1.5cm] (cause3) at (-5,-2.5) {Stress \& Trauma\\(Stress, traumatismes)};

% Substances
\node[draw, rounded corners=8pt, fill=poppink!25, thick, align=center, minimum width=3.2cm, minimum height=1.5cm] (drug1) at (0,-3.5) {Tramadol\\(Opio\"ide de synth\`ese)};
\node[draw, rounded corners=8pt, fill=poppink!25, thick, align=center, minimum width=3.2cm, minimum height=1.5cm] (drug2) at (-3.2,-3.5) {Kobolo\\(Alcool frelat\'e + m\'edocs)};
\node[draw, rounded corners=8pt, fill=poppink!25, thick, align=center, minimum width=3.2cm, minimum height=1.5cm] (drug3) at (3.2,-3.5) {Alcohol / Cannabis\\(Produits accessibles)};

% Consequences
\node[draw, rounded corners=8pt, fill=popgreen!15, thick, align=center, minimum width=3.2cm, minimum height=1.5cm] (cons1) at (5,2.5) {School Dropout\\(D\'ecrochage scolaire)};
\node[draw, rounded corners=8pt, fill=popgreen!15, thick, align=center, minimum width=3.2cm, minimum height=1.5cm] (cons2) at (5,0) {Mental Health Issues\\(Troubles psychiques)};
\node[draw, rounded corners=8pt, fill=popgreen!15, thick, align=center, minimum width=3.2cm, minimum height=1.5cm] (cons3) at (5,-2.5) {Crime \& Violence\\(D\'elinquance, violences)};

% Arrows
\foreach \src/\dst in {cause1/main, cause2/main, cause3/main, main/drug1, main/drug2, main/drug3, main/cons1, main/cons2, main/cons3}
  \draw[->, thick, popdark] (\src) -- (\dst);

% Relative clause examples
\node[align=center, font=\small\itshape, text=popblue!80!black] at (-5,-4.3) {Relative Clause:\\ \og Youths \textbf{who} face unemployment \textbf{turn} to drugs.\fg};
\node[align=center, font=\small\itshape, text=poppink!80!black] at (0,-5.3) {Relative Clause:\\ \og Tramadol, \textbf{which} is cheap, \textbf{destroys} lives.\fg};
\node[align=center, font=\small\itshape, text=popgreen!80!black] at (5,-4.3) {Relative Clause:\\ \og Effects \textbf{that} ruin futures \textbf{must} be stopped.\fg};
\end{tikzpicture}
\legende{Concept map: causes, substances and consequences of substance abuse in Cameroon, with relative clause examples.}
\end{popfigure}

\textbf{Problem Statement:} In Douala, gutters blocked by plastic waste favour the spread of mosquitoes and cholera, while in the secondary schools of Yaound\'e, the consumption of tramadol and \og kobolo\fg{} among students worries health authorities. Faced with these challenges, how can a Terminale technique student express himself or herself in English to alert, prevent and propose concrete solutions? This lesson gives you the linguistic tools to become an agent of change in your community.

\courssub{Vocabulary Building: Substance Abuse Lexicon}

\begin{definition}[Key Vocabulary -- Mots-cl\'es du th\`eme]
The following terms are essential to talk about addictions in a technical and Cameroonian context. Learn them with their meaning and their French equivalent.
\end{definition}

\begin{center}
\begin{tabularx}{\linewidth}{|l|X|l|}\hline
\rowcolor{popblueL}
\textbf{English Term} & \textbf{Definition / Context} & \textbf{Fran\c{c}ais} \\ \hline
\cle{Drug abuse} & Misuse of legal or illegal substances for non-medical purposes & Abus de drogues \\ \hline
\cle{Addiction} & Chronic condition characterized by compulsive drug seeking & D\'ependance, toxicomanie \\ \hline
\cle{Peer pressure} & Influence from friends or age-mates to adopt risky behaviour & Pression des pairs \\ \hline
\cle{Overdose} & Excessive dose causing coma or death & Surdose \\ \hline
\cle{Rehabilitation} & Medical and psychological process to recover from addiction & R\'ehabilitation, sevrage encadr\'e \\ \hline
\cle{Withdrawal symptoms} & Physical and mental effects when stopping a drug (tremors, anxiety) & Sympt\^omes de sevrage \\ \hline
\cle{Narcotics} & Drugs that dull the senses, relieve pain and induce sleep (opioids) & Stup\'efiants, narcotiques \\ \hline
\cle{Tramadol} & Synthetic opioid painkiller widely abused in Cameroon (street name: \og Tramol\fg) & Tramadol (m\'edicament d\'etourn\'e) \\ \hline
\cle{Kobolo} & Local street cocktail: cheap alcohol mixed with crushed pills & Kobolo (m\'elange local dangereux) \\ \hline
\cle{Relapse} & Return to drug use after a period of abstinence & Rechute \\ \hline
\cle{Harm reduction} & Policies reducing negative consequences without requiring abstinence & R\'eduction des risques \\ \hline
\end{tabularx}
\end{center}

\begin{exemplebox}[Vocabulary in Context -- Contexte camerounais]
\textbf{Text:} \og In many neighbourhoods of Yaound\'e and Douala, \textbf{peer pressure} pushes teenagers to try \textbf{tramadol} or \textbf{kobolo}. What starts as curiosity becomes \textbf{addiction}. Users suffer severe \textbf{withdrawal symptoms} like tremors and hallucinations when they try to quit. \textbf{Rehabilitation} centres are overwhelmed. Prevention is cheaper than cure.\fg

\textbf{Traduction:} Dans de nombreux quartiers de Yaound\'e et Douala, la \textbf{pression des pairs} pousse les adolescents \`a essayer le \textbf{tramadol} ou le \textbf{kobolo}. Ce qui commence par curiosit\'e devient de la \textbf{d\'ependance}. Les usagers souffrent de graves \textbf{sympt\^omes de sevrage} comme des tremblements et des hallucinations quand ils essaient d'arr\^eter. Les centres de \textbf{r\'ehabilitation} sont d\'ebord\'es. La pr\'evention co\^ute moins cher que le traitement.
\end{exemplebox}

\begin{exercice}[Vocabulary Practice -- Exercice de vocabulaire]
\textbf{A. Match each term with its definition.}
\begin{enumerate}
\item Addiction \hfill a) Influence from friends to do something risky
\item Peer pressure \hfill b) Return to drug use after stopping
\item Overdose \hfill c) Taking more of a drug than the body can handle
\item Withdrawal symptoms \hfill d) Chronic compulsive drug seeking
\item Relapse \hfill e) Physical effects when stopping a drug
\end{enumerate}

\textbf{B. Complete the sentences with the correct word.}
\begin{enumerate}
\item Many youths start taking tramadol because of \_\_\_\_\_\_\_\_ from friends.
\item The \_\_\_\_\_\_\_\_ centre in Bertoua helps 50 patients per month.
\item \_\_\_\_\_\_\_\_ symptoms include sweating, anxiety, and insomnia.
\item A lethal \_\_\_\_\_\_\_\_ of tramadol can stop breathing.
\item \_\_\_\_\_\_\_\_ is a dangerous local mixture of alcohol and pills.
\end{enumerate}
\end{exercice}

\begin{corrige}[Exercice Vocabulary Practice]
\textbf{A.} 1-d, 2-a, 3-c, 4-e, 5-b \\
\textbf{B.} 1. peer pressure \quad 2. rehabilitation \quad 3. Withdrawal \quad 4. overdose \quad 5. Kobolo
\end{corrige}

\courssub{Reading Comprehension: Testimonials from Former Addicts}

\begin{exemplebox}[Reading Text: \og A Second Chance\fg -- T\'emoignages fictifs inspir\'es de la r\'ealit\'e camerounaise]
\textbf{Text 1 -- Paul, 22, Douala:}
\og I started tramadol at 17 because my friends \textbf{who} worked at the port told me it gives energy for night shifts. I believed them. Soon, I could not work without it. The drug \textbf{which} I took daily destroyed my stomach. I stole money from my mother \textbf{who} cried every night. Last year, I overdosed. The nurse \textbf{whose} shift it was saved my life. Now I am in a rehabilitation centre \textbf{where} I learn carpentry. My advice: do not start.\fg

\textbf{Text 2 -- Amina, 19, Yaound\'e:}
\og Kobolo was everywhere in my quarter. Girls \textbf{who} drank it said it makes you bold. I tried it when I was 15. The alcohol \textbf{that} was mixed with pills made me unconscious. I woke up in hospital \textbf{where} a social worker talked to me. She is the person \textbf{to whom} I owe my life. Today, I speak in schools \textbf{where} I warn students about kobolo.\fg

\textbf{Glossaire:} \cle{night shifts} = quarts de nuit; \cle{bold} = audacieuse, sans peur; \cle{unconscious} = inconsciente; \cle{to owe one's life to} = devoir la vie \`a.
\end{exemplebox}

\begin{exercice}[Reading Comprehension -- Compr\'ehension de lecture]
\textbf{Answer in complete sentences.}
\begin{enumerate}
\item Why did Paul start taking tramadol?
\item What does \og the drug \textbf{which} I took daily\fg refer to?
\item Who saved Paul's life?
\item Where does Amina now speak to warn students?
\item Find two relative clauses in Text 1 and identify their antecedents.
\end{enumerate}
\end{exercice}

\begin{corrige}[Exercice Reading Comprehension]
1. Paul started tramadol because friends who worked at the port told him it gives energy for night shifts. \\
2. \og The drug which I took daily\fg refers to tramadol. \\
3. The nurse whose shift it was saved Paul's life. \\
4. Amina speaks in schools where she warns students about kobolo. \\
5. Examples: \og friends \textbf{who} worked at the port\fg (antecedent: friends); \og the drug \textbf{which} I took daily\fg (antecedent: the drug); \og my mother \textbf{who} cried every night\fg (antecedent: my mother); \og the nurse \textbf{whose} shift it was\fg (antecedent: the nurse); \og a rehabilitation centre \textbf{where} I learn carpentry\fg (antecedent: centre).
\end{corrige}

\courssub{Grammar Focus: Relative Clauses -- Propositions Relatives}

\begin{definition}[Relative Clause -- Proposition relative]
A \textbf{relative clause} (proposition relative) is a subordinate clause that gives information about a noun called the \textbf{antecedent} (ant\'ec\'edent). It begins with a \textbf{relative pronoun} (pronom relatif) which links the clause to the antecedent.
\\[4pt]
\textbf{Structure:} Antecedent + Relative Pronoun + Verb + (Complements)
\\[4pt]
\textbf{Example:} \og The youth \textbf{who} consumes tramadol risks his life.\fg \\
Antecedent: \og The youth\fg \quad Pronoun: \og who\fg \quad Verb: \og consumes\fg
\end{definition}

\begin{propriete}[Choice of Relative Pronoun -- Choix du pronom relatif]
The relative pronoun depends on the \textbf{nature of the antecedent} (person/thing) and on its \textbf{grammatical function} inside the relative clause (subject, object, possession, place, time).
\begin{center}
\begin{tabularx}{\linewidth}{|c|l|X|X|}\hline
\rowcolor{popblueL}
\textbf{Pronoun} & \textbf{Antecedent} & \textbf{Function in the clause} & \textbf{Example} \\ \hline
\cle{who} & Person & Subject & The doctor \textbf{who} treats addicts is busy. \\ \hline
\cle{whom} & Person & Object (formal) & The patient \textbf{whom} the doctor examined... \\ \hline
\cle{which} & Thing/Animal & Subject or Object & The pill \textbf{which} causes addiction... \\ \hline
\cle{that} & Person/Thing & Subject or Object (defining only) & The youth \textbf{that} smokes kobolo... \\ \hline
\cle{whose} & Person/Thing & Possession & The father \textbf{whose} son is addicted... \\ \hline
\cle{where} & Place & Place (replaces in/at/on which) & The centre \textbf{where} he recovers... \\ \hline
\cle{when} & Time & Time (replaces on/in which) & The day \textbf{when} he quit... \\ \hline
\end{tabularx}
\end{center}
\end{propriete}

\begin{attention}[Pi\`eges fr\'equents -- Common Mistakes]
\begin{enumerate}
\item \textbf{Never use \og what\fg as a standard relative pronoun.} \\
\og The thing \textbf{what} I need\fg{} $\rightarrow$ \textbf{The thing \textbf{that/which} I need.}
\item \textbf{Do not keep a redundant personal pronoun.} \\
\og The boy \textbf{who he} is my friend\fg{} $\rightarrow$ \textbf{The boy \textbf{who} is my friend.}
\item \textbf{Whom vs Who:} \og whom\fg{} is formal and used as an object. In speech, we use \og who\fg{} or we omit the pronoun. \\
\og The man \textbf{whom} I saw\fg{} (formal) = \og The man \textbf{who} I saw\fg{} = \og The man I saw\fg.
\item \textbf{\og That\fg{} is NEVER used in non-defining relative clauses (with commas).} \\
\og My father, \textbf{that} is a doctor...\fg{} $\rightarrow$ \textbf{My father, \textbf{who} is a doctor...}
\end{enumerate}
\end{attention}

\courssub{Defining vs Non-Defining Relative Clauses}

\begin{propriete}[Restrictive (Defining) vs Non-Restrictive (Non-Defining)]
The distinction changes the \textbf{meaning} and the \textbf{punctuation}.
\begin{center}
\begin{tabularx}{\linewidth}{|l|X|X|}\hline
\rowcolor{popblueL}
\textbf{Criterion} & \textbf{Defining (Restrictive)} & \textbf{Non-Defining (Non-Restrictive)} \\ \hline
\textbf{Meaning} & \textbf{Essential} information to identify the antecedent & \textbf{Extra} information, not needed for identification \\ \hline
\textbf{Punctuation} & \textbf{No commas} & \textbf{Between commas} (or dashes/brackets) \\ \hline
\textbf{Pronouns} & who, which, that, whose, where, when & who, which, whose, where, when (\textbf{never that}) \\ \hline
\textbf{Omission} & Possible if the pronoun is the object & \textbf{Impossible} \\ \hline
\textbf{Example} & Students \textbf{who} take tramadol risk expulsion. & Paul, \textbf{who} takes tramadol, risks expulsion. \\ \hline
\textbf{Translation} & Les \'el\`eves \textbf{qui} prennent du tramadol risquent l'exclusion. (Only those who take it) & Paul, \textbf{qui} prend du tramadol, risque l'exclusion. (Extra information about Paul) \\ \hline
\end{tabularx}
\end{center}
\end{propriete}

\begin{exemplebox}[Minimal Pair -- Paire minimale: the meaning changes completely]
\begin{enumerate}
\item \og The students \textbf{who} failed the test must retake it.\fg \\
$\rightarrow$ \textbf{Only} those who failed will retake it. (Defining)
\item \og The students, \textbf{who} failed the test, must retake it.\fg \\
$\rightarrow$ \textbf{All} the students failed and all will retake it. (Non-defining)
\end{enumerate}
\textbf{Prevention application:} \og Youths \textbf{who} consume kobolo destroy their health.\fg{} (Only those who consume it) vs \og Youths, \textbf{who} consume kobolo, destroy their health.\fg{} (All youths consume it -- false!)
\end{exemplebox}

\begin{exemplebox}[Relative Clauses with Statistics -- Relatives et statistiques]
\textbf{Raw data (illustrative figure):} \og About one student in five aged 15--19 has already tried tramadol.\fg
\\[4pt]
\textbf{Transformation into relative clauses:}
\begin{itemize}
\item \og One student in five \textbf{who} is aged 15--19 has already tried tramadol.\fg{} (Defining)
\item \og Tramadol, \textbf{which} is a synthetic opioid, is one of the most abused drugs in schools.\fg{} (Non-defining)
\item \og The regions \textbf{where} consumption is highest are the Littoral and the Centre.\fg{} (Place)
\item \og The year \textbf{when} tramadol seizures peaked was 2022.\fg{} (Time)
\item \og Parents \textbf{whose} children are addicted often feel helpless.\fg{} (Possession)
\end{itemize}
\end{exemplebox}

\courssub{Method: Combining Sentences with Relative Clauses}

\begin{methode}[How to combine two sentences into one with a relative clause -- 4 steps]
\textbf{Aim:} Transform two simple sentences into one complex sentence containing a relative clause.
\\[6pt]
\textbf{Starting example:}
\begin{enumerate}
\item The youths smoke kobolo. (Sentence 1)
\item Kobolo destroys their lungs. (Sentence 2)
\end{enumerate}

\textbf{Step 1 -- Identify the common antecedent.} \\
\og Kobolo\fg{} (a thing) is the antecedent. It is the subject in Sentence 2.

\textbf{Step 2 -- Choose the relative pronoun.} \\
Thing + subject $\rightarrow$ \textbf{which} (or \textbf{that} if defining).

\textbf{Step 3 -- Remove the redundant word.} \\
In Sentence 2, \og Kobolo\fg{} becomes \og which\fg. Delete any repeated pronoun such as \og it\fg.

\textbf{Step 4 -- Place the relative clause just after the antecedent.} \\
\og The youths smoke kobolo \textbf{which destroys their lungs}.\fg

\textbf{Final result:} \og The youths smoke kobolo \textbf{which destroys their lungs}.\fg
\end{methode}

\begin{exoresolu}[Sentence Combining -- Combinaison de phrases]
\textbf{Combine each pair using a relative pronoun.}

\textbf{1.} The doctor treated the addict. The doctor works at Laquintinie Hospital.
\tcblower
\textbf{Solution:} The doctor \textbf{who/that} treated the addict works at Laquintinie Hospital. (Defining: identifies which doctor) \\
OR: The doctor, \textbf{who} treated the addict, works at Laquintinie Hospital. (Non-defining: extra information)

\textbf{2.} Tramadol is a painkiller. Young people abuse it.
\\[4pt]
\textbf{Solution:} Tramadol, \textbf{which} young people abuse, is a painkiller. (Non-defining) \\
OR: Tramadol \textbf{that/which} young people abuse is a painkiller. (Defining -- less natural here)

\textbf{3.} The centre helps addicts. My brother works there.
\\[4pt]
\textbf{Solution:} The centre \textbf{where} my brother works helps addicts. (Place $\rightarrow$ where)

\textbf{4.} 2022 was a bad year. Tramadol seizures increased then.
\\[4pt]
\textbf{Solution:} 2022 was the year \textbf{when} tramadol seizures increased. (Time $\rightarrow$ when)

\textbf{5.} The mother cried. Her son was addicted.
\\[4pt]
\textbf{Solution:} The mother \textbf{whose} son was addicted cried. (Possession $\rightarrow$ whose)
\end{exoresolu}

\begin{exercice}[Sentence Combining Practice -- Entra\^inement combinaison]
\textbf{Combine each pair. Indicate if defining (D) or non-defining (ND).}
\begin{enumerate}
\item The social worker helped Amina. The social worker speaks in schools now.
\item Kobolo is a local drink. It contains alcohol and tramadol.
\item The rehabilitation centre is in Bafoussam. Paul stays there.
\item Peer pressure pushes youths to drugs. It is very strong in secondary schools.
\item The law punishes drug traffickers. It was passed in 1997.
\end{enumerate}
\end{exercice}

\begin{corrige}[Exercice Sentence Combining Practice]
1. The social worker \textbf{who/that} helped Amina speaks in schools now. (D) \quad OR \quad The social worker, \textbf{who} helped Amina, speaks in schools now. (ND) \\
2. Kobolo, \textbf{which} contains alcohol and tramadol, is a local drink. (ND) \\
3. The rehabilitation centre \textbf{where} Paul stays is in Bafoussam. (D) \\
4. Peer pressure, \textbf{which} is very strong in secondary schools, pushes youths to drugs. (ND) \\
5. The law \textbf{which/that} was passed in 1997 punishes drug traffickers. (D)
\end{corrige}

\courssub{Speaking Activity: Information Gap -- Prevention Poster Description}

\begin{methode}[Information Gap Activity -- D\'ecrire une image de pr\'evention]
\textbf{Instructions:} Work in pairs. Student A has Image A (healthy lung). Student B has Image B (lung damaged by smoking/kobolo). Without showing your image, describe it to your partner using \textbf{at least 3 relative clauses}. Your partner draws what he or she hears. Compare at the end.
\end{methode}

\begin{popfigure}
\begin{tikzpicture}[scale=0.9, every node/.style={font=\small}]
% Healthy lung
\begin{scope}[xshift=-5cm]
\draw[fill=poppink!30, thick] (0,0) ellipse (1.5 and 2.5);
\draw[fill=poppink!50] (0,0) ellipse (1.2 and 2.2);
\draw[poppink!70!black, thick] (-0.8,1.5) .. controls (-1.2,0.5) .. (-0.5,-0.5);
\draw[poppink!70!black, thick] (0.8,1.5) .. controls (1.2,0.5) .. (0.5,-0.5);
\draw[poppink!70!black, thick] (-0.5,1) .. controls (-0.9,0) .. (-0.4,-1);
\draw[poppink!70!black, thick] (0.5,1) .. controls (0.9,0) .. (0.4,-1);
\node[align=center, font=\bfseries, text=poppink!70!black] at (0,-3.2) {Image A: Healthy Lung\\(Poumon sain)};
\node[align=left, font=\footnotesize, text=popdark] at (2.6,1.5) {Pink, spongy\\Efficient gas exchange\\No tar deposits};
\end{scope}

% Damaged lung
\begin{scope}[xshift=5cm]
\draw[fill=popmuted!40, thick] (0,0) ellipse (1.5 and 2.5);
\draw[fill=popmuted!60] (0,0) ellipse (1.2 and 2.2);
\foreach \x/\y in {-0.8/1.2, -0.3/0.8, 0.5/-0.5, -1.0/-1.2, 0.2/-1.8, -0.6/-0.2, 0.9/0.3}
  \draw[fill=popdark!80, popdark] (\x,\y) circle (0.15);
\draw[popdark, thick] (-0.8,1.5) .. controls (-1.2,0.5) .. (-0.5,-0.5);
\draw[popdark, thick] (0.8,1.5) .. controls (1.2,0.5) .. (0.5,-0.5);
\node[align=center, font=\bfseries, text=popdark] at (0,-3.2) {Image B: Damaged Lung\\(Poumon ab\^im\'e)};
\node[align=left, font=\footnotesize, text=popdark] at (2.6,1.5) {Black/grey spots\\Tar deposits (goudron)\\Reduced capacity\\Emphysema risk};
\end{scope}

% Arrow between
\draw[<->, thick, popblue, dashed] (-3.2,0) -- (3.2,0) node[midway, above, font=\small] {Information Gap: Describe \& Draw};
\end{tikzpicture}
\legende{Information Gap: healthy lung (A) vs lung damaged by smoking/kobolo (B). Students describe using relative clauses.}
\end{popfigure}

\begin{exemplebox}[Model Description -- Description mod\`ele (Student A)]
\og This is a picture of a healthy lung \textbf{which} is pink and spongy. The surface \textbf{where} gas exchange happens is clean. There are no black spots \textbf{that} indicate tar deposits. The alveoli \textbf{which} look like tiny balloons are fully inflated. This lung belongs to a person \textbf{who} does not smoke kobolo or cigarettes.\fg

\textbf{Relatives used:} 1) \og which is pink\fg{} (defining), 2) \og where gas exchange happens\fg{} (place), 3) \og that indicate tar deposits\fg{} (defining), 4) \og which look like balloons\fg{} (defining), 5) \og who does not smoke\fg{} (defining).
\end{exemplebox}

\begin{exemplebox}[Model Description -- Description mod\`ele (Student B)]
\og This image shows a damaged lung \textbf{which} has turned grey and black. The tissue \textbf{where} oxygen should pass is blocked by tar. The black spots \textbf{that} you see are toxic deposits from kobolo and cigarettes. This is the lung of a youth \textbf{who} has smoked for five years. The doctor \textbf{whose} patient owns this lung warns that emphysema, \textbf{which} is incurable, will follow.\fg

\textbf{Relatives used:} 1) \og which has turned\fg, 2) \og where oxygen should pass\fg, 3) \og that you see\fg{} (omissible: \og The black spots you see\fg), 4) \og who has smoked\fg, 5) \og whose patient\fg, 6) \og which is incurable\fg{} (non-defining).
\end{exemplebox}

\begin{exercice}[Speaking Task -- T\^ache orale not\'ee]
\textbf{Prepare a 2-minute presentation of the poster below. Use at least 4 relative clauses (2 defining, 2 non-defining). Record yourself or present to the class.}

\begin{center}
\begin{tikzpicture}[scale=0.7]
% Poster mockup
\draw[fill=popcream, thick, rounded corners] (-4,-3) rectangle (4,3);
\draw[fill=poppink!30, rounded corners] (-3.5,1.5) rectangle (3.5,2.8);
\node[font=\bfseries\large] at (0,2.15) {STOP TRAMADOL \& KOBOLO!};
\draw[fill=popblue!10, rounded corners] (-3.5,-0.5) rectangle (3.5,1.2);
\node[align=center, font=\small] at (0,0.35) {Tramadol destroys your brain.\\Kobolo kills your future.\\YOU HAVE THE POWER TO SAY NO.};
\draw[fill=popgreen!10, rounded corners] (-3.5,-2.8) rectangle (3.5,-0.8);
\node[align=center, font=\small] at (0,-1.8) {Help line: 1510 (Free)\\Rehabilitation Centre -- Laquintinie Hospital\\You are not alone.};
\node[align=center, font=\small\itshape] at (0,-3.8) {Prevention poster -- School health campaign, Cameroon};
\end{tikzpicture}
\end{center}

\textbf{Assessment criteria:} 1) Accurate vocabulary (5 pts) \quad 2) At least 4 correct relative clauses (8 pts) \quad 3) Pronunciation and fluency (4 pts) \quad 4) Logical structure (3 pts) -- Total: 20 pts.
\end{exercice}

\courssub{Cultural Context \& Legal Framework: Cameroon Focus}

\begin{savaistu}[Did you know? -- Contexte l\'egal et social au Cameroun]
\begin{itemize}
\item \textbf{Law No. 97/19 of 10 January 1997:} punishes the use, possession and trafficking of narcotics in Cameroon, with prison sentences and heavy fines.
\item \textbf{Tramadol:} originally a painkiller (WHO level-2 analgesic). Diverted as a \og doping\fg{} product for hard physical work (dockers, drivers, farmers). Street doses (200--400 mg) far exceed therapeutic doses (50--100 mg).
\item \textbf{Kobolo:} a home-made mixture of adulterated alcohol and crushed pills (tramadol, diazepam). Sold very cheaply (about 100--300 FCFA per sachet), it is common in call-box areas and popular neighbourhoods.
\item \textbf{Health data:} health authorities report that a significant proportion of teenagers aged 15--19 have already tried tramadol, and that substance abuse is involved in many psychiatric admissions in reference hospitals such as Laquintinie (Douala) and Jamot (Yaound\'e).
\item \textbf{Emergency numbers:} 1510 (Police), 116 (Children in Danger), 119 (Medical emergencies). Care centres: Laquintinie (Douala), Jamot (Yaound\'e), and regional hospitals in Bafoussam, Garoua and Maroua.
\end{itemize}
\end{savaistu}

\begin{exemplebox}[Authentic Document Analysis -- Analyse de document authentique]
\textbf{Adapted extract -- Public health press release:}
\og The Ministry of Public Health, \textbf{which} coordinates the national drug control strategy, warns that tramadol \textbf{that} is sold illegally in markets contains unknown substances. Parents \textbf{whose} children show signs of addiction (isolation, aggression, weight loss) must contact the nearest health centre \textbf{where} free counselling is available. The campaign \textbf{whose} slogan is \og Zero Tramadol in Schools\fg{} targets secondary students \textbf{who} are the most vulnerable.\fg

\textbf{Activity:} Identify the 6 relative clauses. Classify them as D (defining) or ND (non-defining). Justify the punctuation.
\end{exemplebox}

\begin{corrige}[Analyse de document authentique]
1. \og \textbf{which} coordinates...\fg{} -- ND (comma before it; extra information about the Ministry) \\
2. \og \textbf{that} is sold illegally...\fg{} -- D (identifies which tramadol: the one sold illegally) \\
3. \og \textbf{whose} children show signs...\fg{} -- D (identifies which parents: those whose children show signs) \\
4. \og \textbf{where} free counselling is available\fg{} -- D (identifies which centre: the one where counselling is free) \\
5. \og \textbf{whose} slogan is...\fg{} -- D (identifies which campaign: the one whose slogan is...) \\
6. \og \textbf{who} are the most vulnerable\fg{} -- D (identifies which students: the most vulnerable ones)
\end{corrige}

\courssub{Integration \& Remediation: Mini-Project}

\begin{exercice}[Mini-Project: \og Youth Ambassadors Against Drugs\fg{} -- Projet int\'egr\'e]
\textbf{Scenario:} You are a health ambassador in your technical secondary school. Create a prevention poster (A4 format) in English for the new school year.

\textbf{Compulsory requirements:}
\begin{enumerate}
\item A catchy title (slogan).
\item 3 real causes (peer pressure, unemployment, stress...).
\item 3 consequences (health, studies, law).
\item 1 statistic on drug abuse (mention your source).
\item \textbf{At least 5 relative clauses} (mix of defining/non-defining, varied: who, which, where, whose, when).
\item Emergency numbers (1510, 116, 119).
\item A visual: drawing, diagram or simple infographic.
\end{enumerate}

\textbf{Self-assessment grid (tick):}
\begin{itemize}
\item[$\square$] Unit vocabulary used (addiction, withdrawal, rehabilitation, peer pressure, overdose, kobolo, tramadol)
\item[$\square$] 5+ correct relative clauses (punctuation respected)
\item[$\square$] Clear D / ND distinction (commas in the right place)
\item[$\square$] No \og what\fg{} as a relative pronoun, no redundant pronoun
\item[$\square$] Spelling and grammar checked
\item[$\square$] Clear message, appropriate tone (alerting without stigmatizing)
\end{itemize}

\textbf{Output:} Poster displayed in the school corridor + 2-minute oral presentation in class (continuous assessment).
\end{exercice}

\begin{aretenir}[Section 1 Summary -- Key points for the exam]
\begin{enumerate}
\item \textbf{Key vocabulary:} drug abuse, addiction, peer pressure, overdose, rehabilitation, withdrawal symptoms, narcotics, tramadol, kobolo, relapse, harm reduction.
\item \textbf{Relative clauses:} they give information about an antecedent. Pronouns: who (person/subject), whom (person/object, formal), which (thing), that (person/thing, defining only), whose (possession), where (place), when (time).
\item \textbf{Defining (restrictive):} essential information, \textbf{no commas}, \og that\fg{} allowed, omission possible when the pronoun is the object.
\item \textbf{Non-defining (non-restrictive):} extra information, \textbf{between commas}, \og that\fg{} \textbf{forbidden}, omission impossible.
\item \textbf{Combining sentences:} 1) find the common antecedent 2) choose the pronoun according to its function 3) remove the redundant word 4) place the clause after the antecedent.
\item \textbf{Cameroon context:} Law No. 97/19 of 1997, tramadol/kobolo as youth scourges, emergency numbers 1510 / 116 / 119.
\end{enumerate}
\end{aretenir}

\coursec{Preventing Communicable Diseases : Listening \& Reported Speech}

In the previous section, we analysed the devastating impact of substance abuse on health and practised \cle{relative clauses} to define terms like \emph{addiction} and \emph{withdrawal}. We now shift our focus from non-communicable behavioural risks to \cle{communicable diseases} — illnesses caused by viruses, bacteria, or parasites that spread from person to person or through vectors. In Cameroon, seasonal outbreaks of cholera, malaria, and meningitis require rapid public health responses. Mastering \cle{reported speech} is essential here: as future technicians, health workers, or informed citizens, you must accurately relay official health protocols, radio announcements, or medical advice without distorting the message.

\courssub{Vocabulary: Prevention \& Disease Classification}

Understanding the \cle{mode of transmission} is the first step in prevention. The following table categorises key diseases relevant to the Cameroonian context.

\begin{center}
\begin{tabularx}{\linewidth}{|l|X|X|}\hline
\rowcolor{popblueL}
\textbf{Domain / Measure} & \textbf{Key Terms (English / Français)} & \textbf{Examples / Exemples} \\ \hline
\cle{Vector-borne} & \emph{Vector} (vecteur : moustique, tique), \emph{Stagnant water} (eau stagnante), \emph{Breeding site} (gîte larvaire) & \textbf{Malaria} (paludisme), \textbf{Dengue}, \textbf{Yellow fever} (fièvre jaune) \\ \hline
\cle{Water-borne} & \emph{Contaminated source} (source contaminée), \emph{Latrine} (latrine), \emph{Oral rehydration} (réhydratation orale), \emph{Outbreak} (épidémie) & \textbf{Cholera} (choléra), \textbf{Typhoid} (fièvre typhoïde), \textbf{Hepatitis A} \\ \hline
\cle{Airborne} & \emph{Droplets} (gouttelettes), \emph{Ventilation} (aération), \emph{Quarantine} (quarantaine), \emph{Isolation} (isolement), \emph{Protocol} (protocole) & \textbf{Tuberculosis} (TB), \textbf{COVID-19}, \textbf{Meningitis} (méningite) \\ \hline
\cle{Prevention \& Control} & \emph{Vaccination} (vaccination), \emph{Sanitation} (assainissement), \emph{Hygiene} (hygiène), \emph{Handwashing} (lavage des mains), \emph{Surveillance} (surveillance) & \emph{Immunization campaign} (campagne de vaccination), \emph{WASH strategy} (stratégie EAH) \\ \hline
\end{tabularx}
\end{center}

\begin{definition}[Communicable Disease]
A \cle{communicable disease} is an illness caused by an infectious agent (bacterium, virus, parasite, fungus) that can be transmitted from one host to another, directly or indirectly (via vectors, water, air, surfaces). In French : \emph{Maladie transmissible}.
\end{definition}

\begin{savaistu}[La stratégie WASH au Cameroun]
La stratégie \textbf{WASH} (\textbf{W}ater, \textbf{S}anitation, \textbf{H}ygiene / Eau, Assainissement, Hygiène) est le pilier de la prévention du choléra dans les écoles camerounaises. Le programme \emph{« Écoles Amies de l'Hygiène »}, soutenu par l'UNICEF et le MINSANTE, impose la présence de points d'eau potable, de latrines séparées (filles/garçons) et de dispositifs de lavage des mains (\emph{handwashing stations}) à base de cendres ou de savon. En saison des pluies, la surveillance de la \cle{water quality} (qualité de l'eau) devient critique dans le Septentrion.
\end{savaistu}

\courssub{Listening Skills: Decoding Public Health Announcements}

In the \cle{Probatoire} and \cle{Baccalauréat technique}, listening exams often feature authentic audio: a CRTV news flash, a speech by the \emph{Regional Delegate of Public Health} (Délégué Régional de la Santé Publique), or a community health worker sensitising a market crowd.

\begin{experience}[Simulated Listening Task: Cholera Response in the North Region]
\textbf{Context:} Rainy season (July–October). Unprotected water points (\emph{puits non protégés}, \emph{marigots}) contaminated by flooding.
\textbf{Script (Audio Simulation):}
\begin{quote}
\emph{``Good morning. This is Dr. Aboubakar, Regional Delegate of Health for the North. We confirm a cholera outbreak in Garoua and Figuil health districts. As of today, 47 cases and 3 deaths are recorded. The Ministry urges all populations to strictly observe hygiene rules. You \textbf{must} drink only treated or boiled water. You \textbf{must} use latrines and stop open defecation. Food \textbf{must} be cooked thoroughly and eaten hot. Parents, you \textbf{should} vaccinate your children against cholera at the nearest health centre. Health workers \textbf{have to} report any suspected case immediately. Together, we can stop this epidemic.''}
\end{quote}
\textbf{Task:} Classify the measures into \textbf{Obligatory} (\emph{Must / Have to}) vs \textbf{Recommended} (\emph{Should / Ought to}).
\end{experience}

\begin{exoresolu}[Note-taking Grid: Obligation vs Recommendation]
\textbf{Instructions:} Complete the table based on the audio script above.
\begin{center}
\begin{tabularx}{\linewidth}{|l|X|X|}\hline
\rowcolor{popgreenL}
\textbf{Modality} & \textbf{Measure (Direct Speech)} & \textbf{Target Group} \\ \hline
\cle{Must} (Strong Obligation) & Drink only treated/boiled water & Populations \\ \hline
\cle{Must} & Use latrines; stop open defecation & Populations \\ \hline
\cle{Must} & Cook food thoroughly; eat hot & Populations \\ \hline
\cle{Have to} (External Obligation) & Report suspected cases immediately & Health workers \\ \hline
\cle{Should} (Advice) & Vaccinate children against cholera & Parents \\ \hline
\end{tabularx}
\end{center}
\textbf{Key Distinction:} \emph{Must} reflects the speaker's authority/urgency. \emph{Have to} reflects an external rule (professional protocol). \emph{Should} indicates strong advice for protection.
\end{exoresolu}

\courssub{Grammar Focus: Reported Speech (Discours Rapporté)}

When you write a report for your school journal, a safety logbook, or a technical file, you rarely quote people word-for-word (\cle{Direct Speech}). You summarise using \cle{Reported Speech} (Indirect Speech).

\begin{definition}[Reported Speech]
\cle{Reported Speech} reports what someone said without using their exact words. It typically involves a \textbf{reporting verb} (\emph{said, told, warned, advised, announced, confirmed}) in the past tense, followed by a \textbf{reported clause} introduced by \emph{that} (often omitted). The tense of the verb in the reported clause usually shifts \textbf{backwards} (\cle{backshift}).
\end{definition}

\begin{popfigure}
\begin{tikzpicture}[scale=0.7, every node/.style={font=\small}]
% Direct Speech Box
\node[draw, rectangle, fill=popblue!15, rounded corners, align=center, text width=5cm] (direct) at (-4,0) {
\textbf{Direct Speech}\\
\textquoteleft Wash your hands regularly,\textquoteright\\
\textbf{said} the nurse.
};
% Indirect Speech Box
\node[draw, rectangle, fill=popgreen!15, rounded corners, align=center, text width=5cm] (indirect) at (4,0) {
\textbf{Reported Speech}\\
The nurse \textbf{advised} us\\
\textbf{to wash} our hands regularly.
};
% Arrow
\draw[<->, thick, dashed, popdark] (direct.east) -- node[above, font=\small, popdark] {Backshift Rules} (indirect.west);
% Tense Shifts Box
\node[draw, rectangle, fill=popcream, rounded corners, align=left, text width=6cm, font=\small] (tense) at (0,-2.5) {
\textbf{Tense Shifts (Backshift):}\\
Present Simple $\rightarrow$ Past Simple\\
\emph{``I \textbf{am} tired''} $\rightarrow$ He said he \textbf{was} tired.\\
Present Perfect $\rightarrow$ Past Perfect\\
\emph{``Cases \textbf{have increased}''} $\rightarrow$ She reported cases \textbf{had increased}.\\
Will $\rightarrow$ Would\\
\emph{``It \textbf{will} rain''} $\rightarrow$ He said it \textbf{would} rain.\\
Can $\rightarrow$ Could\\
\emph{``We \textbf{can} help''} $\rightarrow$ They said they \textbf{could} help.\\
Must $\rightarrow$ Had to\\
\emph{``You \textbf{must} boil water''} $\rightarrow$ He said we \textbf{had to} boil water.
};
% Imperative Box
\node[draw, rectangle, fill=poppink!15, rounded corners, align=left, text width=6cm, font=\small] (imper) at (0,-5) {
\textbf{Imperative $\rightarrow$ Infinitive:}\\
\emph{``\textbf{Do} this''} $\rightarrow$ He told me \textbf{to do} this.\\
\emph{``\textbf{Don't} touch''} $\rightarrow$ She warned me \textbf{not to touch}.\\
\emph{Reporting verbs:} \textbf{advise, warn, order, urge, recommend}.
};
\end{tikzpicture}
\legende{Visual summary of Reported Speech transformations: Backshift rules for tenses and conversion of Imperatives to Infinitives.}
\end{popfigure}

\begin{propriete}[La Règle du Backshift (Concordance des Temps)]
Si le verbe introducteur (\emph{reporting verb}) est au \textbf{passé} (\emph{said, told, warned, announced}), les temps de la proposition rapportée reculent d'un cran :
\begin{itemize}
    \item \textbf{Present Simple} $\rightarrow$ \textbf{Past Simple}
    \item \textbf{Present Continuous} $\rightarrow$ \textbf{Past Continuous}
    \item \textbf{Present Perfect} $\rightarrow$ \textbf{Past Perfect}
    \item \textbf{Past Simple} $\rightarrow$ \textbf{Past Perfect}
    \item \textbf{Will} $\rightarrow$ \textbf{Would}
    \item \textbf{Can} $\rightarrow$ \textbf{Could}
    \item \textbf{May} $\rightarrow$ \textbf{Might}
    \item \textbf{Must} $\rightarrow$ \textbf{Had to} (Pas de \emph{musted} !)
\end{itemize}
\textbf{Exceptions (Pas de backshift) :}
\begin{enumerate}
    \item \cle{Vérités générales / Faits scientifiques} : \emph{``Water boils at 100°C''} $\rightarrow$ He said water \textbf{boils} at 100°C.
    \item \cle{Faits encore actuels / Futur immédiat} : \emph{``The meeting is tomorrow''} (dit aujourd'hui) $\rightarrow$ She said the meeting \textbf{is} tomorrow.
    \item \cle{Conditionnels 2 \& 3} : \emph{``If I were you, I'd go''} $\rightarrow$ He said if he \textbf{were} me, he \textbf{'d} go. (Déjà au passé).
\end{enumerate}
\end{propriete}

\begin{methode}[Transformer un ordre/conseil (Impératif) en Discours Rapporté]
Les consignes d'hygiène (affiches, radios) sont souvent à l'impératif. Suivez ces étapes :
\begin{enumerate}
    \item \textbf{Identifiez le verbe introducteur} selon la force : \emph{advise} (conseil), \emph{warn} (avertissement danger), \emph{order} (ordre strict), \emph{urge} (insistance), \emph{recommend} (recommandation officielle).
    \item \textbf{Mettez le verbe principal à l'infinitif} : \emph{to + base verbale}.
    \item \textbf{Gérez la négation} : \emph{not to + base verbale} (placez \emph{not} avant \emph{to}).
    \item \textbf{Adaptez les pronoms et déterminants} : \emph{I $\rightarrow$ he/she, you $\rightarrow$ me/us/them, my $\rightarrow$ his/her, this $\rightarrow$ that, these $\rightarrow$ those}.
    \item \textbf{Adaptez les marqueurs spatio-temporels} (voir tableau ci-dessous).
\end{enumerate}
\end{methode}

\begin{center}
\begin{tabular}{|l|l|}\hline
\rowcolor{poporangeL}
\textbf{Direct Speech (Marqueurs)} & \textbf{Reported Speech (Marqueurs)} \\ \hline
now & then / at that moment \\ \hline
today / tonight & that day / that night \\ \hline
yesterday & the day before / the previous day \\ \hline
tomorrow & the next day / the following day \\ \hline
last week / month & the previous week / month \\ \hline
next week / month & the following week / month \\ \hline
here & there \\ \hline
this / these & that / those \\ \hline
ago & before \\ \hline
\end{tabular}
\end{center}

\begin{attention}[Pièges Fréquents : Say vs Tell \& Pronoms]
\begin{itemize}
    \item \textbf{Say} vs \textbf{Tell} : \emph{Say} ne prend \textbf{JAMAIS} d'objet personne (COI) directement. \emph{Tell} \textbf{OBLIGE} un objet personne.
    \begin{itemize}
        \item \checkmark He \textbf{said} (that) he was sick. \quad \checkmark He \textbf{told me} (that) he was sick.
        \item $\times$ He \textbf{said me} he was sick. \quad $\times$ He \textbf{told} he was sick.
    \end{itemize}
    \item \textbf{Oubli du changement de pronom} : \emph{``I will wash \textbf{my} hands''} $\rightarrow$ He said \textbf{he} would wash \textbf{his} hands (pas \emph{my}).
    \item \textbf{Must $\rightarrow$ Had to} : Ne créez pas \emph{musted}. Pour l'absence d'obligation au passé : \emph{didn't have to}. Pour l'interdiction : \emph{must not} $\rightarrow$ \emph{warned not to} ou \emph{forbade from -ing}.
\end{itemize}
\end{attention}

\begin{exemplebox}[Chiffré : Rapport d'épidémie]
\textbf{Direct :} \emph{``Cholera cases \textbf{have increased} by 15\% \textbf{this} week,''} \textbf{reported} the doctor.
\textbf{Reported :} The doctor \textbf{reported that} cholera cases \textbf{had increased} by 15\% \textbf{that} week.
\begin{itemize}
    \item \emph{have increased} (Present Perfect) $\rightarrow$ \emph{had increased} (Past Perfect).
    \item \emph{this week} $\rightarrow$ \emph{that week}.
    \item Verbe introducteur : \emph{reported} (plus précis que \emph{said} pour un bulletin officiel).
\end{itemize}
\end{exemplebox}

\begin{exemplebox}[Variété des Verbes Introducteurs]
Ne vous limitez pas à \emph{said/told}. Choisissez le verbe qui rend la \textbf{fonction illocutoire} (l'intention) :
\begin{center}
\begin{tabular}{|l|l|l|}\hline
\rowcolor{poppurpleL}
\textbf{Direct Speech} & \textbf{Reporting Verb} & \textbf{Reported Structure} \\ \hline
``Get vaccinated!'' & \textbf{Urged / Advised} & He \textbf{urged us to get} vaccinated. \\ \hline
``Don't drink that water!'' & \textbf{Warned} & She \textbf{warned us not to drink} that water. \\ \hline
``We confirm the outbreak.'' & \textbf{Confirmed} & They \textbf{confirmed that} there \textbf{was} an outbreak. \\ \hline
``I announce free treatment.'' & \textbf{Announced} & The Minister \textbf{announced that} treatment \textbf{was} free. \\ \hline
``Please, use the latrine.'' & \textbf{Recommended / Asked} & He \textbf{recommended using} the latrine / \textbf{asked us to use} it. \\ \hline
\end{tabular}
\end{center}
\end{exemplebox}

\courssub{Writing Practice: From School Poster to Journal Report}

\begin{exoresolu}[Transformation Task: School Hygiene Poster $\rightarrow$ Written Report]
\textbf{Scenario:} You are the editor of the \emph{Technical High School Gazette}. You saw this poster on the wall. Write a short report (50–60 words) for the next edition.

\textbf{Poster Text (Direct Speech):}
\begin{quote}
\textbf{STOP CHOLERA!}
\textbf{Wash} your hands with soap before eating.
\textbf{Do not} eat cold food from street vendors.
\textbf{Drink} only treated water.
\textbf{Report} any diarrhoea case to the infirmary \textbf{immediately}.
\textit{— The School Nurse, Monday 14th October}
\end{quote}

\textbf{Model Report (Reported Speech):}
\begin{quote}
On Monday 14th October, the school nurse \textbf{warned} students about cholera risks. She \textbf{advised} them \textbf{to wash} their hands with soap before meals and \textbf{urged} them \textbf{not to eat} cold food from street vendors. She \textbf{insisted that} they \textbf{had to drink} only treated water. Finally, she \textbf{ordered} anyone with diarrhoea \textbf{to report} to the infirmary immediately.
\end{quote}

\textbf{Analysis of Transformations:}
\begin{enumerate}
    \item \emph{Wash} (Imperative) $\rightarrow$ \emph{advised ... to wash} (Infinitif).
    \item \emph{Do not eat} (Négatif) $\rightarrow$ \emph{urged ... not to eat}.
    \item \emph{Drink} (Imperatif fort) $\rightarrow$ \emph{insisted that ... had to drink} (Subjonctif/Modalité rapportée) ou \emph{ordered ... to drink}. Ici \emph{had to} rend l'obligation \emph{must} implicite.
    \item \emph{Report ... immediately} $\rightarrow$ \emph{ordered ... to report ... immediately} (adverbe conservé car relatif à l'action).
    \item \emph{Monday 14th October} $\rightarrow$ \emph{On Monday 14th October} (marqueur temporel ancré).
\end{enumerate}
\end{exoresolu}

\courssub{Exercices d'Application (Style Probatoire / Bac Tech)}

\begin{exercice}[Grammaire : Backshift \& Verbes Introducteurs]
\textbf{Turn the following direct statements into reported speech. Choose an appropriate reporting verb (warn, advise, confirm, announce, urge, order).}
\begin{enumerate}
    \item ``The cholera vaccine is available at the district hospital,'' the Director said.
    \item ``Do not swim in the river during the rainy season,'' the health worker told the children.
    \item ``We will launch a sanitation campaign next week,'' the Mayor announced.
    \item ``You must boil drinking water for at least 10 minutes,'' the nurse insisted.
    \item ``Malaria cases have dropped by 20\% this year,'' the report stated.
\end{enumerate}
\end{exercice}

\begin{exercice}[Vocabulaire \& Compréhension : Complétion]
\textbf{Complete the sentences with words from the vocabulary table (Vector-borne, Water-borne, Airborne, Outbreak, Quarantine, Sanitation, Vaccination, Latrine).}
\begin{enumerate}
    \item The \_\_\_\_\_\_\_\_ of cholera in the North Region was linked to flooded \_\_\_\_\_\_\_\_.
    \item Sleeping under treated mosquito nets prevents \_\_\_\_\_\_\_\_ diseases like malaria.
    \item TB patients must respect \_\_\_\_\_\_\_\_ to avoid infecting their family.
    \item Improving \_\_\_\_\_\_\_\_ (building toilets, waste management) stops \_\_\_\_\_\_\_\_ diseases.
    \item The \_\_\_\_\_\_\_\_ campaign against meningitis targets students in the meningitis belt.
\end{enumerate}
\end{exercice}

\begin{exercice}[Production Écrite : Rapport Technique]
\textbf{Context:} You are a \emph{HSE (Health, Safety, Environment) Technician} intern at a local agro-industrial company (ex: CDC, SOCAPALM, HEVECAM). The Site Doctor gave this briefing yesterday:
\emph{``Workers must wear masks in the dusty packaging area. They should wash hands before lunch. Anyone with a persistent cough must go to the clinic today. We will fumigate the dormitories tomorrow.''}
\textbf{Task:} Write the \textbf{Daily HSE Logbook Entry} (Reported Speech, formal style, ~60 words).
\end{exercice}

\courssub{Corrigés Détaillés}

\begin{corrige}[Exercice 1 : Grammaire]
\begin{enumerate}
    \item The Director \textbf{confirmed/announced} (that) the cholera vaccine \textbf{was} available at the district hospital. (\emph{is $\rightarrow$ was} ; vérité générale possible mais \emph{was} standard au passé).
    \item The health worker \textbf{warned} the children \textbf{not to swim} in the river during the rainy season. (\emph{Imperatif négatif $\rightarrow$ warned not to}).
    \item The Mayor \textbf{announced} (that) they \textbf{would launch} a sanitation campaign \textbf{the following week}. (\emph{will $\rightarrow$ would ; next week $\rightarrow$ the following week}).
    \item The nurse \textbf{insisted/urged} (that) they \textbf{had to boil} drinking water for at least 10 minutes. (\emph{must $\rightarrow$ had to}).
    \item The report \textbf{stated/reported} (that) malaria cases \textbf{had dropped} by 20\% \textbf{that year}. (\emph{have dropped $\rightarrow$ had dropped ; this year $\rightarrow$ that year}).
\end{enumerate}
\end{corrige}

\begin{corrige}[Exercice 2 : Vocabulaire]
\begin{enumerate}
    \item The \textbf{outbreak} of cholera in the North Region was linked to flooded \textbf{latrines} (or \emph{water points/sources}).
    \item Sleeping under treated mosquito nets prevents \textbf{vector-borne} diseases like malaria.
    \item TB patients must respect \textbf{quarantine} (or \emph{isolation}) to avoid infecting their family.
    \item Improving \textbf{sanitation} (building toilets, waste management) stops \textbf{water-borne} diseases.
    \item The \textbf{vaccination} campaign against meningitis targets students in the meningitis belt.
\end{enumerate}
\end{corrige}

\begin{corrige}[Exercice 3 : Production Écrite]
\textbf{Daily HSE Logbook — Date: [Today's Date]}
The Site Doctor \textbf{ordered} workers \textbf{to wear} masks in the dusty packaging area and \textbf{advised} them \textbf{to wash} their hands before lunch. He \textbf{warned} that anyone with a persistent cough \textbf{had to go} to the clinic \textbf{that day}. He \textbf{announced} that the dormitories \textbf{would be fumigated} \textbf{the next day}.
\textit{(Word count: ~55 words. Tenses backshifted. Pronouns/Time markers adapted. Passive used for 'fumigate' as agent unknown/irrelevant).}
\end{corrige}

\begin{aretenir}[Checklist Rédaction Rapport Technique (HSE / Journal Scolaire)]
Avant de rendre votre copie, vérifiez :
\begin{enumerate}
    \item \textbf{Verbe introducteur} varié et précis (\emph{warned, advised, ordered, confirmed}...).
    \item \textbf{Backshift} appliqué systématiquement (sauf exceptions).
    \item \textbf{Pronoms / Possessifs} adaptés (\emph{I $\rightarrow$ he, my $\rightarrow$ his}).
    \item \textbf{Marqueurs temps/lieu} adaptés (\emph{now $\rightarrow$ then, today $\rightarrow$ that day, tomorrow $\rightarrow$ the next day}).
    \item \textbf{Impératifs} $\rightarrow$ \emph{to + V} / \emph{not to + V}.
    \item \textbf{Registre} formel (pas de contractions \emph{didn't} $\rightarrow$ \emph{did not} si très formel, sinon toléré).
\end{enumerate}
\end{aretenir}

\coursec{Environmental Issues: Climate Change \& Waste Management (Reading \& Conditionals)}

\courssub{Transition from the Previous Section}
In the previous section, we learnt how to \textbf{report} what health experts say about preventing communicable diseases, using \textbf{Direct and Indirect Speech}. We moved from \emph{quoting} exact words to \emph{summarising} messages. Now we shift our focus from human health to \textbf{planetary health}. We will read about environmental challenges in Cameroon and around the world, and we will learn to express \textbf{conditions and consequences} using the \textbf{First and Second Conditionals}. This allows us to talk about \emph{real future plans} (If we act, things will improve) and \emph{imaginary situations} (If everyone cared, the world would be better).

\courssub{Vocabulary: Environmental Issues \& Waste Management}

\begin{definition}[Key Environmental Vocabulary]
The following words are essential for understanding texts on climate change and waste management in the Cameroonian context. Learn each word with its French gloss (\emph{glose française}).
\end{definition}

\begin{center}
\begin{tabularx}{\linewidth}{|l|X|l|}
\hline
\rowcolor{popblueL}
\textbf{English Term} & \textbf{Definition / Context} & \textbf{French Gloss} \\
\hline
Global warming & Long-term rise in the Earth's average temperature. & Réchauffement climatique \\
\hline
Greenhouse gases (GHG) & Gases (carbon dioxide, methane) trapping heat in the atmosphere. & Gaz à effet de serre \\
\hline
Deforestation & Clearing forests for agriculture or logging. & Déforestation \\
\hline
Erosion & Wearing away of topsoil by wind or rain. & Érosion \\
\hline
Desertification & Process by which fertile land becomes desert. & Désertification \\
\hline
Landfill & Site for burying waste (e.g., Nkolfoulou near Yaoundé). & Décharge (d'ordures) \\
\hline
Recycling & Converting waste into reusable material. & Recyclage \\
\hline
Composting & Decomposing organic waste into fertilizer. & Compostage \\
\hline
Biodegradable & Able to decay naturally without harming the environment. & Biodégradable \\
\hline
Single-use plastic & Plastic items used once then thrown away (bags, straws). & Plastique à usage unique \\
\hline
Circular economy & Economic system aimed at eliminating waste and reusing resources. & Économie circulaire \\
\hline
Carbon footprint & Total GHG emissions caused by an individual, event or product. & Empreinte carbone \\
\hline
\end{tabularx}
\end{center}

\begin{exemplebox}[Vocabulary in Context -- Cameroon]
\begin{itemize}
    \item \textbf{Rising sea levels} threaten coastal towns like \textbf{Kribi} and \textbf{Limbé}. (\emph{L'élévation du niveau de la mer menace...})
    \item \textbf{Desertification} advances in the \textbf{Far North Region} (Sahel zone), causing \textbf{food insecurity}. (\emph{La désertification progresse...})
    \item \textbf{HYSACAM} (Hygiène et Salubrité du Cameroun) collects household waste in Yaoundé and Douala and manages the \textbf{Nkolfoulou landfill}, but \textbf{recycling rates} remain low.
    \item Local \textbf{startups} turn \textbf{plastic waste} into \textbf{paving stones} (\emph{pavés}) -- a \textbf{circular economy} initiative.
\end{itemize}
\end{exemplebox}

\courssub{Reading: Environmental Challenges in Cameroon}

\begin{exemplebox}[Reading Text 1: Coastal Threats \& Northern Desertification]
\textbf{Cameroon's Climate Frontlines}

Cameroon is often called ``Africa in miniature'' because of its diverse geography. However, this diversity makes it vulnerable to climate change.

In the \textbf{Southwest} and \textbf{Littoral} regions, the coastal towns of \textbf{Limbé} and \textbf{Kribi} face \textbf{rising sea levels} and \textbf{coastal erosion}. The main causes are \textbf{global warming} (melting ice caps and the thermal expansion of seawater) and local \textbf{deforestation} of mangroves, which normally protect the shore. Consequences include the destruction of fishing infrastructure, the salinization of freshwater, and the displacement of communities.

In the \textbf{Far North} (Sahelian zone), \textbf{desertification} is accelerating. Irregular rainfall, overgrazing, and cutting trees for firewood cause severe \textbf{soil erosion}. This leads to \textbf{food insecurity} and forces \textbf{climate migration} --- people leaving villages for towns like Maroua, or even for Yaoundé.

\textit{Source: Adapted from MINEPDED and UNESCO materials on climate change education.}
\end{exemplebox}

\begin{exemplebox}[Reading Text 2: Waste Management in Yaoundé -- The Nkolfoulou Challenge]
\textbf{From Waste to Resource: The Yaoundé Case}

Yaoundé, the political capital, produces well over \textbf{1,500 tons of waste every day}. Most of it ends up at the \textbf{Nkolfoulou landfill}, north of the city. The composition of this waste is revealing: the largest portion is \textbf{organic} (compostable), followed by \textbf{plastic} (often \textbf{single-use}), paper, and metal.

Currently, only a small fraction is formally recycled. The rest pollutes waterways such as the \textbf{Mfoundi River} and releases \textbf{methane}, a potent \textbf{greenhouse gas}. However, change is coming. The 2014 regulation against \textbf{non-biodegradable plastic packaging} (less than 60 microns thick) was an important legal step. Today, youth-led \textbf{startups} collect plastic bottles and transform them into \textbf{ecological paving stones} or \textbf{polyester fibre} for the textile industry. This \textbf{circular economy} approach creates jobs and reduces the \textbf{carbon footprint}.

\textit{Source: Adapted from HYSACAM data and Cameroon Tribune articles.}
\end{exemplebox}

\begin{exercice}[Reading Comprehension Check]
Answer the questions based on the texts above. Use full sentences.
\begin{enumerate}
    \item Identify two causes of coastal erosion in Limbé/Kribi mentioned in Text 1.
    \item What is the link between desertification in the Far North and migration?
    \item According to Text 2, where does most of Yaoundé's waste end up?
    \item Name two products made from recycled plastic by Cameroonian startups.
    \item \textbf{Vocabulary:} Find words in the texts meaning:
    \begin{itemize}
        \item[a)] Making land into desert (Text 1)
        \item[b)] Able to rot naturally (Text 2)
        \item[c)] Gas released by rotting waste (Text 2)
    \end{itemize}
\end{enumerate}
\end{exercice}

\begin{corrige}[Exercice Reading Comprehension]
\begin{enumerate}
    \item Global warming (melting ice caps, thermal expansion of seawater) and local deforestation of mangroves.
    \item Desertification causes soil erosion and food insecurity, which forces people to leave their villages and migrate to towns (climate migration).
    \item Most of it ends up at the Nkolfoulou landfill, north of the city.
    \item Ecological paving stones and polyester fibre for the textile industry.
    \item a) Desertification; b) Biodegradable; c) Methane.
\end{enumerate}
\end{corrige}

\courssub{Grammar: First \& Second Conditionals}

\begin{definition}[Conditionals: Expressing Cause \& Effect under Condition]
Conditionals are sentences with two parts: an \textbf{If-clause} (the condition) and a \textbf{main clause} (the result). They express \textbf{probability} and \textbf{reality}.
\begin{itemize}
    \item \textbf{Type 1 (Real Future):} The condition is \textbf{possible} and \textbf{likely} to happen. We talk about \emph{real plans, warnings, promises}.
    \item \textbf{Type 2 (Unreal Present/Future):} The condition is \textbf{imaginary}, \textbf{unlikely}, or \textbf{contrary to current reality}. We talk about \emph{dreams, advice, hypothetical scenarios}.
\end{itemize}
\end{definition}

\begin{propriete}[Structures -- First vs Second Conditional]
\begin{center}
\begin{tabularx}{\linewidth}{|l|X|X|}
\hline
\rowcolor{popblueL}
\textbf{Type} & \textbf{If-Clause (Condition)} & \textbf{Main Clause (Result)} \\
\hline
\textbf{Type 1} (Real Future) & If + \textbf{Present Simple} & Subject + \textbf{will / can / must / may} + \textbf{base verb} \\
& \textit{If we recycle more,} & \textit{we \textbf{will reduce} pollution.} \\
\hline
\textbf{Type 2} (Unreal / Hypothetical) & If + \textbf{Past Simple} & Subject + \textbf{would / could / might} + \textbf{base verb} \\
& \textit{If everyone \textbf{composted},} & \textit{landfills \textbf{would shrink}.} \\
& \textit{(Note: \textbf{were} for all persons: If I \textbf{were} you...)} & \\
\hline
\end{tabularx}
\end{center}
\end{propriete}

\begin{methode}[Choosing the Right Conditional: 3-Step Decision]
Follow these steps to decide between Type 1 and Type 2:
\begin{enumerate}
    \item \textbf{Analyse Reality:} Is the condition \emph{really possible} in the near future? (e.g., ``If the council \textbf{organizes} a clean-up next month...'') $\rightarrow$ \textbf{Type 1}.
    \item \textbf{Analyse Imagination:} Is the condition \emph{imaginary, hypothetical, or false now}? (e.g., ``If I \textbf{were} the Minister of Environment...'' or ``If everyone \textbf{sorted} their waste today...'' --- but they don't) $\rightarrow$ \textbf{Type 2}.
    \item \textbf{Check Verb Forms:}
    \begin{itemize}
        \item Type 1: \textbf{No `will' in the If-clause!} (\emph{Incorrect: *If it will rain...} $\rightarrow$ Correct: \textbf{If it rains...})
        \item Type 2: Use \textbf{were} for \emph{I, he, she, it} in formal writing (\emph{If I \textbf{were} you...}).
    \end{itemize}
\end{enumerate}
\end{methode}

\begin{attention}[Common Traps -- Bac Technique Focus]
\begin{itemize}
    \item \textbf{Trap 1: Future in the If-clause.} \textbf{Never} use \emph{will/would} after \emph{If}.
    \begin{itemize}
        \item Wrong: *If the rain \textbf{will stop}, we \textbf{will plant} trees.
        \item Right: If the rain \textbf{stops}, we \textbf{will plant} trees. (Type 1)
    \end{itemize}
    \item \textbf{Trap 2: `Was' vs `Were' in Type 2.} In formal English (exams), always use \textbf{were} for \emph{I/he/she/it}.
    \begin{itemize}
        \item Informal: If I \textbf{was} you...
        \item Formal / Exam: If I \textbf{were} you, I \textbf{would} recycle.
    \end{itemize}
    \item \textbf{Trap 3: Confusing Type 2 and Type 3.} Type 2 = present/future hypothetical (If + Past Simple, would + base verb). Type 3 = past hypothetical (If + Past Perfect, would have + past participle). \textbf{Type 3 is not required in this unit.} Do not use \emph{would have} here.
\end{itemize}
\end{attention}

\begin{exemplebox}[Nuance in Context: Cameroon Policy]
Compare these two sentences about the plastic packaging regulation:
\begin{enumerate}
    \item \textbf{Type 1 (Real / Political Will):} ``If the government \textbf{enforces} the 2014 regulation strictly, plastic pollution \textbf{will decrease}.''
    \begin{itemize}
        \item \textit{Meaning: Enforcement is a real possibility. The speaker believes it can happen.}
    \end{itemize}
    \item \textbf{Type 2 (Hypothetical / Criticism):} ``If the government \textbf{enforced} the regulation strictly, plastic pollution \textbf{would decrease}.''
    \begin{itemize}
        \item \textit{Meaning: The government is NOT enforcing it strictly now. This is a dream or a piece of advice.}
    \end{itemize}
\end{enumerate}
\end{exemplebox}

\begin{exoresolu}[Grammar Application: Conditionals \& Data]
\textbf{Context:} In Yaoundé, only a small percentage of household waste is formally recycled. Imagine the current recycling rate is about 10\% of the roughly 550,000 tons produced each year.
\textbf{Task:} Write two sentences using this data: one Type 1 (realistic plan), one Type 2 (ideal scenario).
\tcblower
\textbf{Solution:}
\begin{enumerate}
    \item \textbf{Type 1 (Real Plan):} ``If the recycling rate \textbf{reaches} 20\% next year, the city \textbf{will recycle} about 110,000 tons of waste instead of 55,000.''
    \begin{itemize}
        \item \textit{Reasoning: Reaching 20\% is a realistic target (doubling the current rate). Present Simple (\textbf{reaches}) + \textbf{will} + base verb. Calculation: $550\,000 \times 20\% = 110\,000$ tons.}
    \end{itemize}
    \item \textbf{Type 2 (Ideal Scenario):} ``If the recycling rate \textbf{reached} 50\%, about 275,000 tons of waste \textbf{would be} recycled every year.''
    \begin{itemize}
        \item \textit{Reasoning: 50\% is currently imaginary/utopian. Past Simple (\textbf{reached}) + \textbf{would be} + past participle (passive). Calculation: $550\,000 \times 50\% = 275\,000$ tons.}
    \end{itemize}
\end{enumerate}
\end{exoresolu}

\courssub{Data Visualization: Waste Composition \& Conditionals}

\begin{popfigure}
\begin{center}
\begin{tikzpicture}
\begin{axis}[
    ybar,
    bar width=0.7cm,
    ymin=0, ymax=60,
    symbolic x coords={Organic, Plastic, Paper, Metal, Other},
    xtick=data,
    ylabel={Percentage (\%)},
    xlabel={Waste category},
    title={Estimated composition of household waste in Yaoundé},
    nodes near coords,
    every node near coord/.append style={font=\small},
    tick label style={font=\small},
    label style={font=\small},
    width=0.85\linewidth,
    height=6cm
]
\addplot[fill=popgreen!70, draw=popdark] coordinates {(Organic,50) (Plastic,20) (Paper,12) (Metal,8) (Other,10)};
\end{axis}
\end{tikzpicture}
\end{center}
\legende{Figure 1: Estimated composition of household waste in Yaoundé (indicative values, in \%). Organic waste (about 50\%) represents a huge potential for composting --- \emph{If everyone composted, landfill volume would drop} (Type 2). Plastic (about 20\%) requires sorting to boost recycling --- \emph{If households sort their waste, recycling rates will increase} (Type 1).}
\end{popfigure}

\begin{attention}[Reading a Bar Chart in English]
When you describe a chart at the Bac technique, always:
\begin{enumerate}
    \item Name the \textbf{subject} and the \textbf{unit} (``The chart shows the composition of household waste, in per cent.'').
    \item Identify the \textbf{highest} and \textbf{lowest} values (use the \textbf{superlative}: ``Organic waste is \textbf{the largest} category.'').
    \item \textbf{Compare} categories (``Plastic is \textbf{twice as high as} metal.'').
    \item Draw a \textbf{conclusion} linked to the lesson (``This is why composting should be encouraged.'').
\end{enumerate}
\end{attention}

\courssub{Grammar Extension: Advanced Adjective Structures}

The syllabus requires mastery of degree structures to emphasise environmental urgency.

\begin{propriete}[Advanced Structures with Adjectives]
\begin{center}
\begin{tabularx}{\linewidth}{|l|X|X|}
\hline
\rowcolor{popblueL}
\textbf{Structure} & \textbf{Form} & \textbf{Meaning / Example (Environment)} \\
\hline
\textbf{How + Adjective} & \textbf{How} + Adj + subject + verb! & Exclamation / strong feeling. \textit{How \textbf{serious} the erosion in Limbé is!} \\
\hline
\textbf{Too + Adj + To + V} & \textbf{too} + Adj + \textbf{to} + base verb & Negative consequence: excess prevents action. \textit{The landfill is \textbf{too full to accept} more trucks.} \\
\hline
\textbf{Adj + Enough + To + V} & Adj + \textbf{enough} + \textbf{to} + base verb & Positive sufficiency. \textit{The compost is \textbf{rich enough to fertilize} the fields.} \\
\hline
\textbf{So + Adj + That} & \textbf{so} + Adj + \textbf{that} + clause & Emphasis on the result. \textit{The heat was \textbf{so intense that} crops failed.} \\
\hline
\textbf{The Superlative} & \textbf{the most} + Adj / \textbf{the} + Adj-\textbf{est} & Comparing three or more items (extreme degree). \textit{Climate change is \textbf{the biggest} threat to the Sahel.} \\
\hline
\end{tabularx}
\end{center}
\end{propriete}

\begin{attention}[Degree Structures: Common Errors]
\begin{itemize}
    \item \textbf{Too vs Very:} \emph{Too} implies a \textbf{negative result} (something cannot be done). \emph{Very} is just intensity.
    \begin{itemize}
        \item Correct: The water is \textbf{too polluted to drink}. (Result: we cannot drink it.)
        \item Illogical: *The water is \textbf{too clean to drink}.
    \end{itemize}
    \item \textbf{Position of Enough:} the adjective comes \textbf{before} \emph{enough} (\emph{strong enough}, NOT *\emph{enough strong}). But \emph{enough} comes \textbf{before} a noun: \emph{enough bins}, \emph{enough money}.
    \item \textbf{Superlative + Present Perfect:} often used with \emph{ever}: \textit{``The hottest year \textbf{we have ever recorded}.''}
    \item \textbf{Irregular superlatives:} good $\rightarrow$ \textbf{the best}; bad $\rightarrow$ \textbf{the worst}; far $\rightarrow$ \textbf{the farthest/furthest}.
\end{itemize}
\end{attention}

\begin{exercice}[Practice: Transform Using the Structures]
Rewrite the sentences using the word(s) in brackets.
\begin{enumerate}
    \item The deforestation rate is alarming. (\textbf{How})
    \item The plastic bags are very thin. They tear easily. (\textbf{Too ... to})
    \item The recycling plant is sufficiently modern. It can process metal. (\textbf{Enough ... to})
    \item Climate change is a bigger threat than pollution and erosion. (\textbf{The biggest})
    \item The drought was extremely severe. The harvests were lost. (\textbf{So ... that})
\end{enumerate}
\end{exercice}

\begin{corrige}[Exercice Degree Structures]
\begin{enumerate}
    \item \textbf{How alarming} the deforestation rate is!
    \item The plastic bags are \textbf{too thin to be reused}. (or: \textbf{too thin to resist} tearing.)
    \item The recycling plant is \textbf{modern enough to process} metal.
    \item Climate change is \textbf{the biggest threat} (of the three / of all).
    \item The drought was \textbf{so severe that} the harvests were lost.
\end{enumerate}
\end{corrige}

\courssub{Guided Writing: ``Solutions for My Neighbourhood''}

\begin{methode}[Writing a ``Problem--Solution'' Paragraph Using Conditionals]
\textbf{Goal:} Write a coherent paragraph (about 100--150 words) proposing solutions for a local environmental issue, mixing Type 1 and Type 2 conditionals.

\textbf{Step 1: Identify the Problem (Observation).}
Choose a specific issue: \textit{plastic litter in the streets, blocked gutters causing floods, smoke from burning trash, lack of bins.}

\textbf{Step 2: Propose Real Actions (Type 1).}
Use \emph{If + Present Simple, ... will/can/must} for things the community or the council \textbf{can actually do}.
\textit{Ex: ``If the council \textbf{places} more bins, people \textbf{will stop} littering.''}

\textbf{Step 3: Express an Ideal Vision / Advice (Type 2).}
Use \emph{If + Past Simple, ... would/could} for dreams or advice to citizens.
\textit{Ex: ``If everyone \textbf{composted} kitchen waste, our gardens \textbf{would be} greener.''}

\textbf{Step 4: Add Emphasis (Degree Structures).}
Use \emph{How / Too / Enough / Superlative}.
\textit{Ex: ``It is \textbf{too dangerous to ignore} this issue.'' / ``\textbf{How clean} our streets \textbf{could be}!''}

\textbf{Step 5: Conclude with a Call to Action.}
\textit{Ex: ``Let's start sorting today.''}
\end{methode}

\begin{exoresolu}[Model Paragraph: ``Cleaning Up Mvog-Ada Quarter'']
\textbf{Topic:} Waste management in a Yaoundé neighbourhood.
\tcblower
\textit{In Mvog-Ada quarter, blocked gutters cause severe floods every rainy season. \textbf{How serious} the situation becomes when plastic bottles clog the drains! If the local council \textbf{organizes} monthly clean-up campaigns (Type 1), the water \textbf{will flow} freely. If residents \textbf{stop} throwing bags into the stream (Type 1), the risk of cholera \textbf{will decrease}. However, habits are hard to change. \textbf{If I were} the quarter chief (Type 2), I \textbf{would fine} anyone dumping waste illegally. \textbf{If every household composted} its organic peelings (Type 2), the amount of trash \textbf{would drop} by half. The Nkolfoulou landfill is \textbf{too full to receive} much more waste. We need a \textbf{circular economy} approach now. \textbf{Let us act before the next rainy season.}}

\textbf{Analysis of the Model:}
\begin{itemize}
    \item \textbf{Type 1 (Real):} \textit{If the council organizes... will flow / If residents stop... will decrease.}
    \item \textbf{Type 2 (Hypothetical / Advice):} \textit{If I were... would fine / If every household composted... would drop.}
    \item \textbf{Degree:} \textit{How serious... / too full to receive.}
    \item \textbf{Vocabulary:} \textit{gutters, clog, cholera, composted, circular economy.}
\end{itemize}
\end{exoresolu}

\begin{exercice}[Writing Task -- Bac Technique Practice]
\textbf{Instruction:} Write a composition of \textbf{150--180 words} on the following topic:
\textit{``Your school environment suffers from plastic pollution and erosion. Write a letter to the Principal proposing concrete solutions. Use \textbf{First and Second Conditionals}, \textbf{Too/Enough/How} structures, and \textbf{environmental vocabulary} from this lesson.''}

\textbf{Suggested Plan:}
\begin{enumerate}
    \item \textbf{Salutation \& Problem Statement:} describe the issue (plastic, erosion) --- use \emph{How / So...that}.
    \item \textbf{Realistic Proposals (Type 1):} bins, an environment club, tree planting --- \emph{If we plant trees, we will stop erosion.}
    \item \textbf{Ideal Changes / Advice (Type 2):} behaviour change, banning single-use plastics --- \emph{If students brought reusable bottles, we would reduce waste.}
    \item \textbf{Urgency (Too/Enough):} \emph{The situation is too urgent to ignore.}
    \item \textbf{Closing:} formal sign-off (Yours faithfully, ...).
\end{enumerate}
\end{exercice}

\begin{savaistu}[Did You Know? Cameroon Green Initiatives]
\begin{itemize}
    \item \textbf{Plastic Regulation (2014):} Cameroon adopted a regulation against \textbf{non-biodegradable plastic packaging} (less than 60 microns thick). Enforcement remains the main challenge.
    \item \textbf{Recycling Startups:} In Douala and Yaoundé, young entrepreneurs collect plastic bottles, shred them, and produce \textbf{ecological paving stones} or \textbf{polyester fibre} for the textile industry --- the \textbf{circular economy} in action!
    \item \textbf{Sorting at Source:} Some councils in Yaoundé have pilot projects for \textbf{sorting at source} (\emph{tri sélectif}) with coloured bins (green = organic, yellow = plastic, blue = paper).
    \item \textbf{The Great Green Wall:} Cameroon participates in the African Union initiative to combat \textbf{desertification} in the Sahel by planting a belt of trees across the continent.
\end{itemize}
\end{savaistu}

\courssub{Integration \& Remediation Checklist}
Before moving to the next section, make sure you can:
\begin{itemize}
    \item[$\square$] Explain \textbf{global warming}, \textbf{deforestation}, \textbf{landfill} and \textbf{circular economy} in English.
    \item[$\square$] Distinguish \textbf{Type 1} (Real: \emph{If + Present Simple $\rightarrow$ will}) from \textbf{Type 2} (Unreal: \emph{If + Past Simple $\rightarrow$ would}).
    \item[$\square$] Avoid putting \emph{will/would} in the \emph{If}-clause.
    \item[$\square$] Use \textbf{were} for all subjects in Type 2 (\emph{If I were...}).
    \item[$\square$] Apply \textbf{Too...To}, \textbf{Enough...To}, \textbf{How + Adjective} and \textbf{superlatives} correctly.
    \item[$\square$] Read and comment on a simple bar chart using comparatives and superlatives.
    \item[$\square$] Write a structured paragraph mixing both conditionals about a local environmental issue.
\end{itemize}

\coursec{Advanced Adjective Structures : Degree \& Comparison (How + Adj, Too...To, Superlatives)}

In Section 3, we read texts about climate change, global warming and waste management, and we practised the first and second conditionals (\emph{If we recycle more, the city will be cleaner}). To write and speak about these environmental problems convincingly, we now need a richer toolbox to express \cle{degree} (degré) and \cle{comparison} (comparaison): how hot? too polluted? aware enough? This section gives you exactly that toolbox. You will learn to say not only \emph{The situation is serious} but also \emph{How serious is the situation?}, \emph{The water is too polluted to drink}, and \emph{The youth are aware enough to act}. These structures are frequent at the Probatoire and Baccalauréat technique, both in reading texts and in the writing paper.

\courssub{Quick Revision: Comparatives and Superlatives}

Before the new structures, let us fix the basics. Adjectives express qualities; comparatives and superlatives express \emph{degrees} of these qualities.

\begin{definition}[Comparative and Superlative]
The \cle{comparative} compares two elements (\emph{hotter than, more dangerous than}). The \cle{superlative} designates the element with the highest degree in a group (\emph{the hottest, the most polluted}).
\end{definition}

\begin{center}
\begin{tabularx}{\linewidth}{|l|X|X|X|}
\hline
\rowcolor{popblueL} \textbf{Adjective type} & \textbf{Base} & \textbf{Comparative} & \textbf{Superlative} \\
\hline
Short (1 syllable) & hot & hotter than & the hottest \\
\hline
Short ending in -y & dirty & dirtier than & the dirtiest \\
\hline
Long (2+ syllables) & dangerous & more dangerous than & the most dangerous \\
\hline
Long & polluted & more polluted than & the most polluted \\
\hline
Irregular & good / bad & better / worse & the best / the worst \\
\hline
\end{tabularx}
\end{center}

\begin{exemplebox}[Environmental context]
\begin{itemize}
\item March is \textbf{hotter than} December in Maroua.
\item Douala is \textbf{more polluted than} Bafoussam because of traffic and industry.
\item The Wouri river is one of \textbf{the most polluted} rivers in the region.
\item Tramadol abuse is \textbf{the most dangerous} drug problem among young workers today.
\end{itemize}
\end{exemplebox}

\begin{attention}[Common errors]
Never mix the two forms: *\emph{more hotter} or *\emph{the most dirtiest} are wrong. Double the final consonant of short adjectives after a single vowel: hot $\rightarrow$ \textbf{hotter}, big $\rightarrow$ \textbf{bigger}. And always use \textbf{than} after a comparative, never \emph{that} or \emph{as}.
\end{attention}

\courssub{How + Adjective: Exclamations and Indirect Questions}

English uses \cle{How + adjective} in two situations that students often confuse.

\begin{definition}[How + adjective]
\textbf{How + adjective} introduces either an \cle{exclamation} (\emph{How beautiful!}) or an \cle{indirect interrogative clause} (\emph{I wonder how difficult it is.}).
\end{definition}

\textbf{1. The exclamation.} Word order: \emph{How + adjective + subject + verb!}

\begin{itemize}
\item \emph{How hot it is today!} (Qu'il fait chaud aujourd'hui !)
\item \emph{How serious the situation is!}
\item \emph{How dirty our school yard is!}
\end{itemize}

\textbf{2. The indirect question.} After verbs like \emph{know, realize, wonder, see, understand, ask}, the word order stays \emph{affirmative}: \emph{how + adjective + subject + verb} (no inversion, no question mark).

\begin{itemize}
\item Direct question: \emph{How serious is the situation?}
\item Indirect question: \emph{I realize \textbf{how serious the situation is}.}
\item \emph{We know \textbf{how dangerous} smoking \textbf{is}.}
\item \emph{The doctor explained \textbf{how harmful} tramadol \textbf{is} for the liver.}
\end{itemize}

\begin{attention}[No inversion in indirect questions]
Wrong: *\emph{I realize how serious is the situation.} Correct: \emph{I realize how serious the situation \textbf{is}.} In an indirect question, the subject comes before the verb, exactly like in a normal affirmative sentence.
\end{attention}

\courssub{Too + Adjective + To + Infinitive: Excess and Negative Result}

Imagine a dustbin in the school yard. It is full of metal scraps. A student tries to lift it and fails. How do we describe this? \emph{The waste is \textbf{too heavy to lift}.} The structure \cle{too + adjective + to + infinitive} expresses an \textbf{excess} (excès) that produces a \textbf{negative or impossible result}.

\begin{propriete}[Equivalence of Too...to and So...that...not]
\emph{Too + adjective + to + infinitive} (excess, negative meaning) is equivalent to \emph{so + adjective + that + subject + cannot/could not + verb} (consequence with negation):
\begin{center}
The water is \textbf{too polluted to drink}. $\Longleftrightarrow$ The water is \textbf{so polluted that} we \textbf{cannot drink} it.
\end{center}
\end{propriete}

\begin{exemplebox}[A numerical example from North Cameroon]
The average temperature in Maroua reaches \SI{38}{\celsius} in March.
\begin{itemize}
\item \emph{It is \textbf{too hot to work} outside at noon.}
\item \emph{It is \textbf{so hot that} workers \textbf{cannot work} outside between 12 p.m. and 3 p.m.}
\end{itemize}
Both sentences carry the same message: the heat is excessive, so the activity is impossible.
\end{exemplebox}

\begin{methode}[Transforming Too...to into So...that]
\begin{enumerate}
\item Keep the same subject and the same adjective.
\item Replace \emph{too + adjective + to + verb} by \emph{so + adjective + that + subject + cannot/could not + verb}.
\item Check that the \emph{that}-clause contains the negation (\emph{cannot, could not, do not}): this negation is the hidden meaning of \emph{too}.
\end{enumerate}
Example: \emph{The smoke is too thick to breathe.} $\rightarrow$ Step 1: subject \emph{the smoke}, adjective \emph{thick}. Step 2: \emph{The smoke is so thick that we cannot breathe.} Step 3: negation \emph{cannot} present. Correct.
\end{methode}

\begin{attention}[Too is negative, very is neutral]
A very frequent error is to confuse \textbf{too} (negative, excessive) with \textbf{very} (neutral intensity). *\emph{He is too strong to lift it} means \emph{He is so strong that he CANNOT lift it} — absurd! Say \emph{He is \textbf{strong enough to} lift it} or \emph{He is \textbf{very} strong}. Use \emph{too} only when the result is impossible or undesirable: \emph{The air in Douala is too polluted for morning sports.}
\end{attention}

\courssub{Adjective + Enough + To + Infinitive: Sufficiency and Positive Result}

The opposite of \emph{too} is \cle{enough}. It expresses \textbf{sufficiency} (suffisance) and a \textbf{possible} result: \emph{The youth are \textbf{aware enough to act}.} Note the position: \emph{enough} comes \textbf{after} the adjective.

\begin{propriete}[Position of enough]
\emph{Enough} follows the adjective (\emph{adjective + enough}: \emph{strong enough, clean enough}) but precedes the noun (\emph{enough + noun}: \emph{enough time, enough bins, enough information}).
\end{propriete}

\begin{exemplebox}[Contextual examples]
\begin{itemize}
\item \emph{The students are \textbf{responsible enough to} sort the waste.} (adjective + enough)
\item \emph{We have \textbf{enough bins} to separate plastic from paper.} (enough + noun)
\item \emph{The well is \textbf{clean enough to} provide drinking water.}
\item \emph{Is the clinic \textbf{big enough to} receive all the patients?}
\end{itemize}
\end{exemplebox}

\begin{attention}[No to without an infinitive]
Do not add \emph{to} after \emph{enough} when there is no infinitive verb: \emph{We have enough time} (not *\emph{enough to time}). And never place \emph{enough} before the adjective: *\emph{enough strong} is wrong; say \emph{strong \textbf{enough}}.
\end{attention}

\courssub{So + Adjective + That + Clause: Consequence}

We have already met \emph{so...that} as the equivalent of \emph{too...to}. It is a general structure of \cle{consequence} (conséquence): \emph{so + adjective + that + clause}.

\begin{itemize}
\item \emph{Climate change is \textbf{so rapid that} adaptation is difficult.}
\item \emph{The flood was \textbf{so violent that} the bridge collapsed.}
\item \emph{Tramadol is \textbf{so addictive that} many young workers cannot stop.}
\end{itemize}

The figure below summarises the whole system of degree structures.

\begin{popfigure}
\begin{center}
\begin{tikzpicture}[scale=0.8]
\node[draw, rounded corners, fill=popgold!25] (center) at (0,0) {Structure \textbf{Too ... to}};
\node[draw, rounded corners, fill=poppink!25] (neg) at (-3,-2) {Meaning: \textbf{Impossible} / Negative result};
\node[draw, rounded corners, fill=popgreen!20] (pos) at (3,-2) {Equivalent: \textbf{So ... that ... not}};
\node[draw, rounded corners, fill=popblue!15] (ex1) at (-3,-4) {Ex: The waste is \textbf{too heavy} \textbf{to lift}.};
\node[draw, rounded corners, fill=popblue!15] (ex2) at (3,-4) {Ex: The waste is \textbf{so heavy that} we \textbf{cannot lift} it.};
\draw[->] (center) -- (neg);
\draw[->] (center) -- (pos);
\draw[->] (neg) -- (ex1);
\draw[->] (pos) -- (ex2);
\node[align=center, font=\small] at (0,2) {\textbf{How + Adj} (Exclamation/Indirect Question)\\ \textit{How hot it is today!} / \textit{We know \textbf{how dangerous} drugs are.}};
\end{tikzpicture}
\end{center}
\legende{Mind map of degree structures: \emph{too...to} (excess, negative result) and its equivalent \emph{so...that...not}; above, the two uses of \emph{how + adjective}.}
\end{popfigure}

\begin{savaistu}[How come?]
In informal spoken English, \emph{How come?} replaces \emph{Why?} but keeps the affirmative word order: \emph{How come you are late?} (Pourquoi es-tu en retard ?). Compare: \emph{Why are you late?} (inversion) versus \emph{How come you are late?} (no inversion). Very useful for the Speaking paper!
\end{savaistu}

\courssub{Solved Exercise and Guided Practice}

\begin{exoresolu}[Rewriting a report sentence]
Énoncé — Rewrite this sentence from an environmental report in two different ways: \emph{The air in Douala is very polluted; people with asthma cannot go out during traffic peaks.} Use (a) \emph{too...to} and (b) \emph{so...that}.
\tcblower
Solution détaillée.
(a) Identify the adjective (\emph{polluted}) and the impossible action (\emph{go out}): \emph{The air in Douala is \textbf{too polluted for} asthmatic people \textbf{to go out} during traffic peaks.}
(b) Apply the method: keep the subject, use \emph{so + adjective + that} and put the negation in the that-clause: \emph{The air in Douala is \textbf{so polluted that} asthmatic people \textbf{cannot go out} during traffic peaks.}
Both sentences are grammatically correct and stylistically richer than the original.
\end{exoresolu}

\begin{exoresolu}[How + adjective]
Énoncé — Transform the direct question into an indirect question after \emph{The journalist wants to know...}: \emph{How dangerous is tramadol for young workers?}
\tcblower
Solution détaillée. In an indirect question there is no inversion and no question mark. The subject (\emph{tramadol}) comes before the verb (\emph{is}): \emph{The journalist wants to know \textbf{how dangerous tramadol is} for young workers.}
\end{exoresolu}

\begin{exercice}[Transformation practice]
Rewrite each sentence using the structure in brackets.
\begin{enumerate}
\item The school yard is very dirty; we cannot play football there. (\emph{too...to})
\item The students are very motivated; they can organise a clean-up day. (\emph{adjective + enough + to})
\item The heat in the Far North is extreme; farmers start work at 5 a.m. (\emph{so...that})
\item How serious is global warming? The minister wonders... (indirect question)
\item The rubbish heap is very high; the children cannot climb over it safely. (\emph{too...to} then \emph{so...that...not})
\end{enumerate}
\end{exercice}

\begin{corrige}[Exercice : Transformation practice]
\begin{enumerate}
\item \emph{The school yard is \textbf{too dirty to play} football there.} (or: \emph{...too dirty for us to play football there.})
\item \emph{The students are \textbf{motivated enough to organise} a clean-up day.}
\item \emph{The heat in the Far North is \textbf{so extreme that} farmers start work at 5 a.m.}
\item \emph{The minister wonders \textbf{how serious global warming is}.} (no inversion)
\item \emph{The rubbish heap is \textbf{too high for} the children \textbf{to climb} over safely.} / \emph{The rubbish heap is \textbf{so high that} the children \textbf{cannot climb} over it safely.}
\end{enumerate}
\end{corrige}

\begin{exercice}[Contextual application: describe your environment]
Write five sentences about the following situations, using a different degree structure each time (comparative, superlative, \emph{how + adjective}, \emph{too...to}, \emph{enough...to}):
\begin{enumerate}
\item the state of your school yard after the rainy season;
\item the air quality in Douala or Yaoundé during rush hour;
\item the danger of tramadol for young apprentices;
\item the temperature in your region in March;
\item the motivation of your classmates to protect the environment.
\end{enumerate}
\end{exercice}

\begin{corrige}[Exercice : Contextual application (sample answers)]
\begin{enumerate}
\item \emph{After the rainy season, our school yard is \textbf{muddier than} the football pitch.} (comparative)
\item \emph{Rush hour is \textbf{the most polluted} moment of the day in Douala.} (superlative)
\item \emph{We all know \textbf{how dangerous} tramadol \textbf{is} for young apprentices.} (indirect question)
\item \emph{In March, it is \textbf{too hot to study} in the classroom at noon.} (too...to)
\item \emph{My classmates are \textbf{motivated enough to plant} trees around the school.} (enough...to)
\end{enumerate}
\end{corrige}

\begin{aretenir}[Key points of the section]
\begin{itemize}
\item Comparatives: \emph{hotter than, more dangerous than}; superlatives: \emph{the hottest, the most polluted}.
\item \emph{How + adjective}: exclamation (\emph{How hot it is!}) or indirect question without inversion (\emph{I know how dangerous it is}).
\item \emph{Too + adjective + to + infinitive} = excess, negative result $\Longleftrightarrow$ \emph{so + adjective + that + cannot}.
\item \emph{Adjective + enough + to + infinitive} = sufficiency, positive result; but \emph{enough + noun}.
\item \emph{So + adjective + that + clause} expresses a consequence.
\end{itemize}
\end{aretenir}

These structures will be your main tools in Section 5, where you will write a full composition about an environmental issue and carry out the integration project: describing a real problem, evaluating its degree of seriousness, and proposing realistic actions.

\coursec{Writing \& Integration : From Analysis to Action (Composition \& Project)}

In the previous section, we mastered advanced adjective structures to express degree and comparison precisely. Now, we bring together \textbf{all} the grammar, vocabulary, and functions from this module — Relative Clauses, Reported Speech, Conditionals, and Adjective Structures — to produce coherent, high-level written compositions and carry out an integrated project. This is the ultimate test of your communicative competence for the Probatoire and Baccalauréat Technique.

\courssub{Analysing Examination Prompts (Sujets Types Bac/Probatoire)}

Understanding the \textbf{task}, \textbf{format}, \textbf{audience}, and \textbf{register} is the first step to success. Technical exam questions typically fall into three main formats within this module.

\begin{exemplebox}[Common Exam Formats for Module 3]
\textbf{1. Formal Letter (Administrative / Civic Action)}
\begin{itemize}
    \item \textit{Prompt:} ``Write a letter to the Mayor of your council about the poor waste management in your neighbourhood. Propose concrete solutions.''
    \item \textit{Format:} Sender's address, Date, Receiver's address, Subject line, Formal salutation (Dear Sir/Madam), Body (3-4 paragraphs), Formal closing (Yours faithfully), Signature.
    \item \textit{Register:} Formal, polite but firm, persuasive.
\end{itemize}

\textbf{2. Magazine / Newspaper Article (Awareness / Opinion)}
\begin{itemize}
    \item \textit{Prompt:} ``Write an article for your school magazine on ``The Dangers of Drug Abuse among Youths''.''
    \item \textit{Format:} Catchy Title, By-line (By [Name], Form/Class), Introduction (Hook), Body (Arguments + Examples), Conclusion (Call to action).
    \item \textit{Register:} Semi-formal, engaging, informative, slightly persuasive.
\end{itemize}

\textbf{3. Speech / Talk (Ceremonial / Advocacy)}
\begin{itemize}
    \item \textit{Prompt:} ``As the President of the Environmental Club, deliver a speech on World Environment Day (June 5th) on the theme: ``Beat Plastic Pollution''.''
    \item \textit{Format:} Protocol (Recognising authorities), Hook (Quote/Statistic), Thesis, Body (PEEL paragraphs), Conclusion (Rallying cry/Call to action), Thanks.
    \item \textit{Register:} Oral style (rhetorical questions, direct address ``Dear friends'', repetition, inclusive ``we/our''), inspiring, authoritative.
\end{itemize}
\end{exemplebox}

\begin{attention}[Decoding the Prompt - Pièges Fréquents]
\begin{itemize}
    \item \textbf{Ignoring the Role/Persona:} If the prompt says ``As a Health Prefect...'', you \textbf{must} write in the first person (``I'') and use the authority of that role.
    \item \textbf{Missing the Target Audience:} Writing a letter to the Mayor requires formal administrative language (``I wish to draw your attention to...''\ ), not slang or over-familiarity.
    \item \textbf{Word Count:} Respect the limit (usually \textbf{250--300 words} for Bac Tech). Writing 400 words wastes time and risks penalisation for lack of concision; writing 150 words lacks development.
    \item \textbf{Off-topic:} ``Waste management'' $\neq$ ``Deforestation''. Stay focused on the specific issue (drugs, diseases, waste, climate).
\end{itemize}
\end{attention}

\courssub{The Architecture of a High-Scoring Composition}

A well-structured composition follows a logical funnel: General $\rightarrow$ Specific $\rightarrow$ General.

\begin{popfigure}
\begin{tikzpicture}[scale=0.7, every node/.style={font=\small}]
% Nodes
\node[draw, rectangle, rounded corners, fill=popblueL, align=center, text width=3.5cm] (intro) at (0,3.5) {\textbf{INTRODUCTION}\\($\approx$ 40-50 words)\\\textbf{1. Hook:} Stat / Quote / Rhetorical Q\\\textbf{2. Context:} Cameroon reality\\\textbf{3. Thesis Statement:} Main argument + 3 key points};
\node[draw, rectangle, rounded corners, fill=popgreenL, align=center, text width=3.5cm] (body1) at (-5,0) {\textbf{BODY 1: CAUSES / PROBLEM}\\($\approx$ 70 words)\\\textbf{P}oint: Topic Sentence\\\textbf{E}vidence: Local Data / Fact\\\textbf{E}xplanation: Analysis\\\textbf{L}ink: $\rightarrow$ Effects};
\node[draw, rectangle, rounded corners, fill=popgreenL, align=center, text width=3.5cm] (body2) at (0,0) {\textbf{BODY 2: CONSEQUENCES / IMPACT}\\($\approx$ 70 words)\\\textbf{P}oint: Topic Sentence\\\textbf{E}vidence: Case Study (Logone flood, School dropout)\\\textbf{E}xplanation: Why it matters\\\textbf{L}ink: $\rightarrow$ Solutions};
\node[draw, rectangle, rounded corners, fill=popgreenL, align=center, text width=3.5cm] (body3) at (5,0) {\textbf{BODY 3: SOLUTIONS / RECOMMENDATIONS}\\($\approx$ 70 words)\\\textbf{P}oint: Measure proposed\\\textbf{E}vidence: Example (Green Sahel, Eco-Collect)\\\textbf{E}xplanation: Feasibility\\\textbf{L}ink: $\rightarrow$ Conclusion};
\node[draw, rectangle, rounded corners, fill=poppinkL, align=center, text width=3.5cm] (conc) at (0,-3.5) {\textbf{CONCLUSION}\\($\approx$ 40 words)\\\textbf{1. Restate Thesis} (other words)\\\textbf{2. Call to Action} (Govt / Youth / You)\\\textbf{3. Final Thought} (Proverb / Vision)};
% Arrows
\draw[->, thick, popblue] (intro.south) -- ++(0,-0.3) -| (body1.north);
\draw[->, thick, popblue] (intro.south) -- ++(0,-0.3) -| (body2.north);
\draw[->, thick, popblue] (intro.south) -- ++(0,-0.3) -| (body3.north);
\draw[->, thick, popgreen] (body1.south) -- ++(0,-0.3) -| (conc.north);
\draw[->, thick, popgreen] (body2.south) -- ++(0,-0.3) -| (conc.north);
\draw[->, thick, popgreen] (body3.south) -- ++(0,-0.3) -| (conc.north);
\end{tikzpicture}
\legende{Structure type d'une composition d'examen (250--300 mots). La flèche centrale représente la Thesis Statement qui irrigue les trois paragraphes de développement.}
\end{popfigure}

\begin{definition}[Thesis Statement (Phrase Thèse)]
Une \cle{Thesis Statement} résume l'argument principal de la composition en une seule phrase, généralement placée à la fin de l'introduction. Elle annonce le plan (souvent 3 points).
\begin{itemize}
    \item \textbf{Weak:} ``There are many problems with waste in Cameroon.'' (Trop vague, pas de plan).
    \item \textbf{Strong:} ``\textbf{Effective waste management in Cameroon requires strict law enforcement, community recycling initiatives, and environmental education in schools.}'' (Clair, argumenté, plan en 3 points = 3 paragraphes Body).
\end{itemize}
\end{definition}

\begin{methode}[La Méthode PEEL pour les paragraphes de développement]
Chaque paragraphe \textbf{Body} doit suivre la structure \textbf{PEEL} :
\begin{enumerate}
    \item \textbf{P -- Point (Phrase Sujet) :} Une phrase claire qui annonce l'idée unique du paragraphe. \textit{Ex: "Firstly, weak enforcement of existing laws worsens plastic pollution."}
    \item \textbf{E -- Evidence (Preuve / Exemple) :} Donnée chiffrée, fait local précis, citation, exemple camerounais. \textit{Ex: "In Yaoundé, only 30\% of waste is collected regularly (HYSACAM 2022 report)."}
    \item \textbf{E -- Explanation (Analyse) :} Explique \textbf{POURQUOI} et \textbf{COMMENT} la preuve soutient le point. \textit{Ex: "This low rate forces citizens to burn or dump trash in streams, blocking drainage systems."}
    \item \textbf{L -- Link (Transition) :} Relie au paragraphe suivant ou revient à la thèse. \textit{Ex: "Consequently, these blocked drains lead to devastating floods during the rainy season."} (Transition vers Body 2 : Effects).
\end{enumerate}
\end{methode}

\courssub{Cohesion \& Coherence : Logical Connectors (Connecteurs Logiques)}

Examiners reward \textbf{cohesion} (linking words) and \textbf{coherence} (logical flow). Avoid basic connectors (``and, but, because, also''). Upgrade to B2/C1 level.

\begin{propriete}[Connecteurs Essentiels pour l'Argumentation (Niveau Bac Tech)]
\begin{center}
\begin{tabularx}{\linewidth}{|l|X|X|}
\hline
\rowcolor{popblueL}
\textbf{Fonction} & \textbf{Connecteurs (Mots de liaison)} & \textbf{Exemple Contextualisé (Module 3)} \\
\hline
\textbf{Addition / Renforcement} & \cle{Furthermore}, \cle{Moreover}, \cle{In addition}, \cle{What is more}, \cle{Besides} & ``\textbf{Moreover}, recycling creates green jobs for unemployed youths.''\\
\hline
\textbf{Contraste / Concession} & \cle{However}, \cle{Nevertheless}, \cle{Although} / \cle{Even though} + Clause, \cle{Despite} / \cle{In spite of} + Noun/-ing, \cle{On the one hand... On the other hand} & ``\textbf{Although} the 2014 plastic ban exists, \textbf{nevertheless}, thin bags flood markets.''\\
\hline
\textbf{Cause / Effet} & \cle{Consequently}, \cle{Therefore}, \cle{As a result}, \cle{Thus}, \cle{Hence}, \cle{Due to} / \cle{Owing to} + Noun & ``Drains are blocked; \textbf{consequently}, floods destroy homes in Douala.''\\
\hline
\textbf{Condition / Hypothèse} & \cle{Provided that}, \cle{As long as}, \cle{Unless}, \cle{If} (Cond 1/2) & ``\textbf{Provided that} funds are released, the recycling plant will open.''\\
\hline
\textbf{Exemplification} & \cle{For instance}, \cle{For example}, \cle{Such as}, \cle{Namely}, \cle{A case in point is} & ``\textit{A case in point is} the ``Eco-Collect'' initiative in Buea.''\\
\hline
\textbf{Conclusion} & \cle{To conclude}, \cle{In conclusion}, \cle{Ultimately}, \cle{All things considered} & ``\textbf{Ultimately}, a clean Cameroon is everyone's responsibility.''\\
\hline
\end{tabularx}
\end{center}
\end{propriete}

\begin{attention}[Erreurs de Cohésion -- Pièges à Éviter]
\begin{itemize}
    \item \textbf{Le "And/But/So" Syndrome :} Ne commencez \textbf{jamais} une phrase par ``And'', ``But'', ``So'' dans une composition formelle. Utilisez ``Furthermore'', ``However'', ``Therefore''.
    \item \textbf{Mauvaise ponctuation des connecteurs :} ``However'' (adverbe conjonctif) $\rightarrow$ Point-virgule avant, virgule après. ``Although'' (conjonction) $\rightarrow$ Pas de virgule avant si en milieu de phrase.
    \item \textit{Correct:} ``The law exists; \textbf{however}, it is ignored.'' / ``\textbf{Although} the law exists, it is ignored.''
    \item \textbf{Répétition lexicale :} Variez : \cle{waste} / \cle{garbage} / \cle{trash} / \cle{rubbish} / \cle{refuse} / \cle{litter} (déchets, ordures, détritus). \cle{Solutions} / \cle{measures} / \cle{strategies} / \cle{remedies} / \cle{actions}.
\end{itemize}
\end{attention}

\courssub{Grammatical Integration Checklist (La Check-list de Révision)}

C'est ici que vous gagnez les points de \textbf{Grammaire} et \textbf{Vocabulaire}. Avant de rendre votre copie, cochez ces 5 cases. Si une case est vide, \textbf{réécrivez la phrase concernée}.

\begin{definition}[Check-list d'Intégration Grammaticale Obligatoire (Cible : 5/5)]
\begin{enumerate}
    \item \textbf{\cle{Relative Clauses} (Min. 2) :} Une définissante (sans virgule) + Une non-définissante (avec virgules).
    \item \textbf{\cle{Reported Speech} (Min. 1) :} Rapport d'expert, de témoin, ou de loi. \textit{Ex: "The Minister stated that stricter measures would be taken."}
    \item \textbf{\cle{First Conditional} (Real Future) (Min. 1) :} \textit{Ex: "If the government enforces the law, pollution will decrease."}
    \item \textbf{\cle{Second Conditional} (Hypothetical Present) (Min. 1) :} \textit{Ex: "If I were the Mayor, I would launch a composting programme."}
    \item \textbf{\cle{Advanced Adjective Structures} (Min. 2 structures différentes) :}
    \begin{itemize}
        \item \cle{How + Adjective} (Exclamatif) : ``\textbf{How alarming} the rate of deforestation is!''
        \item \cle{Too ... to + Verb} (Excès négatif) : ``The waste is \textbf{too toxic to handle} without protection.''
        \item \cle{So ... that / Such ... that} (Conséquence) : ``It was \textbf{such a severe drought that} crops failed.''
        \item \cle{Superlatif + Ever / Of all} : ``The \textbf{worst flood ever} recorded in the Far North.''
    \end{itemize}
\end{enumerate}
\end{definition}

\begin{exoresolu}[Intégration Grammaticale dans un Paragraphe]
\textbf{Sujet :} ``Write a body paragraph on solutions to plastic pollution.''\\
\textbf{Production élève (Brut) :}
``We must recycle. The government must ban plastic. People must clean. If we do this, Cameroon will be clean.''

\textbf{Analyse (Check-list) :} 0 Relative, 0 Reported, 0 Cond 1, 0 Cond 2, 0 Adv Adj. \textbf{Note : 4/20 (Grammaire)}.

\textbf{Réécriture Intégrée (Cible 20/20 Grammaire) :}
\begin{quote}
\textbf{Furthermore}, \textbf{community recycling initiatives}, \textbf{which empower local women}, offer a sustainable alternative. \textit{(Relative non-définissante \checkmark)} The Director of HYSACAM \textbf{recently announced that} \textbf{new sorting centres would open} in 2025. \textit{(Reported Speech \checkmark)} \textbf{If these centres are accessible}, citizens \textbf{will sort} their waste at source. \textit{(Cond 1 \checkmark)} \textbf{If every household composted} organic waste, landfill volume \textbf{would drop} by 50\%. \textit{(Cond 2 \checkmark)} \textbf{How simple} this habit seems, \textbf{yet how powerful} its impact!\textit{(How + Adj \checkmark)} We \textbf{must not be too passive to act}; the future depends on us. \textit{(Too...to \checkmark)}
\end{quote}
\end{exoresolu}

\courssub{Integrated Project : ``Clean \& Healthy School Campaign'' (Projet Intégré)}

This project synthesises \textbf{Speaking, Listening, Reading, Writing, Grammar, Vocabulary}. Work in groups of 4-5. Duration: 3-4 sessions.

\begin{methode}[Étapes du Projet ``Clean \& Healthy School Campaign'']
\begin{enumerate}
    \item \textbf{Planning (Session 1) :} Assign roles: \textit{Project Manager, Graphic Designer, Speech Writer, Journalist, Presenter}. Choose a specific focus: ``Zero Plastic'', ``No to Drugs'', ``Stop Cholera''.
    \item \textbf{Production (Session 2-3) :} Create the 4 deliverables below. \textbf{Every member contributes written text.}
    \item \textbf{Presentation \& Evaluation (Session 4) :} Gallery Walk / Oral Defense. Peer assessment + Teacher Grid.
\end{enumerate}
\end{methode}

\begin{exemplebox}[Les 4 Livrables (Deliverables) -- Contraintes Grammaticales]
\begin{enumerate}
    \item \textbf{Affiche de Sensibilisation (Poster) -- \textit{Focus: Vocabulaire + Impératifs + Slogans}}
    \begin{itemize}
        \item Visuel fort (Dessiné ou Canva/PPT).
        \item 5 Impératifs variés : ``Sort your waste!'', ``Don't litter!'', ``Report leaks!'', ``Use the bin!'', ``Join the club!''.
        \item Slogan avec \textbf{How + Adj} ou \textbf{Too...to} : ``\textbf{How clean} our school can be!'' / ``\textbf{Too precious to pollute}.''
        \item Vocabulaire module : \cle{recycling bin} (poubelle de tri), \cle{biodegradable} (biodégradable), \cle{landfill} (décharge), \cle{contamination} (contamination), \cle{hygiene} (hygiène).
    \end{itemize}

    \item \textbf{Discours du Délégué / Président de Club (Speech) -- \textit{Focus: Reported Speech + Rhetorical Devices}}
    \begin{itemize}
        \item 2-3 minutes oral.
        \item Inclure \textbf{au moins 2 Reported Speech} : ``\textbf{The Principal emphasized that} health is our wealth.'' / ``\textbf{WHO warns that} tobacco kills 8 million yearly.''
        \item Adresse directe : ``Dear friends, \textbf{imagine a school where}...'' (Cond 2).
        \item Call to action final.
    \end{itemize}

    \item \textbf{Article pour le Journal Scolaire (Magazine Article) -- \textit{Focus: Conditionnels + Structures Adjectivales + PEEL}}
    \begin{itemize}
        \item Titre accrocheur. 250 mots.
        \item Structure PEEL x 3.
        \item \textbf{1 Cond 1}, \textbf{1 Cond 2}, \textbf{1 How+Adj}, \textbf{1 Too...to}, \textbf{2 Relatives}.
        \item Exemple titre : ``\textbf{How Green} Can We Go? The \textbf{Too Urgent to Ignore} Truth About Our Playground.''
    \end{itemize}

    \item \textbf{Présentation Orale de Groupe (Oral Defense) -- \textit{Focus: Fluency + Interaction + Grammar Spotlight}}
    \begin{itemize}
        \item 5 min présentation + 3 min Q\&A.
        \item Chaque membre parle.
        \item Le jury (classe/prof) note l'utilisation spontanée des structures cibles.
    \end{itemize}
\end{enumerate}
\end{exemplebox}

\courssub{Self-Assessment \& Remediation (Auto-évaluation \& Remédiation)}

Devenir autonome, c'est savoir \textbf{où} l'on perd des points pour les corriger \textbf{avant} l'examen.

\begin{definition}[Grille d'Auto-évaluation Critériée (Sur 20 -- Standard Bac)]
Utilisez cette grille pour noter votre propre composition (ou celle d'un pair) \textbf{avant} la correction prof.

\begin{center}
\begin{tabularx}{\linewidth}{|l|X|c|}
\hline
\rowcolor{popblueL}
\textbf{Critère} & \textbf{Indicateurs de Réussite (Niveau "Bien" / 4/4)} & \textbf{/4} \\
\hline
\textbf{1. Contenu / Task Achievement} & Répond exactement au sujet. Idées pertinentes, développées, exemples locaux (Cameroun). Ton/Registre adaptés. & \\
\hline
\textbf{2. Organisation / Cohérence} & Structure Intro/Body/Conclusion respectée. Paragraphes PEEL. Connecteurs variés (Add/Contrast/Cause/Cond). Thesis Statement claire. & \\
\hline
\textbf{3. Grammaire / Structures} & \textbf{Check-list 5/5 validée} (2 Relatives, 1 Reported, 1 Cond1, 1 Cond2, 2 Adv Adj). Phrases complexes, peu d'erreurs bloquantes. & \\
\hline
\textbf{4. Vocabulaire / Lexique} & Vocabulaire spécifique module (mitigation, resilience, withdrawal, contamination, landfill, enforcement...). Pas de répétition. Collocations justes. & \\
\hline
\textbf{5. Orthographe / Ponctuation / Lisibilité} & Mots usuels corrects. Majuscules/ponctuation anglaises respectées. Écriture lisible. Respect nombre de mots ($\pm$ 10\%). & \\
\hline
\rowcolor{popblueL}
\textbf{TOTAL} &  & \textbf{/20} \\
\hline
\end{tabularx}
\end{center}
\end{definition}

\begin{methode}[Procédure de Remédiation Ciblée (Post-Auto-éval)]
\begin{enumerate}
    \item \textbf{Identifiez la note la plus basse} (ex: Grammaire 1/4).
    \item \textbf{Ciblez la lacune précise :} ``Je n'ai pas mis de Conditionnel 2'' ou ``Je confonds Although / However''.
    \item \textbf{Action corrective immédiate (15 min) :}
    \begin{itemize}
        \item \textit{Grammaire :} Refaire 3 exercices de transformation (Active $\rightarrow$ Passive, Direct $\rightarrow$ Reported, Cond 1 $\leftrightarrow$ Cond 2) sur le manuel / OnBuch+.
        \item \textit{Vocabulaire :} Créer 5 phrases personnelles avec les 5 mots manquants (Flashcards).
        \item \textit{Structure :} Réécrire \textbf{un seul paragraphe} faible en appliquant PEEL + Check-list grammaire.
    \end{itemize}
    \item \textbf{Re-vérification :} Re-noter ce paragraphe réécrit avec la grille.
\end{enumerate}
\end{methode}

\begin{attention}[Remédiation : Erreurs Systémiques Fréquentes en Terminale Tech]
\begin{itemize}
    \item \textbf{Subject-Verb Agreement :} ``The effects of climate change \textbf{is} devastating'' $\rightarrow$ \textbf{are}. (Le sujet est ``effects'', pluriel).
    \item \textbf{Article Usage :} ``\textbf{The} pollution is bad'' (Général $\rightarrow$ Zéro article) vs ``\textbf{The} pollution \textbf{in our city} is bad'' (Spécifique $\rightarrow$ The).
    \item \textbf{Prepositions après Adjectifs/Verbes :} \textit{Interested \textbf{in}}, \textit{Good \textbf{at}}, \textit{Responsible \textbf{for}}, \textit{Aware \textbf{of}}, \textit{Depend \textbf{on}}, \textit{Consist \textbf{of}}.
    \item \textbf{False Friends :} ``Actually'' = \textit{En fait} (pas \textit{Actuellement} = Currently). ``Eventually'' = \textit{Finalement} (pas \textit{Éventuellement} = Possibly).
\end{itemize}
\end{attention}

\courssub{Exam Simulation : Timed Writing Practice (Entraînement Chronométré)}

\begin{exoresolu}[Sujet Type Bac Tech 2023 -- Simulation 45 min]
\textbf{Prompt :} ``Write an article for your school magazine on \textbf{``The Role of Youth in Fighting Climate Change in Cameroon''}.'' \textbf{Length:} 250--300 words. \textbf{Constraints:} Integrate 1 Relative, 1 Cond 1, 1 ``Too...to'', Specific Vocab (mitigation, adaptation, resilience).

\textbf{Plan Rapide (5 min) :}
\begin{itemize}
    \item \textbf{Thesis:} Youth are the driving force of climate action through education, innovation, and advocacy.
    \item \textbf{Body 1 (Education/Awareness):} School clubs, Green Sahel. \textit{Vocab: mitigation, awareness.}
    \item \textbf{Body 2 (Innovation/Action):} Recycling startups, tree planting. \textit{Vocab: adaptation, resilience.}
    \item \textbf{Body 3 (Advocacy/Policy):} Voice at COP, pressure leaders. \textit{Grammar targets.}
\end{itemize}

\textbf{Modèle de Réponse (Corrigé Intégré) :}
\begin{quote}
\textbf{The Role of Youth in Fighting Climate Change in Cameroon}

\textbf{How crucial} is the energy of youth in the battle for our planet's survival? In Cameroon, where floods in the Far North and landslides in the West signal a climate emergency, young people \textbf{who understand the stakes} are no longer waiting for adults to act. They are the engine of \textbf{mitigation} (atténuation) and \textbf{adaptation} (adaptation).

\textbf{Firstly}, education builds \textbf{resilience} (résilience). School environmental clubs, \textbf{which multiply across the country}, turn theory into practice. \textbf{If students learn} composting and tree planting today, they \textbf{will protect} their farms tomorrow. \textit{(Cond 1 \checkmark)} For instance, the ``Green Sahel'' initiative sees thousands of youths planting acacia trees to stop desertification.

\textbf{Moreover}, innovation drives the circular economy. Young entrepreneurs in Douala and Yaoundé transform plastic waste into paving stones and fuel. \textbf{If I were} the Minister of Environment, I \textbf{would subsidize} these green startups massively. \textit{(Cond 2 \checkmark)} \textbf{The Minister of Youth recently declared that} government grants for green projects \textbf{would increase} next year. \textit{(Reported Speech \checkmark)}

\textbf{However}, advocacy remains \textbf{too vital to ignore}. Youths \textbf{must not be too silent to influence} policy. \textit{(Too...to \checkmark)} \textbf{Provided that} they organise via social media and vote, they can force leaders to honour the 35\% emission reduction pledge (NDC).

\textbf{In conclusion}, Cameroon's climate future rests on its youth. \textbf{Let us not be spectators of our own destiny.} Plant a tree, recycle a bottle, raise your voice. The time is now.
\end{quote}

\textbf{Analyse du Corrigé :} $\approx$ 260 mots. Thesis claire. 3 Body PEEL. Check-list Grammaire : Relative (\textit{who understand}, \textit{which multiply}) \checkmark, Reported (\textit{declared that... would increase}) \checkmark, Cond 1 \checkmark, Cond 2 \checkmark, Too...to \checkmark, How+Adj \checkmark. Vocab : \textit{mitigation, adaptation, resilience, desertification, circular economy, pledge, NDC}. Connecteurs : \textit{Firstly, Moreover, However, In conclusion, Provided that}.
\end{exoresolu}

\begin{savaistu}[Le Cameroun et l'Action Climatique : Le Rôle des Jeunes]
Le Cameroun a soumis sa \textbf{CDN (Contribution Déterminée au niveau National)} à l'ONU, s'engageant à réduire ses émissions de gaz à effet de serre de \textbf{35\% d'ici 2030} (scénario conditionnel : avec appui international).
\begin{itemize}
    \item \textbf{Grande Muraille Verte :} Initiative panafricaine (Sahara au Sahel). Au Nord-Cameroun, des milliers de jeunes (Lycéens, étudiants, associations) participent aux campagnes de reboisement (Acacia, Neem) pour fixer les dunes et restaurer les sols.
    \item \textbf{Innovation Locale :} Des startups comme ``Kemit Ecology'' (Douala) transforment les déchets ménagers en charbon écologique (briquettes), luttant contre la déforestation (bois de chauffe) et l'insalubrité.
    \item \textbf{Données clés (Source: MINEDUB / HYSACAM / OMS) :} Production annuelle de déchets solides $\approx$ 6 millions de tonnes ; taux de collecte formelle $\approx$ 30-40\% dans les grandes villes (Yaoundé, Douala) ; le plastique représente $\sim$ 10-15\% du flux.
    \item \textbf{Votre rôle :} En tant qu'élèves de Terminale Technique, vous avez des compétences pratiques (électricité, mécanique, génie civil, commerce...). \textbf{La transition écologique a besoin de techniciens :} installateurs solaires, gestionnaires de déchets, agriculteurs climato-intelligents. Votre Bac Tech est un passeport pour l'économie verte !
\end{itemize}
\end{savaistu}

\begin{exercice}[Entraînement Final -- Épreuve Complète 45 min]
Choisissez \textbf{UN} sujet. Respectez le format, le nombre de mots (250-300) et la \textbf{Check-list Grammaticale (5/5)}.

\textbf{Sujet A (Lettre Formelle) :} Vous êtes le Président du Club ``Health \& Environment'' de votre établissement. Écrivez une lettre au \textbf{Ministre de la Santé Publique} pour dénoncer la vente de drogues (Tramadol, cannabis, alcool frelaté) aux abords des écoles. Proposez 3 mesures concrètes.

\textbf{Sujet B (Article Magazine) :} ``\textbf{Turning Trash into Treasure: The Power of Recycling in Cameroon.}'' Article pour le journal de l'école. Montrez comment le recyclage crée de l'emploi, protège l'environnement et génère des revenus.

\textbf{Sujet C (Discours) :} ``En tant que Délégué de la classe de Terminale, prononcez un discours lors de la \textbf{Journée Mondiale de l'Environnement (5 Juin)} sur le thème : ``Notre École, Notre Planète : Zéro Plastique, Zéro Drogue, Zéro Déchet.''''

\textbf{Consigne :} Écrivez votre production sur une feuille d'examen. Chronométrez-vous (45 min). Puis auto-évaluez-vous avec la \textbf{Grille (5.6)}.
\end{exercice}

\begin{corrige}[Exercice 5.7 -- Pistes de Correction (Sujet B - Article)]
\textbf{Structure Attendue :}
\begin{itemize}
    \item \textbf{Titre:} Turning Trash into Treasure: The Power of Recycling in Cameroon.
    \item \textbf{Intro:} Hook (Stat: 6M tons waste/yr in CMR). Context: Urban pollution. Thesis: Recycling = Economic + Ecological + Social win.
    \item \textbf{Body 1 (Economic/Jobs):} Point: Job creation. Evidence: Eco-Collect, Kemit Ecology. Explanation: Value chain (collection $\rightarrow$ sorting $\rightarrow$ processing). Link: But environment gains too.
    \item \textbf{Body 2 (Environment):} Point: Less landfill/emissions. Evidence: Methane reduction, saved forests (charcoal alternative). Explanation: Circular economy. Grammar: \textit{If waste is sorted, methane will drop (Cond 1)}.
    \item \textbf{Body 3 (Social/Community):} Point: Community pride/health. Evidence: Clean neighbourhoods = less cholera/malaria. Grammar: \textit{The Mayor stated that bins would be provided (Reported)}. \textit{How clean our streets become! (How+Adj)}.
    \item \textbf{Conclusion:} Restate. Call to Action: Sort at source. Final: \textit{Waste is too valuable to waste (Too...to)}.
\end{itemize}
\textbf{Vérif Check-list :} 2 Relatives? 1 Reported? 1 Cond 1? 1 Cond 2? 2 Adv Adj? Vocab: \textit{Value chain, landfill, methane, circular economy, sorting, processing}.
\end{corrige}

\newpage

\coursec{Activité d'intégration}

\coursec{Environment, Well-being and Health : Integration Project — « Jeunesse, Assainissement et Prévention »}

\courssub{Objectives of the Integration Activity}
\begin{itemize}
    \item Mobilise all linguistic skills of the unit: vocabulary (health, environment, addictions), grammar (relative clauses, reported speech, conditionals 1 \& 2, advanced adjective structures: \cle{how + adjective}, \cle{too ... to}, superlatives).
    \item Develop the targeted savoir-faire: \emph{apply concepts to real-life situations} and \emph{write and justify a process or conclusion}.
    \item Produce a complete \textbf{multimedia advocacy campaign}: field survey (oral), technical policy brief (written), visual poster, and oral pitch.
    \item Assess transversal skills: teamwork, critical thinking, active citizenship, digital communication.
\end{itemize}

\courssub{Context \& Situation}
\textbf{Context:} The Arrondissement Commune of \textbf{Nkolbisson} (Yaoundé 6) faces a dual health and environmental crisis. The Departmental Delegate for Environment, Nature Protection and Sustainable Development (MINEPDED) and the Departmental Delegate for Public Health (MINSANTE) jointly launch the Call for Projects \og \emph{Youth, Sanitation and Prevention} \fg{} targeting technical high schools in the arrondissement.

\textbf{Simulated Real Data (based on typical local reports):}
\begin{center}
\begin{tabularx}{\linewidth}{|l|X|}\hline
\textbf{Indicator} & \textbf{Nkolbisson Data — Etoudi / Melen Quarter} \\ \hline
Illegal dumpsites & 3 major unauthorized sites (one 50m from Nkolbisson Technical High School). Estimated volume: \num{15} tons/month. \\ \hline
Waste collection & HYSACAM coverage rate: 40\%. Frequency: 1 pass/week (insufficient). \\ \hline
Malaria (declared cases at CSPS) & \num{320} cases/month (45\% children < 15 years). Peak in rainy season (stagnant water in waste). \\ \hline
Cholera / Typhoid & \num{12} confirmed cases last quarter (link: polluted water / waste). \\ \hline
Tramadol / \emph{Tramol} consumption & Anonymous school survey (2023): 18\% of Terminale students have used. Sales points: 3 \emph{call-boxes} / kiosks < 200m from school. \\ \hline
Illicit alcohol (\emph{Odontol}, \emph{Arki}) & Easy availability. 2 unlicensed bars open from 6am near student entrance. \\ \hline
Municipal sanitation budget & \num{5.000.000} FCFA/year (insufficient for bins + sorting center). \\ \hline
\end{tabularx}
\end{center}

\textbf{Mission:} As Terminale Technique students (Civil Engineering, Electrotechnics, Commerce, Secretariat), you form \textbf{project teams (4-5 students)}. You must submit a \textbf{complete dossier} before the deadline (end of term). A jury (Teachers, Student Delegates, MINEPDED/MINSANTE Reps) will select the best project for a pilot grant of \num{500.000} FCFA.

\courssub{Instructions \& Timeline}
\textbf{Duration: 3 sessions (2h + 2h + 2h) + personal work.}

\courssub{Step 1: Field Survey \& Reported Speech (Session 1 - 2h)}
\begin{enumerate}
    \item \textbf{Preparation (Pairs):} Design an interview grid (10 questions) for two targets:
    \begin{itemize}
        \item Target A: School nurse (or Etoudi CSPS health worker) — focus: diseases linked to insanitation, observed addictions.
        \item Target B: Local resident / parent / shopkeeper — focus: daily experience (smells, mosquitoes, insecurity, drug sales).
    \end{itemize}
    \item \textbf{Execution (Out of class / Dedicated slot):} Conduct interviews (10 min each). Record (phone) or take detailed notes. \textbf{Key vocabulary to use:} \emph{breeding ground, stagnant water, waste disposal, drug abuse, peer pressure, withdrawal symptoms, rehabilitation, awareness campaign}.
    \item \textbf{Written Production (Class):} Write a \textbf{survey report (150 words)} in \textbf{Reported Speech}.
    \begin{itemize}
        \item Example: \og The nurse \textbf{stated that} malaria cases \textbf{had increased} by 20\% \textbf{due to} stagnant water in illegal dumps. \fg{}
        \item Example: \og A resident \textbf{reported that} drug dealers \textbf{sold} tramadol openly near the school gate. \fg{}
    \end{itemize}
    \item \textbf{Deliverable:} Audio file/notes + Written report (Submitted to teacher).
\end{enumerate}

\courssub{Step 2: Data Analysis \& Technical Diagnosis (Session 2 - 1h)}
\begin{enumerate}
    \setcounter{enumi}{4}
    \item \textbf{Exploit the data table (above):} Write the \textbf{\og Diagnosis \fg{}} section of the Policy Brief (200 words).
    \item \textbf{Mandatory Grammar Constraints:}
    \begin{itemize}
        \item Minimum \textbf{3 Relative Clauses} (defining \& non-defining). \emph{Ex: The illegal dumps \textbf{which are located near the school} attract mosquitoes.}
        \item Minimum \textbf{2 \cle{How + adjective} structures}. \emph{Ex: \textbf{How critical} the situation is! / \textbf{How urgent} action is needed!}
        \item Precise vocabulary: \emph{tonnage, prevalence rate, coverage ratio, hotspot, unlicensed premises}.
    \end{itemize}
\end{enumerate}

\courssub{Step 3: Policy Brief — Conditional Scenarios \& Argumentation (Session 2 - 1h + Homework)}
\begin{enumerate}
    \setcounter{enumi}{5}
    \item Write the complete briefing note (2 pages max, font 12, 1.5 spacing) addressed to the Mayor of Nkolbisson. Imposed structure:
    \begin{enumerate}
        \item \textbf{Title \& Addressee.}
        \item \textbf{Executive Summary (5 lines).}
        \item \textbf{Diagnosis} (Step 2).
        \item \textbf{Key Testimonies} (Extract from Step 1 — Reported Speech).
        \item \textbf{Intervention Scenarios (Grammar: Conditionals):}
        \begin{itemize}
            \item \textbf{Scenario A (Short term — Conditional 1 — Real Future):} \og \textbf{If the Council provides} 10 additional bins and \textbf{organizes} weekly cleaning, \textbf{we will reduce} illegal dumping by 50\% within 3 months. \fg{}
            \item \textbf{Scenario B (Long term — Conditional 2 — Hypothetical):} \og \textbf{If the Council built} a local recycling center, \textbf{it would create} jobs for youth and \textbf{would stop} open burning. \fg{}
        \end{itemize}
        \item \textbf{Final Argument (Adjective Structures):} Mandatory use of:
        \begin{itemize}
            \item \cle{Too ... to}: \og The risk is \textbf{too high to ignore}. \fg{}
            \item \cle{Superlative}: \og This is \textbf{the most sustainable} solution for our community. \fg{}
            \item \cle{How + adj} (recall).
        \end{itemize}
        \item \textbf{Concrete Action Plan (3 costed measures, responsible parties, deadlines).}
    \end{enumerate}
\end{enumerate}

\courssub{Step 4: Visual Production \& Oral Pitch (Session 3 - 2h)}
\begin{enumerate}
    \setcounter{enumi}{6}
    \item \textbf{Awareness Poster (A3 format — Canva / Paper):} Catchy slogan + 3 key messages (Visuals + Short text). Themes: \og Selective Sorting: I Sort, I Win \fg{} / \og No to Tramadol, Yes to My Future \fg{} / \og My Clean Quarter, My Pride \fg{}. \textbf{Constraint:} Include at least one \cle{Too ... to} or \cle{How + adj} structure on the poster.
    \item \textbf{Oral Pitch (3 min / team):} Project presentation to jury.
    \begin{itemize}
        \item 1 min: Problem statement (Data + Reported testimonies).
        \item 1 min: Solution (Conditional Scenarios + Action Plan).
        \item 1 min: Call to action (Strong argument: Superlatives, \cle{How + adj}).
    \end{itemize}
    \textbf{Criteria:} Fluency, pronunciation, correct use of target structures, eye contact, time management.
\end{enumerate}

\courssub{Step 5: Assessment \& Remediation (Post-session)}
\begin{enumerate}
    \setcounter{enumi}{8}
    \item Assessment via criteria grid (out of 20): Technical Content/Data (5) — Language Quality/Target Grammar (5) — Feasibility/Realism (5) — Creativity/Supports (5).
    \item Individual remediation: Feedback sheet per student targeting recurring grammar errors (e.g., Cond 1/Cond 2 confusion, omitted relative pronoun, incorrect backshifting).
\end{enumerate}

\courssub{Resources to Mobilise}

\begin{definition}[Relative Clauses]
\textbf{Defining (Restrictive):} Identifies the antecedent. No commas. Pronouns: \cle{who} (people), \cle{which} (things), \cle{that} (people/things), \cle{whose} (possession), \cle{where} (place), \cle{when} (time). \emph{Ex: The bins \textbf{which/that are full} smell bad.}
\textbf{Non-defining (Non-restrictive):} Adds extra info. Commas mandatory. \cle{That} forbidden. \emph{Ex: The Nkolbisson Council, \textbf{which lacks funds}, cannot buy new trucks.}
\end{definition}

\begin{definition}[Reported Speech — Main Backshifting Rules]
\begin{center}
\begin{tabular}{|l|l|}\hline
\textbf{Direct Speech} & \textbf{Reported Speech} \\ \hline
Present Simple (\og I \textbf{live} \fg{}) & Past Simple (\og He \textbf{lived} \fg{}) \\ \hline
Present Continuous (\og I \textbf{am working} \fg{}) & Past Continuous (\og He \textbf{was working} \fg{}) \\ \hline
Present Perfect (\og I \textbf{have seen} \fg{}) & Past Perfect (\og He \textbf{had seen} \fg{}) \\ \hline
Past Simple (\og I \textbf{saw} \fg{}) & Past Perfect (\og He \textbf{had seen} \fg{}) \\ \hline
Will (\og I \textbf{will} go \fg{}) & Would (\og He \textbf{would} go \fg{}) \\ \hline
Can (\og I \textbf{can} swim \fg{}) & Could (\og He \textbf{could} swim \fg{}) \\ \hline
\textbf{Time/Place:} Now $\to$ Then, Today $\to$ That day, Here $\to$ There, Tomorrow $\to$ The next day. & \\ \hline
\end{tabular}
\end{center}
\textbf{Rule:} The reporting verb (\emph{said, told, reported, stated, explained}) is usually past $\to$ triggers backshifting.
\end{definition}

\begin{definition}[Conditionals: Type 1 vs Type 2]
\begin{itemize}
    \item \textbf{Type 1 (Real Future — Possible):} \cle{If + Present Simple, ... Will + Base Verb}. \emph{Use: Promises, warnings, real plans.} \og \textbf{If we clean} the gutters, \textbf{water will flow} normally. \fg{}
    \item \textbf{Type 2 (Hypothetical / Unreal Present):} \cle{If + Past Simple (were for \emph{be}), ... Would + Base Verb}. \emph{Use: Dreams, advice (\og If I were you \fg{}), unlikely hypotheses.} \og \textbf{If the Council built} a center, \textbf{it would solve} the problem. \fg{}
\end{itemize}
\end{definition}

\begin{definition}[Advanced Adjective Structures]
\begin{itemize}
    \item \textbf{How + Adjective (+ Subject + Verb):} Exclamation. \emph{Ex: \textbf{How dangerous} this waste is! / \textbf{How urgent} the need for action is!}
    \item \textbf{Too + Adjective + To + Verb:} Negative excess $\to$ Impossibility. \emph{Ex: The dump is \textbf{too dangerous to approach}.}
    \item \textbf{Adjective + Enough + To + Verb:} Sufficiency. \emph{Ex: The budget is \textbf{large enough to buy} bins.}
    \item \textbf{Superlative:} \cle{The most + adj (long)} / \cle{The -est (short)}. \emph{Ex: \textbf{The most effective} measure. \textbf{The cleanest} area.}
\end{itemize}
\end{definition}

\begin{definition}[Key Vocabulary — Memo]
\textbf{Health/Addictions:} \emph{Substance abuse, addiction, withdrawal, overdose, peer pressure, rehabilitation, counselling, narcotics, tramadol, cannabis, alcohol abuse, liver damage, mental health.}
\textbf{Environment/Waste:} \emph{Illegal dumping, landfill, recycling, composting, waste sorting/segregation, incineration, carbon footprint, climate change, global warming, greenhouse effect, stagnant water, breeding ground, vector-borne diseases (malaria, cholera), sanitation, hygiene, HYSACAM.}
\end{definition}

\courssub{Step-by-Step Solution \& Commentary}

\begin{exoresolu}{Model Production — Team \og Green Tech Nkolbisson \fg{} (Simulated Correction)}
\textbf{Step 1: Survey Report Extract (Reported Speech)}
\textit{Context:} Interview with Mr \textsc{Atangana}, Head Nurse CSPS Etoudi.
The nurse \textbf{stated that} they \textbf{received} at least 15 malaria cases daily during the rainy season. He \textbf{explained that} the stagnant water in the illegal dump behind the school \textbf{was} the main breeding ground. He \textbf{added that} three students \textbf{had been brought} in the previous month for tramadol overdose. He \textbf{reported that} they \textbf{bought} it at the kiosk near the gate.
\par \textbf{Justification:} \cle{Receive $\to$ Received} (Past Simple), \cle{Is $\to$ Was} (Present $\to$ Past), \cle{Buy $\to$ Had bought} (Past Simple $\to$ Past Perfect), \cle{Last month $\to$ The previous month}. \textbf{Note:} \emph{That} optional but recommended for formal style.
\tcblower
\textbf{Step 2: Diagnosis Extract (Relatives + How + Adj)}
The district of Nkolbisson faces a sanitation crisis \textbf{that demands} immediate intervention. \textbf{How critical} the situation is! Three major illegal dumps, \textbf{which are located less than 100 meters from residential areas and the Technical High School}, accumulate approximately 15 tons of waste monthly. The HYSACAM collection coverage, \textbf{whose current ratio stands at only 40\%}, is insufficient to prevent open burning. Consequently, the prevalence rate of malaria, \textbf{which is a vector-borne disease linked to stagnant water}, reaches 320 cases per month at the local health center. \textbf{How alarming} these figures are for a school environment!
\par \textbf{Analysis:} 4 Relatives (\textbf{that demands}, \textbf{which are located}, \textbf{whose current ratio}, \textbf{which is a vector-borne disease}). 2 \cle{How + Adj} (\textbf{How critical}, \textbf{How alarming}). Technical vocab: \emph{prevalence rate, vector-borne, coverage ratio, residential areas}.

\textbf{Step 3: Policy Brief — Conditional Scenarios \& Argument (Extract)}
\textbf{Scenario A (Immediate Action — Conditional 1):}
\textbf{If the Municipal Council provides} ten reinforced concrete bins at strategic hotspots \textbf{and organizes} a bi-weekly collection schedule with HYSACAM, \textbf{we will reduce} illegal dumping by 50\% within the first quarter. \textbf{This will allow} students to attend school in a healthier environment.

\textbf{Scenario B (Structural Solution — Conditional 2):}
\textbf{If the Council secured} funding from FEICOM or international partners (AFD, World Bank) \textbf{and built} a micro-recycling center managed by a youth cooperative, \textbf{it would create} decent jobs for unemployed graduates \textbf{and would transform} waste into resources (compost, plastic pellets). \textbf{This would address} both the environmental and the socio-economic roots of the crisis.

\textbf{Final Argument:}
The current situation is \textbf{too critical to ignore}. The health of our youth is \textbf{too precious to sacrifice} on the altar of administrative inertia. \textbf{How urgent} the need for political courage is! Investing in sanitation today is \textbf{the most profitable} investment for the municipality's future. \textbf{The cleanest} neighbourhoods attract businesses and tourism. We urge the Mayor to adopt Scenario A immediately and launch the feasibility study for Scenario B.
\par \textbf{Analysis:} Cond 1 (\textbf{provides/organizes $\to$ will reduce/will allow}): Concrete actions, existing budget. Cond 2 (\textbf{secured/built $\to$ would create/would transform/would address}): Hypothesis of external funding, heavy structure. Adj structures: \textbf{too critical to ignore}, \textbf{too precious to sacrifice}, \textbf{How urgent}, \textbf{the most profitable}, \textbf{The cleanest}. All constraints met.

\textbf{Step 4: Poster Design (Textual Description — Canva Simulation)}
\begin{center}
\begin{tikzpicture}
\node[draw=popdark, thick, rounded corners, inner sep=10pt, text width=0.9\linewidth, align=center, fill=popcream] (affiche) {
\textbf{\Large \textcolor{popgreen}{Nkolbisson PROPRE : MON QUARTIER, MA FIERTÉ}} \\[0.5em]
\textcolor{poporange}{\Large \textbf{TOO DIRTY TO LIVE IN? ACT NOW!}} \\[1em]
\begin{minipage}{0.3\linewidth}\centering
\textcolor{popblue}{\textbf{[Recycle Logo]}} \\ \small \textbf{TRIEZ} \\ Plastiques / Organiques
\end{minipage}
\begin{minipage}{0.3\linewidth}\centering
\textcolor{popred}{\textbf{[Pill + Cross]}} \\ \small \textbf{NON AU TRAMADOL} \\ \emph{Your brain is too valuable to lose.}
\end{minipage}
\begin{minipage}{0.3\linewidth}\centering
\textcolor{popgreen}{\textbf{[Tree Icon]}} \\ \small \textbf{1 ARBRE / ÉLÈVE} \\ Reboisement berges Mfoundi
\end{minipage} \\[1em]
\textit{Projet Terminale Tech - Lycée de Nkolbisson - Partenariat MINEPDED / MINSANTE / Commune}
};
\legende{Poster mock-up. Integration of \cle{Too ... to}: \og Your brain is too valuable to lose \fg{}. Icons replaced by text placeholders for compilation safety.}
\end{tikzpicture}
\end{center}

\textbf{Step 4: Oral Pitch Script (3 min — Extract 1 min 30)}
\og Good morning Honorable Jury, Dear Mayor's Representative. \textbf{How critical} is the situation in Etoudi today? Let me show you. \emph{[Shows malaria graph]}. 320 cases monthly. \textbf{If we do nothing}, this number \textbf{will double} next year (Cond 1 prediction). Our team, \textbf{who conducted} interviews with the school nurse and residents (Relative), found that drug dealers operate 50m from our gate. The nurse \textbf{reported that} tramadol was sold like candy (Reported Speech). 

Our proposal is simple. \textbf{Scenario A: If the Council provides} bins this month, \textbf{we will organize} student brigades for sorting (Cond 1). \textbf{Scenario B: If a recycling center were built} (Cond 2), it \textbf{would employ} 20 youths. 

The cost of inaction is \textbf{too high to pay}. Our health is \textbf{too precious to lose}. \textbf{How powerful} youth action can be! We are ready to work. Thank you. \fg{}
\par \textbf{Commentary:} Time respected. All target structures present. Persuasive tone. Precise vocabulary.
\end{exoresolu}

\courssub{Activity Review}

\begin{aretenir}[Validated Competencies]
\begin{itemize}
\item \textbf{Oral (Speaking/Listening):} Ability to conduct a survey (questionnaire design, interview, note-taking), report information in \textbf{Reported Speech} (mastered backshifting), present a structured project (\emph{Pitch}) with fluency and technical vocabulary.
\item \textbf{Written (Reading/Writing):} Exploitation of real data to produce a standard technical document (\emph{Policy Brief}): precise diagnosis (Relatives, \cle{How + adj}), prospective scenarios (Conditionals 1 \& 2 discriminated), persuasive argumentation (\cle{Too ... to}, Superlatives).
\item \textbf{Grammar in Context:} Structures are no longer decontextualized exercises but \textbf{tools for thought and action}: relatives for precise diagnosis, reported speech for giving voice to actors, conditionals for simulating the future, adjective structures for impact.
\item \textbf{Transversal Skills:} Teamwork (role distribution: investigator, writer, designer, speaker), citizenship (ownership of local problems, concrete proposals), digital skills (Canva, audio recording, word processing).
\end{itemize}
\end{aretenir}

\begin{attention}[Watch Points \& Typical Remediation]
\begin{itemize}
\item \textbf{Conditional 1 vs 2:} Frequent error \og \textbf{If the Council will provide} \fg{} (Incorrect). \emph{Remediation:} Rule \og \textbf{Never \textbf{will} after \textbf{if}} \fg{} (except rare volitional sense). Drill: \emph{If + Present $\to$ Future} / \emph{If + Past $\to$ Would}.
\item \textbf{Reported Speech:} Forgotten backshifting on modals (\emph{can $\to$ could, will $\to$ would}) and time/place markers (\emph{now, today, here}). \emph{Remediation:} Conversion table displayed in class during writing.
\item \textbf{Relative Clauses:} Use of \emph{that} in non-defining (\og The Council, \textbf{that} lacks funds... \fg{}). \emph{Remediation:} Comma test / removability. \emph{Which/Who} mandatory with commas.
\item \textbf{Too ... to:} Wrong positive meaning (\og He is \textbf{too strong to lift} the box \fg{} = He is too strong to lift? No, structure = negative excess). \emph{Remediation:} Systematic translation \og Trop ... pour [pouvoir] \fg{}.
\item \textbf{Vocabulary:} Confusion \emph{waste / waist}, \emph{dump / damp}, \emph{trial / trail}. \emph{Remediation:} Minimal pairs list for revision.
\end{itemize}
\end{attention}

\begin{savaistu}[Local Anchoring \& Follow-up]
This project aligns with the \textbf{National Development Strategy (SND30)} and the \textbf{Communal Development Plan (PCD)} of Nkolbisson. The best \emph{Policy Briefs} may be forwarded to the Mayor and the MINEPDED Delegate for real follow-up (FEICOM grant application, partnership with NGO \emph{Green Youth Cameroon}). The Terminale Technique student thus moves from learner posture to \textbf{local development actor}, using English as a technical working language and a tool for civic advocacy.
\end{savaistu}

\newpage

\coursec{Méthodes et exercices résolus}

\courssub{Substance Abuse : Description and Relative Clauses}

\begin{definition}[Key vocabulary : substance abuse]
\cle{drug abuse} (abus de drogues), \cle{addiction} (dépendance), \cle{addict} (toxicomane), \cle{tobacco} (tabac), \cle{alcohol} (alcool), \cle{peer pressure} (pression du groupe), \cle{rehabilitation centre} (centre de désintoxication), \cle{counsellor} (conseiller), \cle{to quit smoking} (arrêter de fumer), \cle{liver} (foie), \cle{overdose} (surdose), \cle{withdrawal} (sevrage).
\end{definition}

\begin{methode}[Building a complex sentence with a relative pronoun]
\textbf{Aim :} Link two simple sentences to describe a person, a thing, a place, a time or a reason related to substance abuse.
\begin{enumerate}
    \item \textbf{Identify the antecedent} (the noun common to both sentences). Ex: \textit{The teenager smokes. The teenager is my brother.} $\rightarrow$ Antecedent = \cle{the teenager}.
    \item \textbf{Choose the right relative pronoun :}
    \begin{itemize}
        \item \cle{who} / \cle{that} : person (subject) $\rightarrow$ \textit{The teenager \textbf{who} smokes...}
        \item \cle{whom} / \cle{that} / \cle{\O} : person (object) $\rightarrow$ \textit{The doctor \textbf{(whom)} I met...}
        \item \cle{which} / \cle{that} : thing/animal (subject or object) $\rightarrow$ \textit{The drug \textbf{which} destroys...}
        \item \cle{whose} : possession $\rightarrow$ \textit{The man \textbf{whose} son is addicted...}
        \item \cle{where} : place $\rightarrow$ \textit{The centre \textbf{where} they treat addicts...}
        \item \cle{when} : time $\rightarrow$ \textit{The day \textbf{when} he quit...}
        \item \cle{why} : reason $\rightarrow$ \textit{The reason \textbf{why} he started...}
    \end{itemize}
    \item \textbf{Delete the repeated noun} (or personal pronoun) in the second sentence.
    \item \textbf{Place the relative clause immediately after the antecedent.}
    \item \textbf{Punctuation :} No comma if the clause is \textbf{defining/restrictive} (essential to the meaning). Commas if \textbf{non-defining/non-restrictive} (extra information). Note : \textit{that} is never used in non-restrictive clauses.
\end{enumerate}
\textbf{Formula to remember :} \fbox{Antecedent + Relative Pronoun + Verb + Complements (no repetition)}
\end{methode}

\begin{exemplebox}[Dialogue : Talking about substance abuse]
\textit{-- Did you hear about the boy \textbf{who} was taken to hospital yesterday?}\\
\textit{-- Yes. He is the student \textbf{whose} father sells cigarettes near the school.}\\
\textit{-- The health club, \textbf{which} meets every Friday, is organising a talk on drug abuse.}\\
\textit{-- That is the reason \textbf{why} I joined it. Young people need information.}
\end{exemplebox}

\begin{exoresolu}[Joining sentences with relative pronouns (Bac type)]
\textbf{Task :} Join the following pairs of sentences using the appropriate relative pronoun (\textit{who, which, whose, where, when, why}). State whether the clause is restrictive (R) or non-restrictive (NR).
\begin{enumerate}
    \item Alcohol damages the liver. The liver is a vital organ.
    \item The counsellor helped my friend. Her father is an alcoholic.
    \item The rehabilitation centre is far from here. Addicts receive treatment there.
    \item Many young people start smoking at parties. Peer pressure is high there.
    \item The year 2020 was difficult. The pandemic increased substance abuse then.
    \item Dr. Smith gave a lecture. He is a specialist in addiction.
\end{enumerate}

\tcblower
\textbf{Detailed solution :}
\begin{enumerate}
    \item \textbf{Alcohol damages the liver, \cle{which} is a vital organ.} (NR -- \textit{the liver} is unique; the information is additional). \textit{Note :} \textit{that} is impossible in a non-restrictive clause.
    \item \textbf{The counsellor \cle{whose} father is an alcoholic helped my friend.} (R -- identifies which counsellor). \textit{Structure :} \textit{whose} + noun (father).
    \item \textbf{The rehabilitation centre \cle{where} addicts receive treatment is far from here.} (R -- \textit{where} = \textit{in which}).
    \item \textbf{Many young people start smoking at parties \cle{where} peer pressure is high.} (R -- \textit{where} = \textit{at which}).
    \item \textbf{The year 2020, \cle{when} the pandemic increased substance abuse, was difficult.} (NR -- the year is already identified).
    \item \textbf{Dr. Smith, \cle{who} is a specialist in addiction, gave a lecture.} (NR -- Dr. Smith is identified by his name).
\end{enumerate}
\textbf{Conclusion :} The choice of pronoun depends on the nature of the antecedent (person/thing/place/time/reason/possession) and on its grammatical function inside the relative clause.
\end{exoresolu}

\begin{exercice}[Practice : Relative clauses]
Join each pair of sentences with a suitable relative pronoun.
\begin{enumerate}
    \item Cigarettes contain nicotine. Nicotine causes addiction.
    \item The girl won the prize. Her essay was about alcohol abuse.
    \item This is the hospital. Drug addicts are treated there.
    \item 2015 was the year. Our school created an anti-drug club then.
\end{enumerate}
\end{exercice}

\begin{corrige}[Exercice : Relative clauses]
\begin{enumerate}
    \item Cigarettes contain nicotine \cle{which/that} causes addiction.
    \item The girl \cle{whose} essay was about alcohol abuse won the prize.
    \item This is the hospital \cle{where} drug addicts are treated.
    \item 2015 was the year \cle{when} our school created an anti-drug club.
\end{enumerate}
\end{corrige}

\courssub{Preventing Communicable Diseases : Listening and Reported Speech}

\begin{definition}[Key vocabulary : preventing communicable diseases]
\cle{communicable disease} (maladie transmissible), \cle{hygiene} (hygiène), \cle{to wash one's hands} (se laver les mains), \cle{to disinfect} (désinfecter), \cle{vaccination} (vaccination), \cle{symptom} (symptôme), \cle{fever} (fièvre), \cle{outbreak} (épidémie), \cle{quarantine} (quarantaine), \cle{health centre} (centre de santé), \cle{to boil water} (faire bouillir l'eau), \cle{contaminated} (contaminé).
\end{definition}

\begin{methode}[Converting direct speech into reported speech]
\textbf{Aim :} Report advice, warnings or information about disease prevention without quoting word for word.
\begin{enumerate}
    \item \textbf{Identify the reporting verb} (often in the past : \cle{said}, \cle{told}, \cle{advised}, \cle{warned}, \cle{explained}, \cle{asked}).
    \item \textbf{Apply the \cle{backshifting} rule} (one tense back) when the reporting verb is in the past :
    \begin{center}
    \begin{tabular}{|c|c|}\hline
    \textbf{Direct speech} & \textbf{Reported speech} \\ \hline
    Present Simple & Past Simple \\
    Present Continuous & Past Continuous \\
    Present Perfect & Past Perfect \\
    Past Simple & Past Perfect \\
    \cle{will} & \cle{would} \\
    \cle{can} & \cle{could} \\
    \cle{must} / \cle{have to} & \cle{had to} \\
    \cle{may} & \cle{might} \\ \hline
    \end{tabular}
    \end{center}
    \item \textbf{Adapt personal and possessive pronouns} (\textit{I $\rightarrow$ he/she, my $\rightarrow$ his/her, we $\rightarrow$ they, our $\rightarrow$ their}).
    \item \textbf{Adapt time and place markers :} \textit{now $\rightarrow$ then, today $\rightarrow$ that day, yesterday $\rightarrow$ the day before, tomorrow $\rightarrow$ the next day, here $\rightarrow$ there, this $\rightarrow$ that, next Monday $\rightarrow$ the following Monday}.
    \item \textbf{Structure according to sentence type :}
    \begin{itemize}
        \item \textbf{Statement :} \textit{Subject + reporting verb + (that) + clause.} (\textit{He said (that) he was sick.})
        \item \textbf{Wh-question :} \textit{Subject + reporting verb + Wh-word + subject + verb} (affirmative order). (\textit{She asked where the clinic was.})
        \item \textbf{Yes/No question :} \textit{Subject + reporting verb + if/whether + clause.} (\textit{He asked if I had a fever.})
        \item \textbf{Order/Advice :} \textit{Subject + reporting verb + object + to-infinitive.} (\textit{The nurse advised us \cle{to wash} our hands.})
    \end{itemize}
\end{enumerate}
\textbf{Exception :} No backshifting if the fact is still true at the moment of reporting (\textit{The teacher said that water \cle{boils} at 100°C}).
\end{methode}

\begin{exoresolu}[Reporting a health talk (Bac type)]
\textbf{Task :} Report the following statements made by a nurse during an educational talk in a technical school. Use the suggested reporting verbs.
\begin{enumerate}
    \item \textit{"Wash your hands regularly with soap,"} she \textbf{advised} the students.
    \item \textit{"I have disinfected the classrooms,"} she \textbf{said}.
    \item \textit{"Where is the nearest health centre?"} a student \textbf{asked}.
    \item \textit{"Do not share personal items like towels,"} she \textbf{warned} them.
    \item \textit{"The vaccination campaign will start next Monday,"} she \textbf{announced}.
\end{enumerate}

\tcblower
\textbf{Detailed solution :}
\begin{enumerate}
    \item \textbf{She \cle{advised} the students \cle{to wash} their hands regularly with soap.}
    \textit{Analysis :} Order/Advice $\rightarrow$ \textit{advised} + object + \cle{to-infinitive}. \textit{your $\rightarrow$ their}.
    \item \textbf{She \cle{said} (that) she \cle{had disinfected} the classrooms.}
    \textit{Analysis :} Statement. Present Perfect (\textit{have disinfected}) $\rightarrow$ Past Perfect (\textit{had disinfected}). \textit{I $\rightarrow$ she}.
    \item \textbf{A student \cle{asked} \cle{where} the nearest health centre \cle{was}.}
    \textit{Analysis :} Wh-question $\rightarrow$ \textit{asked} + Wh-word + subject + verb (affirmative order). \textit{is $\rightarrow$ was} (backshifting).
    \item \textbf{She \cle{warned} them \cle{not to share} personal items like towels.}
    \textit{Analysis :} Negative order $\rightarrow$ \textit{warned} + object + \cle{not to-infinitive}.
    \item \textbf{She \cle{announced} (that) the vaccination campaign \cle{would start} \cle{the following Monday}.}
    \textit{Analysis :} Future statement. \textit{will start $\rightarrow$ would start}. \textit{next Monday $\rightarrow$ the following Monday}.
\end{enumerate}
\textbf{Conclusion :} Faithfulness to the meaning matters more than word-for-word translation ; adapting deictics (persons, time, place) is crucial for the coherence of the report.
\end{exoresolu}

\begin{exercice}[Practice : Reported speech]
Report these sentences. The reporting verb is in the past.
\begin{enumerate}
    \item \textit{"Malaria is transmitted by mosquitoes,"} the doctor explained.
    \item \textit{"We will organise a clean-up day tomorrow,"} the principal said.
    \item \textit{"Have you washed your hands?"} the teacher asked the pupils.
    \item \textit{"Don't drink water from the stream,"} the nurse told the children.
\end{enumerate}
\end{exercice}

\begin{corrige}[Exercice : Reported speech]
\begin{enumerate}
    \item The doctor explained (that) malaria \cle{was} transmitted by mosquitoes. (\textit{is} also accepted : permanent fact.)
    \item The principal said (that) they \cle{would organise} a clean-up day \cle{the next day}.
    \item The teacher asked the pupils \cle{if/whether} they \cle{had washed} their hands.
    \item The nurse told the children \cle{not to drink} water from the stream.
\end{enumerate}
\end{corrige}

\courssub{Environmental Issues : Climate Change, Waste Management and Conditionals}

\begin{definition}[Key vocabulary : environment]
\cle{climate change} (changement climatique), \cle{global warming} (réchauffement climatique), \cle{recycling} (recyclage), \cle{waste disposal} (élimination des déchets), \cle{garbage collection} (ramassage des ordures), \cle{landfill} (décharge), \cle{biodegradable} (biodégradable), \cle{carbon footprint} (empreinte carbone), \cle{renewable energy} (énergie renouvelable), \cle{waste management} (gestion des déchets), \cle{pollution} (pollution), \cle{dump site} (dépôt sauvage), \cle{to sort waste} (trier les déchets).
\end{definition}

\begin{methode}[Using conditionals to express real or hypothetical consequences]
\textbf{Aim :} Discuss climate change and waste management by distinguishing the probable from the imaginary.
\begin{enumerate}
    \item \textbf{Identify the nature of the situation :}
    \begin{itemize}
        \item \textbf{Real / Possible / Probable (future) :} \cle{First Conditional}. \textit{If we recycle, we will save resources.}
        \item \textbf{Unreal / Hypothetical / Imaginary (present/future) :} \cle{Second Conditional}. \textit{If I were the Minister, I would ban plastic bags.}
    \end{itemize}
    \item \textbf{Apply the strict structure :}
    \begin{center}
    \begin{tabular}{|l|l|l|}\hline
    \textbf{Type} & \textbf{If-clause (condition)} & \textbf{Main clause (result)} \\ \hline
    \textbf{First (real)} & \cle{If + Present Simple} & \cle{will / can / may / must + V} \\ \hline
    \textbf{Second (unreal)} & \cle{If + Past Simple (\textit{were} for all persons)} & \cle{would / could / might + V} \\ \hline
    \end{tabular}
    \end{center}
    \item \textbf{Order of the clauses :} \textit{If-clause, main clause} (with a comma) OR \textit{main clause + if-clause} (no comma).
    \item \textbf{Trap :} Never put \textit{will/would} in the \textit{if-clause}. \textit{*If it will rain...} $\rightarrow$ \textbf{If it rains...}
    \item \textbf{\cle{Unless} = If... not :} \textit{Unless we act, the situation will get worse.} (= \textit{If we do not act...})
\end{enumerate}
\textbf{Bac formula :} \fbox{If + Present $\rightarrow$ Will + V} \quad | \quad \fbox{If + Past (were) $\rightarrow$ Would + V}
\end{methode}

\begin{exoresolu}[Verb forms and choice of conditional type (Bac type)]
\textbf{Task :} Put the verbs in brackets in the correct tense (First or Second Conditional). Briefly justify your choice.
\begin{enumerate}
    \item If the factory \textit{(stop)} dumping waste into the river, the water \textit{(be)} clean again.
    \item If everyone \textit{(use)} public transport, carbon emissions \textit{(decrease)} significantly.
    \item What \textit{you / do} if you \textit{(see)} someone throwing plastic in the street?
    \item Unless the government \textit{(invest)} in recycling plants, the landfill \textit{(overflow)} soon.
    \item If I \textit{(be)} you, I \textit{(start)} composting organic waste immediately.
\end{enumerate}

\tcblower
\textbf{Detailed solution :}
\begin{enumerate}
    \item \textbf{If the factory \cle{stops} dumping waste into the river, the water \cle{will be} clean again.}
    \textit{Justification :} \textbf{First Conditional}. Real, possible future situation. \textit{If + Present Simple $\rightarrow$ will + V}.
    \item \textbf{If everyone \cle{used} public transport, carbon emissions \cle{would decrease} significantly.}
    \textit{Justification :} \textbf{Second Conditional}. Hypothetical situation (not everyone does it now). \textit{If + Past Simple $\rightarrow$ would + V}.
    \item \textbf{\cle{What would you do} if you \cle{saw} someone throwing plastic in the street?}
    \textit{Justification :} \textbf{Second Conditional}. Hypothetical question about an imaginary reaction.
    \item \textbf{Unless the government \cle{invests} in recycling plants, the landfill \cle{will overflow} soon.}
    \textit{Justification :} \textbf{First Conditional}. \textit{Unless = If not}. Real threat based on a present condition.
    \item \textbf{If I \cle{were} you, I \cle{would start} composting organic waste immediately.}
    \textit{Justification :} \textbf{Second Conditional}. Classic piece of advice (\textit{If I were you}). \textit{Were} for all persons.
\end{enumerate}
\textbf{Conclusion :} The \textit{First} conditional expresses a commitment or a prediction (\textit{If we act now...}) ; the \textit{Second} expresses a wish, advice or simulation (\textit{If we acted now...}).
\end{exoresolu}

\begin{exercice}[Practice : Conditionals]
Complete with the correct form of the verbs in brackets.
\begin{enumerate}
    \item If we \textit{(plant)} more trees, the air \textit{(be)} cleaner.
    \item If I \textit{(have)} a recycling bin at home, I \textit{(sort)} my waste every day.
    \item The floods \textit{(destroy)} the crops if the rains \textit{(continue)}.
    \item If people \textit{(not / throw)} plastic everywhere, our town would be beautiful.
\end{enumerate}
\end{exercice}

\begin{corrige}[Exercice : Conditionals]
\begin{enumerate}
    \item If we \cle{plant} more trees, the air \cle{will be} cleaner. (First)
    \item If I \cle{had} a recycling bin at home, I \cle{would sort} my waste every day. (Second)
    \item The floods \cle{will destroy} the crops if the rains \cle{continue}. (First)
    \item If people \cle{did not throw} plastic everywhere, our town would be beautiful. (Second)
\end{enumerate}
\end{corrige}

\courssub{Advanced Adjective Structures : Degree and Comparison}

\begin{methode}[Mastering \textit{How + Adj}, \textit{Too...To}, \textit{Enough...To} and the Superlative]
\textbf{Aim :} Express intensity, excess/insufficiency and the supreme degree to describe environmental or health problems.
\begin{enumerate}
    \item \textbf{\cle{How + Adjective (+ subject + verb)} :} exclamation about degree.
    \begin{itemize}
        \item \textit{How \cle{serious} the pollution is!}
        \item \textit{How \cle{difficult} it is to quit smoking!}
    \end{itemize}
    \item \textbf{\cle{Too + Adjective + To-Infinitive} :} excess $\rightarrow$ impossibility (negative meaning).
    \begin{itemize}
        \item \textit{The water is \cle{too polluted to drink}.} (= so polluted that we cannot drink it).
        \item \textit{He is \cle{too weak to resist} peer pressure.}
        \item \textbf{Warning :} the subject of the main clause is the logical subject of the infinitive. If the subjects are different $\rightarrow$ \textit{Too... For + object + To}. \textit{The box is too heavy \cle{for me to lift}.}
    \end{itemize}
    \item \textbf{\cle{Adjective + Enough + To-Infinitive} :} sufficiency $\rightarrow$ possibility (positive meaning).
    \begin{itemize}
        \item \textit{The measures are \cle{strong enough to reduce} waste.}
        \item \textit{Enough} comes \textbf{after} the adjective but \textbf{before} a noun : \textit{enough bins}, \textit{enough money}.
    \end{itemize}
    \item \textbf{Superlative \cle{The + most + long adjective} / \cle{The + short adjective-est} :} comparison within a group.
    \begin{itemize}
        \item \textit{Climate change is \cle{the most urgent} challenge of our time.}
        \item \textit{Plastic is \cle{the worst} pollutant for oceans.} (irregular : \textit{bad $\rightarrow$ worse $\rightarrow$ worst} ; \textit{good $\rightarrow$ better $\rightarrow$ best})
    \end{itemize}
    \item \textbf{Correlated structure \cle{The + comparative..., the + comparative...} :}
    \begin{itemize}
        \item \textit{\cle{The more} we recycle, \cle{the less} waste we produce.}
    \end{itemize}
\end{enumerate}
\textbf{Golden rule :} \textit{Too} = negative (excess $\rightarrow$ not possible). \textit{Enough} = positive (sufficient $\rightarrow$ possible). \textit{Very} = neutral intensity only.
\end{methode}

\begin{exoresolu}[Sentence transformation and error correction (Bac type)]
\textbf{Task :}
\begin{enumerate}
    \item Rewrite with \textit{Too... To} : \textit{The air is so polluted that we cannot breathe easily.}
    \item Rewrite with \textit{How + Adjective} : \textit{It is very shocking to see the amount of waste.}
    \item Complete with the correct form (superlative or correlated comparative) : \textit{Deforestation is \_\_\_\_\_\_ (bad) environmental disaster in the region. The \_\_\_\_\_\_ (many) trees we cut, \_\_\_\_\_\_ (fast) the soil erodes.}
    \item Correct the errors : \textit{This chemical is enough dangerous to kill fish.} / \textit{He is too strong to lift the box.}
\end{enumerate}

\tcblower
\textbf{Detailed solution :}
\begin{enumerate}
    \item \textbf{The air is \cle{too polluted to breathe} easily.} (Or \textit{too polluted for us to breathe easily}). \textit{So... that + negative $\rightarrow$ Too... to}.
    \item \textbf{\cle{How shocking} it is to see the amount of waste!} Exclamation about degree.
    \item \textbf{Deforestation is \cle{the worst} environmental disaster in the region. \cle{The more} trees we cut, \cle{the faster} the soil erodes.}
    \textit{Analysis :} irregular superlative \textit{bad $\rightarrow$ worse $\rightarrow$ worst}. Correlation \textit{The more..., the faster...}.
    \item \textbf{Corrections :}
    \begin{itemize}
        \item \textit{This chemical is \cle{dangerous enough to kill} fish.} \textit{Lesson :} \textit{enough} comes \textbf{after} the adjective. \textit{*Enough dangerous} is wrong.
        \item \textit{He is \cle{not strong enough to lift} the box} (or \textit{too weak to lift it}). \textit{Too strong} would absurdly mean he is so strong that he \textbf{cannot} lift it.
    \end{itemize}
\end{enumerate}
\textbf{Conclusion :} The \textit{Too/Enough} distinction radically changes the meaning (failure vs success). The superlative requires \textit{the} + extreme form. The \textit{The... The...} correlation structures cause-effect argumentation.
\end{exoresolu}

\begin{exercice}[Practice : Degree and comparison]
\begin{enumerate}
    \item Rewrite with \textit{enough} : \textit{The water is not clean, so we cannot drink it.}
    \item Rewrite with \textit{too... to} : \textit{The rubbish heap is so high that the truck cannot pass.}
    \item Complete : \textit{Yaoundé is one of \_\_\_\_\_\_ (polluted) cities in the country.}
    \item Complete : \textit{\_\_\_\_\_\_ (early) we act, \_\_\_\_\_\_ (good) it is for the planet.}
\end{enumerate}
\end{exercice}

\begin{corrige}[Exercice : Degree and comparison]
\begin{enumerate}
    \item The water is \cle{not clean enough to drink}.
    \item The rubbish heap is \cle{too high for the truck to pass}.
    \item Yaoundé is one of \cle{the most polluted} cities in the country.
    \item \cle{The earlier} we act, \cle{the better} it is for the planet.
\end{enumerate}
\end{corrige}

\courssub{Writing and Integration : From Analysis to Action (Composition)}

\begin{methode}[Writing an argumentative composition on the environment (Bac standard)]
\textbf{Aim :} Produce a coherent text (150--200 words) on climate change or waste management, integrating the grammar of the unit (conditionals, relatives, reported speech) and specific vocabulary.
\begin{enumerate}
    \item \textbf{Analyse the topic :} identify the \cle{theme} (e.g. \textit{plastic waste}), the \cle{text type} (argumentative / problem-solution), the \cle{audience} (community, authorities) and the \cle{tone} (formal, persuasive).
    \item \textbf{Plan (5 min) :} make a \cle{three-part plan} :
    \begin{itemize}
        \item \textbf{Introduction :} hook (shocking fact, question) $\rightarrow$ definition of the problem $\rightarrow$ \cle{thesis statement} (e.g. \textit{Urgent action is needed to manage plastic waste through recycling and education.}).
        \item \textbf{Body (2--3 paragraphs) :} one main idea per paragraph.
        \begin{itemize}
            \item \textit{Paragraph 1 : Causes/Effects} (First Conditional : \textit{If we continue..., sea levels will rise.} / Relative clauses : \textit{Factories which pollute...}).
            \item \textit{Paragraph 2 : Existing solutions} (Reported Speech : \textit{Experts say that...} / Modals : \textit{We must/should...}).
            \item \textit{Paragraph 3 : Personal/collective call to action} (Second Conditional : \textit{If everyone recycled...} / \textit{Too/Enough} : \textit{It is not too late to act.}).
        \end{itemize}
        \item \textbf{Conclusion :} summary $\rightarrow$ restated thesis $\rightarrow$ \cle{opening} (future question, proverb, final appeal).
    \end{itemize}
    \item \textbf{Write (25 min) :} use \cle{linking words} (\textit{Firstly, Furthermore, Consequently, However, For instance, In conclusion}). Vary the structures (passive, conditionals, relatives).
    \item \textbf{Proofread (5 min) :} check \textbf{COPS} : \cle{C}apitalization, \cle{O}rganization, \cle{P}unctuation, \cle{S}pelling and grammar (agreements, tenses, prepositions).
\end{enumerate}
\textbf{Bac criteria :} Relevance + Coherence/Cohesion + Grammar/Vocabulary + Spelling/Punctuation.
\end{methode}

\begin{exoresolu}[Guided writing : "Managing Waste in Our Community" (Bac type)]
\textbf{Task :} \textit{Write a composition of about 180 words on the following topic : "Effective waste management is a shared responsibility. Discuss the problems caused by poor waste disposal in your community and suggest practical solutions."}

\tcblower
\textbf{Model answer (corrected) :}

\textbf{Effective Waste Management : A Shared Duty}

\textit{Introduction :} How shocking it is to see piles of garbage blocking drains in our neighbourhoods! Poor waste disposal has become a major environmental issue in many Cameroonian towns. It causes floods, spreads diseases like cholera and malaria, and pollutes the air we breathe. \textbf{Effective waste management is not solely the government's duty ; it requires active participation from every citizen.}

\textit{Body 1 -- Problems :} Firstly, indiscriminate dumping blocks drainage systems. \textbf{If heavy rains fall}, the water \textbf{will not flow}, leading to devastating floods that destroy houses and roads. Moreover, decomposing waste attracts mosquitoes and rats \textbf{which} spread communicable diseases. The air becomes \textbf{too polluted to breathe} comfortably, affecting the health of children and the elderly.

\textit{Body 2 -- Solutions :} However, solutions exist. The authorities \textbf{must provide} enough public bins and regular collection trucks. \textbf{Experts advise that} recycling plants should be built to transform plastic and organic waste into useful products. \textbf{If the government invested} in these facilities, unemployment \textbf{would decrease} too.

\textit{Body 3 -- Individual action :} Citizens also have a role. We \textbf{should sort} our waste at home : organic waste for compost, plastic for recycling. \textbf{The more} we recycle, \textbf{the cleaner} our environment becomes. It is \textbf{not too late to change} our habits.

\textit{Conclusion :} In conclusion, a clean environment is a healthy environment. Let us not wait for the next epidemic to act. As the saying goes, \textit{"Cleanliness is next to godliness."} Together, we can make our community a model of sustainability.

\textbf{Integrated grammar analysis :}
\begin{itemize}
    \item \textit{How shocking...} (exclamation \textit{How + Adj}).
    \item \textit{...which spread...} (restrictive relative clause).
    \item \textit{If heavy rains fall... will not flow} (First Conditional).
    \item \textit{too polluted to breathe} (Too + Adj + To).
    \item \textit{Experts advise that...} (reporting verb + clause).
    \item \textit{If the government invested... would decrease} (Second Conditional).
    \item \textit{The more... the cleaner} (correlated comparatives).
    \item \textit{not too late to change} (negation of Too... To).
\end{itemize}
\textbf{Conclusion :} This model respects the structure and the word limit, and naturally integrates the grammar points of the module.
\end{exoresolu}

\courssub{Integration : Listening and Speaking Simulation}

\begin{methode}[Succeeding in the oral test : health talk simulation]
\textbf{Aim :} Restitute heard information (Listening) and produce a structured oral exchange (Speaking) on prevention.
\begin{enumerate}
    \item \textbf{Before listening (preparation) :} read the questions/grid. Predict the vocabulary (\textit{symptoms, prevention, vaccine, hygiene, outbreak, quarantine, immunity}).
    \item \textbf{While listening (note-taking -- \cle{keywords only}) :} do not write full sentences. Note : \textit{Who? What? Where? When? How many? Key figures (dates, percentages).} Use abbreviations (\textit{govt, hosp, prev, b/c, w/}).
    \item \textbf{After listening (restitution / reported speech) :} turn the notes into full sentences in \textbf{reported speech}. \textit{The speaker explained that... / He warned against... / She recommended...}
    \item \textbf{Oral production (dialogue / talk) :} use the \textbf{PREP} structure : \cle{P}oint (opinion), \cle{R}eason (cause), \cle{E}xample (concrete example), \cle{P}roposal (solution).
    \item \textbf{Interaction :} use \cle{language functions} : agreeing/disagreeing politely, asking for clarification, giving advice (\textit{If I were you...}), expressing obligation (\textit{We must...}).
\end{enumerate}
\textbf{Self-assessment grid :} fluency / pronunciation / grammatical accuracy (conditionals, relatives, reported speech) / lexical richness / time management.
\end{methode}

\begin{exoresolu}[Oral presentation simulation : "Preventing Cholera in Schools" (Bac oral type)]
\textbf{Task :} You are the class representative of a Terminale class. Using the notes below, give a two-minute talk to sensitise your classmates on cholera prevention. Use reported speech to quote the district nurse.

\textbf{Notes (simulated listening) :}
\begin{itemize}
    \item District Hospital nurse : "Cholera spreads via contaminated water/food."
    \item "Wash hands with soap before eating/after the toilet."
    \item "Boil drinking water / use chlorine tablets."
    \item "Symptoms : acute watery diarrhoea, vomiting $\rightarrow$ dehydration $\rightarrow$ risk of death."
    \item "Report cases immediately $\rightarrow$ treatment is free."
    \item Statistics : 15 cases last month in the district ; 2 deaths (late reporting).
\end{itemize}

\tcblower
\textbf{Structured model talk (corrected) :}

\textbf{"Good morning Principal, Teachers, and fellow students.}

\textbf{Introduction :} \textit{As the class representative, I want to share vital information from the district nurse regarding the cholera outbreak in our area.}

\textbf{Reported facts :} \textit{The nurse \cle{explained that} cholera \cle{spreads} through contaminated water and food. She \cle{warned} that the symptoms include acute watery diarrhoea and vomiting, \cle{which} lead to severe dehydration and can cause death if untreated.}

\textbf{Statistics :} \textit{She \cle{revealed that} there \cle{had been} 15 cases in our district last month, and sadly, 2 people \cle{had died} because they \cle{had reported} late.}

\textbf{Advice / Prevention :} \textit{She \cle{advised} us \cle{to wash} our hands with soap before eating and after using the toilet. She \cle{recommended boiling} drinking water or using chlorine tablets.}

\textbf{Call to action :} \textit{\cle{If we follow} these rules, we \cle{will protect} ourselves and our families. \cle{It is not too late to act}. Remember : treatment is free, so \cle{do not hesitate to go} to the health centre immediately if you see symptoms.}

\textbf{Conclusion :} \textit{Let us keep our school clean and safe. Thank you."}

\textbf{Targeted grammar items :}
\begin{itemize}
    \item \textit{explained that / warned that / revealed that / advised us to / recommended + V-ing} (varied reported speech).
    \item \textit{which lead to} (relative clause).
    \item \textit{had been / had died / had reported} (backshifting to Past Perfect).
    \item \textit{If we follow... will protect} (First Conditional -- real possibility).
    \item \textit{not too late to act} (Too... To structure).
\end{itemize}
\end{exoresolu}

\courssub{Reading and Vocabulary : Climate Change and Waste Management Text Analysis}

\begin{methode}[Active reading strategies for technical texts (Bac reading comprehension)]
\textbf{Aim :} Extract explicit and implicit information from a text on the environment (climate change, recycling) and mobilise specialised vocabulary.
\begin{enumerate}
    \item \textbf{Survey :} read the \textbf{title}, \textbf{subtitles}, \textbf{first and last paragraphs}. Identify the \cle{genre} (scientific article, report, blog, letter) and the \cle{purpose} (to inform, to alert, to persuade).
    \item \textbf{Question :} turn the headings into questions (\textit{What are the causes of...? How does recycling work?}).
    \item \textbf{Active reading :}
    \begin{itemize}
        \item \textit{Skimming :} fast reading for the general idea (main idea).
        \item \textit{Scanning :} targeted search (keywords, figures, proper nouns, \textbf{bold/italic words}).
    \end{itemize}
    \item \textbf{Dealing with unknown vocabulary (inference) :}
    \begin{itemize}
        \item Context : \textit{The \cle{landfill} was full, so they opened a new dump site.} $\rightarrow$ \textit{landfill} = décharge.
        \item Word formation : \textit{\cle{bio}degradable, \cle{re}cycle, \cle{non}-renewable}.
        \item Cognates (English/French) : \textit{pollution, environment, atmosphere, carbon, deforestation, erosion, biodiversity}.
    \end{itemize}
    \item \textbf{Typical Bac question types :}
    \begin{itemize}
        \item \textit{Right/Wrong/Not mentioned} (justify with a quotation).
        \item \textit{Wh-questions} (short answers, words from the text).
        \item \textit{Vocabulary matching / synonyms in context}.
        \item \textit{Grammar in context} (find a conditional, a relative clause, a passive).
        \item \textit{Summary / title choice}.
    \end{itemize}
\end{enumerate}
\textbf{Golden rule :} The answer is \textbf{always in the text}. Never bring outside knowledge for literal comprehension questions.
\end{methode}

\begin{exoresolu}[Reading comprehension : "The Plastic Crisis" (Bac type)]
\textbf{Task :} Read the extract and answer the questions.
\textit{"Global plastic production has doubled since 2000, reaching about 460 million tonnes in 2019. Only about 9\% is successfully recycled. The majority ends up in \cle{landfills} or \cle{leaches} into rivers and oceans, threatening marine \cle{biodiversity}. Microplastics have now entered the human food chain. \textbf{If current trends continue}, there \textbf{will be} more plastic than fish in the sea by 2050. Governments \textbf{must implement} Extended Producer Responsibility (EPR) schemes, \textbf{which} make manufacturers responsible for the end-of-life of their products. Citizens \textbf{should reduce} single-use plastics."}

\textbf{Questions :}
\begin{enumerate}
    \item Find the sentence which expresses a \textbf{prediction based on a real condition} (First Conditional). Quote it.
    \item Find the \textbf{relative clause} defining "EPR schemes". Quote the clause only.
    \item Explain the meaning of \textit{"leaches"} in context. (Choose : A) flows slowly / B) burns quickly / C) disappears).
    \item Rewrite using \textit{Too... To} : \textit{The amount of waste is so huge that we cannot ignore it.}
    \item What does the pronoun \textit{"their"} refer to in \textit{"their products"}?
\end{enumerate}

\tcblower
\textbf{Detailed solution :}
\begin{enumerate}
    \item \textbf{Quotation :} \textit{"\textbf{If current trends continue}, there \textbf{will be} more plastic than fish in the sea by 2050."}
    \textit{Justification :} structure \textit{If + Present Simple (continue), ... will + be}. A real prediction.
    \item \textbf{Clause :} \textit{"\textbf{which make manufacturers responsible for the end-of-life of their products}"}.
    \textit{Justification :} restrictive relative clause (no comma) introduced by \textit{which} (antecedent \textit{EPR schemes} = thing), subject of the verb \textit{make}.
    \item \textbf{Answer : A) flows slowly} (s'infiltre, s'écoule).
    \textit{Justification :} context \textit{"ends up in landfills OR leaches into rivers"} : parallel end-points in water.
    \item \textbf{Rewrite :} \textit{"The amount of waste is \textbf{too huge to ignore}."} (or \textit{too huge for us to ignore}).
    \item \textbf{Reference :} \textit{"their"} refers to \textbf{manufacturers} (nearest plural antecedent, owner of the products).
\end{enumerate}
\textbf{Conclusion :} Active reading combines grammatical spotting (conditionals, relatives) and contextual lexical inference.
\end{exoresolu}

\begin{attention}[The 5 mistakes that cost marks at the Bac]
\begin{enumerate}
    \item \textbf{Conditionals : "If + will/would".} \textit{*If it will rain, I stay home.} $\rightarrow$ \textbf{If it rains, I will stay home.} (Never a modal in the \textit{if-clause} for the First/Second conditional.)
    \item \textbf{Reported speech : forgetting backshifting or deictics.} \textit{*He said he is sick yesterday.} $\rightarrow$ \textbf{He said he was sick the day before.} (Tense + \textit{yesterday $\rightarrow$ the day before}.)
    \item \textbf{Relative clauses : doubling the subject/object.} \textit{*The man who he is talking...} $\rightarrow$ \textbf{The man who is talking...} / \textit{*The book which I read it...} $\rightarrow$ \textbf{The book which I read...} (Delete the redundant pronoun.)
    \item \textbf{Too / Enough : position and meaning.} \textit{*He is enough strong.} $\rightarrow$ \textbf{He is strong enough.} \textbf{Rule :} \textit{Adjective + enough} (positive) vs \textit{too + adjective} (negative/impossible). Use \textit{for + object} when the subjects differ.
    \item \textbf{Writing : off-topic / word count / no structure.}
    \begin{itemize}
        \item Writing 80 words instead of 180 = heavy penalty.
        \item No visible paragraphs = loss of coherence marks.
        \item Copying the topic as the introduction = no relevance mark.
        \item Forgetting linkers (\textit{However, Therefore, For example}) = choppy text.
    \end{itemize}
\end{enumerate}
\textbf{Final tip :} proofread your paper \textbf{backwards} (last sentence $\rightarrow$ first) to check only grammar and spelling without being distracted by the meaning.
\end{attention}

\begin{exercice}[Integration task : end-of-unit evaluation]
\textbf{Situation :} Your school is celebrating "Health and Environment Day". As a member of the English Club, you must contribute.
\begin{enumerate}
    \item \textbf{Speaking :} prepare a one-minute dialogue with a partner about the dangers of tobacco and alcohol among students (use at least two relative clauses).
    \item \textbf{Grammar :} report this sentence : \textit{"Our school will build an incinerator next year," the principal announced.}
    \item \textbf{Writing :} write a 150-word composition entitled \textit{"If everyone recycled, our town would change"} using at least one First Conditional, one Second Conditional, one relative clause and one \textit{too...to} structure.
\end{enumerate}
\end{exercice}

\begin{corrige}[Integration task : elements of correction]
\begin{enumerate}
    \item \textbf{Possible dialogue :} \textit{-- Do you know students who smoke behind the classrooms? -- Yes, and I know a boy whose health has been destroyed by alcohol. -- That is the reason why our club organises this campaign.}
    \item \textbf{Reported speech :} The principal announced (that) their school \cle{would build} an incinerator \cle{the following year}.
    \item \textbf{Writing :} expected elements -- clear three-part plan ; \textit{If everyone recycled, our town would be cleaner} (Second Conditional) ; \textit{If we do nothing, pollution will increase} (First Conditional) ; \textit{the waste which litters our streets} (relative) ; \textit{the problem is too serious to ignore} (too...to) ; linkers ; correct spelling and punctuation.
\end{enumerate}
\end{corrige}

\newpage

\coursec{Exercices}

\courssub{Key Vocabulary \& Grammar Reminders}

\begin{aretenir}[Relative Clauses]
\textbf{Defining relative clauses} give essential information about the noun. No commas.
\begin{itemize}
\item \textbf{Who} / \textbf{that} : people (The nurse \textbf{who} works here is kind.)
\item \textbf{Which} / \textbf{that} : things/animals (The waste \textbf{which} pollutes the river is plastic.)
\item \textbf{Whose} : possession (The students \textbf{whose} school won the prize planted trees.)
\item \textbf{Where} : places (Douala is a city \textbf{where} waste collection is a problem.)
\end{itemize}
\textit{Note :} \emph{That} can replace \emph{who/which} in defining clauses. Never use \emph{that} after a preposition or in non-defining clauses (with commas).
\end{aretenir}

\begin{aretenir}[Reported Speech -- Main Changes]
\begin{center}
\begin{tabularx}{\linewidth}{|l|X|}\hline
\textbf{Direct Speech} & \textbf{Reported Speech} \\ \hline
Present Simple & Past Simple \\
Present Continuous & Past Continuous \\
Present Perfect & Past Perfect \\
Past Simple & Past Perfect \\
\textbf{will} & \textbf{would} \\
\textbf{can} & \textbf{could} \\
\textbf{must} & \textbf{had to} \\
\textbf{now} & \textbf{then} \\
\textbf{today} & \textbf{that day} \\
\textbf{tomorrow} & \textbf{the next day} \\
\textbf{yesterday} & \textbf{the day before} \\
\textbf{here} & \textbf{there} \\ \hline
\end{tabularx}
\end{center}
\textbf{Imperative} $\to$ \emph{told/ordered/advised + object + to-infinitive} (negative: \emph{not to}).
\textbf{Questions} $\to$ \emph{asked + if/whether} (yes/no) or \emph{wh-word} (open).
\end{aretenir}

\begin{aretenir}[Conditionals : First vs Second]
\begin{center}
\begin{tabularx}{\linewidth}{|l|X|X|}\hline
\textbf{Type} & \textbf{Structure} & \textbf{Use / Meaning} \\ \hline
\textbf{First (Real)} & If + \textbf{Present Simple}, \textbf{will} + \textbf{base verb} & Real possibility in the future. \\
\textbf{Second (Unreal)} & If + \textbf{Past Simple}, \textbf{would} + \textbf{base verb} & Hypothetical, unlikely, or impossible present/future. \\ \hline
\end{tabularx}
\end{center}
\textit{Key tip :} \emph{If I were you} (not \emph{was}) for advice. \emph{Were} is used for all persons in formal grammar.
\end{aretenir}

\begin{aretenir}[Adjective Structures : Degree \& Comparison]
\begin{itemize}
\item \textbf{Too + adjective + to-infinitive} : Excess, negative result. (\emph{The water is too dirty to drink.})
\item \textbf{Adjective + enough + to-infinitive} : Sufficiency, positive result. (\emph{The campaign was effective enough to reach 5000 students.})
\item \textbf{How + adjective} : Exclamation or embedded question about degree. (\emph{How deep is the well?} / \emph{I was surprised at how high the heap was.})
\item \textbf{Superlative} : \emph{The most + long adj} / \emph{The + short adj + -est}. With \emph{one of} $\to$ \emph{one of the most dangerous diseases}.
\item \textbf{So + adjective + that} : Cause/Consequence (often replaces \emph{too...to}). (\emph{The dump is so dangerous that children cannot play near it.})
\end{itemize}
\end{aretenir}

\begin{definition}[Key Vocabulary -- Environment, Health \& Waste Management]
\begin{itemize}
\item \textbf{Landfill} (\emph{décharge contrôlée}) : A site for the disposal of waste materials by burial.
\item \textbf{Leachate} (\emph{lixiviat / jus de décharge}) : Toxic liquid produced when water filters through waste.
\item \textbf{Biodegradable} (\emph{biodégradable}) : Capable of being decomposed by bacteria or other living organisms.
\item \textbf{Compost} (\emph{compost}) : Decayed organic material used as fertiliser.
\item \textbf{Recycling} (\emph{recyclage}) : Process of converting waste into reusable material.
\item \textbf{Incineration} (\emph{incinération}) : Destruction of waste by burning.
\item \textbf{Methane} (\emph{méthane}) : Gas produced by anaerobic decomposition of organic waste (greenhouse gas).
\item \textbf{Circular economy} (\emph{économie circulaire}) : Model where waste becomes a resource (reduce, reuse, recycle).
\item \textbf{Sustainability} (\emph{durabilité}) : Meeting present needs without compromising future generations.
\item \textbf{Communicable disease} (\emph{maladie transmissible}) : Illness caused by infectious agents transmitted between people/animals.
\item \textbf{Rehabilitation centre} (\emph{centre de réadaptation}) : Facility helping people recover from addiction.
\item \textbf{Substance abuse} (\emph{abus de substances}) : Harmful use of drugs, alcohol, or tobacco.
\end{itemize}
\end{definition}

\courssub{Je vérifie mes connaissances}

\begin{exercice}[QCM : Grammar and vocabulary check]
Choose the correct answer (a, b or c).

\begin{enumerate}
\item The man \textbf{......} works at the health centre is my uncle.
\begin{itemize}
\item[a)] which \item[b)] who \item[c)] whose
\end{itemize}
\item ``I will clean the school tomorrow,'' the prefect said. The prefect said that he ...... the school ...... .
\begin{itemize}
\item[a)] will clean / tomorrow \item[b)] would clean / the next day \item[c)] cleaned / that day
\end{itemize}
\item If we recycle plastic, we ...... pollution.
\begin{itemize}
\item[a)] will reduce \item[b)] would reduce \item[c)] reduced
\end{itemize}
\item A material that can decompose naturally is called ...... .
\begin{itemize}
\item[a)] biodegradable \item[b)] toxic \item[c)] synthetic
\end{itemize}
\item The water is ...... dirty ...... drink.
\begin{itemize}
\item[a)] too / to \item[b)] so / that \item[c)] enough / to
\end{itemize}
\end{enumerate}
\end{exercice}

\begin{exercice}[True or False --- justify your answer]
Say whether each statement is True or False and justify with one sentence in English.

\begin{enumerate}
\item A relative clause introduced by \emph{whose} expresses possession.
\item In reported speech, ``today'' becomes ``tomorrow''.
\item The second conditional describes a real and probable situation.
\item Tobacco, alcohol and drugs are examples of substances that can be abused.
\item \emph{How + adjective} is used to ask about the degree or measurement of something (e.g. \emph{How tall is the building?}).
\end{enumerate}
\end{exercice}

\begin{exercice}[Definitions]
Define the following words in one complete English sentence each.

\begin{enumerate}
\item landfill
\item recycling
\item communicable disease
\item global warming
\item rehabilitation centre
\end{enumerate}
\end{exercice}

\begin{exercice}[Quick transformations]
\begin{enumerate}
\item Complete with the correct relative pronoun: ``The river ...... crosses our village is polluted.''
\item Put into reported speech: ``Wash your hands regularly,'' the doctor told us.
\item Complete with the superlative: ``Cholera is one of ...... (dangerous) diseases in the rainy season.''
\item Rewrite with \emph{too...to}: ``The bin is very small. It cannot contain all our waste.''
\item Complete the conditional: ``If I were the mayor, I ...... (build) more waste treatment plants.''
\end{enumerate}
\end{exercice}

\courssub{Je m'entraîne}

\begin{exercice}[Relative clauses in context]
Join each pair of sentences using a relative pronoun (\emph{who, which, whose, where}).

\begin{enumerate}
\item The nurse works at the district hospital. She vaccinates children every week.
\item Douala is a city. Waste collection is a serious problem there.
\item The students planted trees. Their school won the environmental prize.
\item Plastic bags block the drains. They cause floods in the neighbourhood.
\item The campaign started last Monday. It sensitises youths against drug abuse.
\end{enumerate}
\end{exercice}

\begin{exercice}[Reported speech practice]
Report the following sentences. Begin as indicated.

\begin{enumerate}
\item ``You must boil water before drinking it,'' the nurse said. \\
\emph{The nurse advised us ......}
\item ``We collected two tons of plastic yesterday,'' the club president said. \\
\emph{The club president said that ......}
\item ``Do you sort your waste at home?'' the journalist asked me. \\
\emph{The journalist asked me ......}
\item ``Don't throw garbage into the river,'' the mayor ordered. \\
\emph{The mayor ordered the villagers ......}
\item ``I am organising a clean-up campaign next week,'' she announced. \\
\emph{She announced that ......}
\end{enumerate}
\end{exercice}

\begin{exercice}[Conditionals : first and second]
Complete each sentence with the correct form of the verb in brackets, then say whether it is a first or a second conditional.

\begin{enumerate}
\item If the government builds more landfills, pollution ...... (decrease).
\item If I were you, I ...... (stop) smoking immediately.
\item We ...... (not / suffer) from cholera if we drank clean water.
\item If students plant trees this year, the school ...... (become) greener.
\item If factories treated their waste, our rivers ...... (not / be) so dirty.
\end{enumerate}
\end{exercice}

\begin{exercice}[Degree and comparison]
\begin{enumerate}
\item Rewrite using \emph{too...to}: ``The patient is very weak. He cannot walk alone.''
\item Rewrite using \emph{enough...to}: ``The campaign was very effective. It reached 5000 students.''
\item Ask a question with \emph{How + adjective} about the depth of the well.
\item Complete with the superlative of the adjective in brackets: ``Malaria is ...... (common) communicable disease in our region.''
\item Rewrite using \emph{how + adjective}: ``I was surprised by the height of the rubbish heap.'' (Begin: \emph{I was surprised at how ......})
\end{enumerate}
\end{exercice}

\courssub{Je me prépare au Bac}

\begin{exercice}[Type Bac -- Grammar : transformations]
Rephrase each sentence as instructed. Do not change the meaning. (2 marks each)

\begin{enumerate}
\item The nurse said: ``You must boil water before drinking it.'' \\
(Begin: \emph{The nurse advised ......})
\item ``Plastic bags block drains. This causes floods.'' \\
(Join the two sentences using a relative pronoun.)
\item ``We don't have enough bins. We cannot sort waste.'' \\
(Rewrite using \emph{too...to}.)
\item ``If we don't act now, the situation will get worse,'' the expert warned. \\
(Put the underlined part into reported speech. Begin: \emph{The expert warned that ......})
\item ``The dump is extremely dangerous. Children cannot play near it.'' \\
(Rewrite using \emph{so...that}.)
\item ``Burning waste in the open air is forbidden,'' the mayor said. \\
(Begin: \emph{The mayor said that ......})
\end{enumerate}
\end{exercice}

\begin{exercice}[Type Bac -- Vocabulary : cloze test]
Read the text below and fill in the \textbf{ten gaps} with the correct word chosen from the box. Each word is used only once. (1 mark each)

\begin{center}
\begin{tabularx}{\linewidth}{|X|X|X|X|X|}
\hline
\rowcolor{popblueL} landfill & biodegradable & sorting & recycling & methane \\
\hline
\rowcolor{popblueL} leachate & incineration & compost & circular economy & sustainability \\
\hline
\end{tabularx}
\end{center}

\textbf{Waste Management in Douala}

Every day, the city of Douala produces more than 1,500 tons of household waste. Most of it ends up in an open \textbf{(1) ......} at the outskirts of the city, where it pollutes the soil and the air. When it rains, a toxic black liquid called \textbf{(2) ......} flows out of the dump and contaminates the groundwater.

To solve this problem, the city council wants to promote the \textbf{(3) ......} of waste at home: families should separate \textbf{(4) ......} materials, such as food leftovers, from plastics and metals. Organic waste can be transformed into \textbf{(5) ......}, a natural fertiliser that farmers can use in their fields. Plastics, glass and paper can be sent to \textbf{(6) ......} factories, where they are transformed into new products.

Some experts also suggest the \textbf{(7) ......} of non-recyclable waste in modern ovens, although this method can release dangerous gases. In addition, the gas \textbf{(8) ......}, which is produced by decomposing waste, could be captured and used to produce electricity.

All these measures are part of a new model called the \textbf{(9) ......}, in which waste becomes a resource instead of a problem. According to the mayor, this is the only way to guarantee the \textbf{(10) ......} of the city for future generations.
\end{exercice}

\begin{exercice}[Type Probatoire -- Reading comprehension]
Read the following text carefully and answer the questions in complete sentences.

\textbf{The Green Sahel Initiative}

In the far north of Cameroon, the advance of the desert threatens the lives of thousands of families. Every year, the dry season becomes longer, the rains become rarer, and the fertile land slowly turns into sand. To fight this phenomenon, the government, together with local communities and international partners, launched the Green Sahel Initiative in 2019.

The objective of the project is ambitious: to plant ten million trees in ten years along the northern border of the country. These trees, mostly acacias and neems, were chosen because they can survive with very little water. Their roots hold the soil together and prevent the wind from carrying it away. Moreover, their leaves provide shade for crops and food for animals.

But the initiative is not only about planting trees. In each village, young people are trained to manage nurseries, to dig wells and to build small dams that collect rainwater. Women receive solar cookers, which reduce the need for firewood and therefore protect the few remaining forests. ``Before, I walked five kilometres every day to find wood,'' explains Aissatou, a mother of four from Maroua. ``Now I cook with the sun, and I use my free time to sell vegetables at the market.''

The results are already visible. According to the Ministry of Environment, more than three million trees have been planted so far, and the survival rate has reached 70 per cent. Some villages that were almost abandoned five years ago are now green again, and children have returned to school because their parents no longer need them to fetch wood and water.

However, challenges remain. Climate change continues to accelerate, and the funds promised by international donors arrive slowly. ``We cannot win this battle alone,'' admits the national coordinator of the project. ``If the whole world does not reduce its emissions, our trees will not be enough. But we have the duty to try, for our children and for our land.''

\textbf{Questions}

\begin{enumerate}
\item Name two problems caused by the advance of the desert in the far north of Cameroon. (\emph{literal})
\item Why were acacias and neems chosen for the project? Give two reasons. (\emph{literal})
\item What does the word ``nurseries'' (paragraph 3) mean in this context? (\emph{inferential -- vocabulary in context})
\item What is the author's attitude towards the Green Sahel Initiative: optimistic, pessimistic or balanced? Justify with one element from the text. (\emph{inferential})
\item Give a suitable title to paragraph 4. (\emph{synthesis})
\end{enumerate}
\end{exercice}

\begin{exercice}[Type Bac -- Writing : Composition]
Choose \textbf{ONE} of the following topics and write a composition of about 150--180 words. Pay attention to paragraphing, linking words, and the grammar structures studied in this unit (relative clauses, conditionals, reported speech, adjective structures).

\begin{enumerate}
\item \textbf{Topic A:} ``Waste management is everyone's responsibility.'' Write an article for your school magazine explaining the problems of waste in your community (Douala, Yaoundé, or your local town) and proposing concrete solutions (sorting, recycling, composting, sensitisation). Use at least \textbf{one first conditional}, \textbf{one second conditional}, and \textbf{one relative clause}.
\item \textbf{Topic B:} You attended a talk by a health expert on ``The Dangers of Substance Abuse among Youths''. Write a report for your principal summarising the main points: types of substances (drugs, alcohol, tobacco), health consequences, and prevention strategies. Use \textbf{reported speech} to relay what the expert said (e.g., \emph{The expert warned that...}, \emph{He advised students to...}).
\end{enumerate}
\end{exercice}

\courssub{Corrigés}



\newpage

\coursec{Corrigés des exercices}

\courssub{Corrigés détaillés}

\begin{corrige}[Exercice 1 --- QCM : Grammar and vocabulary check]
\begin{enumerate}
\item \textbf{b) who} --- The noun \emph{the man} refers to a \cle{person (personne)}, so the relative pronoun is \cle{who (qui)} (or \cle{that (que/qui)}). \emph{Which} is for \cle{things (choses)}, \emph{whose} expresses \cle{possession (possession)}.
\item \textbf{b) would clean / the next day} --- In \cle{reported speech (discours rapporté)}: \emph{will} $\to$ \emph{would} and \emph{tomorrow} $\to$ \emph{the next day (le lendemain)}.
\item \textbf{a) will reduce} --- \cle{First conditional (première conditionnelle)}: \emph{If + Present Simple, will + base verb} (real future possibility / possibilité réelle future).
\item \textbf{a) biodegradable} --- By definition, a material that \cle{decomposes naturally (se décompose naturellement)} (by bacteria or living organisms) is \cle{biodegradable (biodégradable)}.
\item \textbf{a) too / to} --- Structure \cle{too + adjective + to-infinitive (trop + adjectif + infinitif)} expresses \cle{excess with a negative result (excès avec résultat négatif)}: \emph{too dirty to drink (trop sale pour être bu)}.
\end{enumerate}
\end{corrige}

\begin{corrige}[Exercice 2 --- True or False]
\begin{enumerate}
\item \textbf{True.} \emph{Whose} replaces a \cle{possessive adjective (adjectif possessif)} (his, her, their) and expresses \cle{possession (possession)}: \emph{The students whose school won the prize planted trees.}
\item \textbf{False.} In \cle{reported speech (discours rapporté)}, ``today'' becomes ``\cle{that day (ce jour-là)}'' (``tomorrow'' becomes ``\cle{the next day (le lendemain)}'').
\item \textbf{False.} The \cle{second conditional (deuxième conditionnelle)} describes an \emph{unreal, hypothetical or unlikely} situation (\emph{If + Past Simple, would + base verb}); the \cle{first conditional (première conditionnelle)} describes a \cle{real, probable (réelle, probable)} one.
\item \textbf{True.} \cle{Substance abuse (toxicomanie / abus de substances)} means the \cle{harmful use (usage nocif)} of \cle{psychoactive substances (substances psychoactives)}, which include \cle{tobacco (tabac)}, \cle{alcohol (alcool)} and \cle{illicit drugs (drogues illicites)}.
\item \textbf{True.} \cle{How + adjective} asks about \cle{degree or measurement (degré ou mesure)}: \emph{How tall is the building? How deep is the well?}
\end{enumerate}
\end{corrige}

\begin{corrige}[Exercice 3 --- Definitions (suggested answers / réponses suggérées)]
\begin{enumerate}
\item A \cle{landfill (décharge / site d'enfouissement)} is a site where \cle{waste materials (déchets)} are \cle{buried in the ground (enterrés)} for \cle{disposal (élimination)}.
\item \cle{Recycling (recyclage)} is the \cle{process (processus)} of \cle{collecting (collecter)} used materials and \cle{transforming them into (les transformer en)} new products instead of \cle{throwing them away (les jeter)}.
\item A \cle{communicable disease (maladie transmissible / contagieuse)} is an \cle{illness (maladie)} caused by \cle{infectious agents (agents infectieux)} (bacteria or viruses) that can be \cle{transmitted (transmise)} from one person or animal to another.
\item \cle{Global warming (réchauffement climatique)} is the \cle{long-term rise (augmentation à long terme)} in the \cle{average temperature (température moyenne)} of the Earth's climate, mainly caused by \cle{human activities (activités humaines)} such as \cle{burning fossil fuels (combustion de combustibles fossiles)}.
\item A \cle{rehabilitation centre (centre de réadaptation / de désintoxication)} is a \cle{facility (structure / établissement)} where people \cle{receive treatment and support (reçoivent un traitement et un soutien)} to \cle{recover from (guérir de / se remettre de)} drug, alcohol or tobacco \cle{addiction (dépendance / toxicomanie)}.
\end{enumerate}
\emph{Point de vigilance :} a definition must be a \cle{complete sentence (phrase complète)}, not a translation; accept any grammatically correct sentence that gives the \cle{essential meaning (sens essentiel)}.
\end{corrige}

\begin{corrige}[Exercice 4 --- Quick transformations]
\begin{enumerate}
\item The river \textbf{which/that} crosses our village is polluted. (\emph{Thing (chose)} $\to$ \emph{which/that}.)
\item The doctor told us \textbf{to wash our hands regularly}. (Imperative $\to$ \emph{told + object + to-infinitive}.)
\item Cholera is one of \textbf{the most dangerous} diseases in the rainy season. (Long adjective $\to$ \emph{the most + adjective}; note the \cle{plural (pluriel)} ``diseases'' after \emph{one of}.)
\item The bin is \textbf{too small to contain} all our waste.
\item If I were the mayor, I \textbf{would build} more waste treatment plants. (\cle{Second conditional (deuxième conditionnelle)}: \emph{If + Past Simple, would + base verb}.)
\end{enumerate}
\end{corrige}

\begin{corrige}[Exercice 5 --- Relative clauses in context]
\begin{enumerate}
\item The nurse \textbf{who/that} works at the district hospital vaccinates children every week. (\emph{Person (personne)} $\to$ \emph{who}.)
\item Douala is a city \textbf{where} waste collection is a serious problem. (\emph{Place (lieu)} $\to$ \emph{where} = \emph{in which}.)
\item The students \textbf{whose} school won the environmental prize planted trees. (\emph{Possession} $\to$ \emph{whose} = their school.)
\item Plastic bags \textbf{which/that} block the drains cause floods in the neighbourhood. (\emph{Things (choses)} $\to$ \emph{which/that}.)
\item The campaign \textbf{which/that} started last Monday sensitises youths against drug abuse. (\emph{Thing (chose)} $\to$ \emph{which/that}.)
\end{enumerate}
\emph{Point de vigilance :} these are \textbf{defining relative clauses (propositions relatives restrictives / définies)} (essential information / information essentielle), so \textbf{no commas (pas de virgules)} are used, and \emph{that} is possible for people and things.
\end{corrige}

\begin{corrige}[Exercice 6 --- Reported speech practice]
\begin{enumerate}
\item The nurse advised us \textbf{to boil water before drinking it}. (Imperative $\to$ \emph{advised + object + to-infinitive}; \emph{must} disappears in the infinitive form.)
\item The club president said that \textbf{they had collected two tons of plastic the day before}. (Past Simple $\to$ \cle{Past Perfect (plus-que-parfait)}; \emph{yesterday} $\to$ \emph{the day before (la veille)}; \emph{we} $\to$ \emph{they}.)
\item The journalist asked me \textbf{if/whether I sorted my waste at home}. (Yes/no question $\to$ \emph{asked + if/whether}; \cle{statement word order (ordre affirmatif)}, no question mark; Present $\to$ Past.)
\item The mayor ordered the villagers \textbf{not to throw garbage into the river}. (Negative imperative $\to$ \emph{ordered + object + not to-infinitive}.)
\item She announced that \textbf{she was organising a clean-up campaign the following week}. (Present Continuous $\to$ \cle{Past Continuous (passé progressif)}; \emph{next week} $\to$ \emph{the following week (la semaine suivante)}.)
\end{enumerate}
\emph{Point de vigilance :} in \cle{reported questions (questions rapportées)}, keep the \textbf{affirmative word order (ordre affirmatif)} (\emph{asked me if I sorted}, never \emph{if did I sort}).
\end{corrige}

\begin{corrige}[Exercice 7 --- Conditionals : first and second]
\begin{enumerate}
\item If the government builds more landfills, pollution \textbf{will decrease}. --- \textbf{First conditional (première conditionnelle)} (real future possibility / possibilité réelle future).
\item If I were you, I \textbf{would stop} smoking immediately. --- \textbf{Second conditional (deuxième conditionnelle)} (hypothetical advice / conseil hypothétique; note \emph{were}, not \emph{was}).
\item We \textbf{would not suffer (wouldn't suffer)} from cholera if we drank clean water. --- \textbf{Second conditional (deuxième conditionnelle)} (unreal present / présent irréel: we do not drink clean water).
\item If students plant trees this year, the school \textbf{will become} greener. --- \textbf{First conditional (première conditionnelle)} (real future possibility / possibilité réelle future).
\item If factories treated their waste, our rivers \textbf{would not be (wouldn't be)} so dirty. --- \textbf{Second conditional (deuxième conditionnelle)} (unreal present / présent irréel: factories do not treat their waste).
\end{enumerate}
\emph{Point de vigilance :} never use \emph{will} or \emph{would} in the \emph{if}-clause: \emph{If the government builds...}, not \emph{If the government will build...}.
\end{corrige}

\begin{corrige}[Exercice 8 --- Degree and comparison]
\begin{enumerate}
\item The patient is \textbf{too weak to walk} alone. (\cle{Too + adjective + to-infinitive}.)
\item The campaign was \textbf{effective enough to reach} 5000 students. (\cle{Adjective + enough + to-infinitive (adjectif + enough + infinitif)} --- \emph{enough} comes \textbf{after} the adjective.)
\item \textbf{How deep is the well?} (\cle{How + adjective} + \cle{inverted word order (inversion)} for a \cle{direct question (question directe)}.)
\item Malaria is \textbf{the most common} communicable disease in our region. (\cle{Superlative (superlatif)} of a \cle{long adjective (adjectif long)}: \emph{the most + adjective}.)
\item I was surprised at \textbf{how high the rubbish heap was}. (Embedded exclamation / structure enchevêtrée: \emph{how + adjective} + \textbf{affirmative word order (ordre affirmatif)}, no inversion.)
\end{enumerate}
\emph{Point de vigilance :} contrast item 3 (direct question, inversion: \emph{How deep is...?}) with item 5 (embedded structure / structure enchevêtrée, no inversion: \emph{...how high the heap was}).
\end{corrige}

\begin{corrige}[Exercice 9 --- Type Bac : Grammar transformations]
\textbf{Barème indicatif : 2 marks par item (1 mark for the correct structure, 1 mark for accuracy and unchanged meaning). Total : 12 marks.}
\begin{enumerate}
\item The nurse advised \textbf{us to boil water before drinking it}. \\
\emph{Attendu :} imperative $\to$ \emph{advised + object + to-infinitive}. Accept also \emph{advised that we should boil water...}.
\item Plastic bags block drains, \textbf{which causes floods}. \\
\emph{Attendu :} \cle{non-defining relative clause (proposition relative non restrictive / explicative)} with \emph{which} referring to the \cle{whole first clause (proposition principale entière)}; the \cle{comma (virgule)} is compulsory here.
\item We have \textbf{too few bins to sort waste}. \\
\emph{Attendu :} \emph{too + adjective (few) + to-infinitive}. Accept also \emph{Our bins are too few for us to sort waste}.
\item The expert warned that \textbf{if they didn't act then, the situation would get worse}. \\
\emph{Attendu :} full \cle{backshift (report du temps)} inside the conditional: \emph{don't} $\to$ \emph{didn't}, \emph{now} $\to$ \emph{then}, \emph{will} $\to$ \emph{would}. The first conditional becomes a ``reported'' conditional with past forms.
\item The dump is \textbf{so dangerous that children cannot play near it}. \\
\emph{Attendu :} \cle{So + adjective + that-clause (so + adjectif + que)}; the meaning (\cle{impossibility (impossibilité)}) must be preserved.
\item The mayor said that \textbf{burning waste in the open air was forbidden}. \\
\emph{Attendu :} Present Simple $\to$ Past Simple (\emph{is} $\to$ \emph{was}); the \cle{gerund subject (sujet gérondif)} \emph{burning waste} is kept.
\end{enumerate}
\emph{Points de vigilance (Bac) :} the meaning must not change; check the backshift of \textbf{all} time markers and verbs; in item 2, \emph{that} is impossible after the comma.
\end{corrige}

\begin{corrige}[Exercice 10 --- Type Bac : Cloze test]
\textbf{Barème indicatif : 1 mark par item. Total : 10 marks.}
\begin{enumerate}
\item \textbf{landfill} --- waste is buried ``at the outskirts of the city''.
\item \textbf{leachate} --- the ``toxic black liquid'' that flows out of a dump when it rains.
\item \textbf{sorting} --- families ``separate'' materials at home = \cle{sorting of waste (tri des déchets)}.
\item \textbf{biodegradable} --- food leftovers are materials that \cle{decompose naturally (se décomposent naturellement)}.
\item \textbf{compost} --- organic waste transformed into ``a natural fertiliser''.
\item \textbf{recycling} --- plastics, glass and paper sent to factories to become ``new products''.
\item \textbf{incineration} --- destruction of waste ``in modern ovens'' (burning).
\item \textbf{methane} --- the gas ``produced by decomposing waste'' that can produce electricity.
\item \textbf{circular economy} --- the model ``in which waste becomes a resource''.
\item \textbf{sustainability} --- guaranteeing the city's future ``for future generations''.
\end{enumerate}
\emph{Point de vigilance :} each word is used only once; no \cle{spelling mistake (faute d'orthographe)} is accepted at the Bac (copy the words exactly from the box).
\end{corrige}

\begin{corrige}[Exercice 11 --- Type Probatoire : Reading comprehension]
\textbf{Barème indicatif : 2 marks par question (1 mark for the correct information, 1 mark for a complete, correct sentence). Total : 10 marks.}
\begin{enumerate}
\item \textbf{(Literal / Explicite)} Two problems caused by the advance of the desert: \emph{the dry season becomes longer and the rains become rarer}; \emph{the fertile land slowly turns into sand} (desertification), which \cle{threatens (menace)} the lives of thousands of families. (Any two of these.)
\item \textbf{(Literal / Explicite)} Acacias and neems were chosen because \emph{they can survive with very little water}, and because \emph{their roots hold the soil together and prevent the wind from carrying it away}. (Also acceptable: their leaves provide \cle{shade (ombre)} for crops and \cle{food (nourriture)} for animals.)
\item \textbf{(Inferential / Vocabulaire)} In this context, ``nurseries'' means \emph{places where young trees (seedlings / plants) are grown and looked after before being planted in the ground} --- not places for babies.
\item \textbf{(Inferential / Inférentiel)} The author's attitude is \textbf{balanced (nuancée / équilibrée)}: the text presents \cle{encouraging results (résultats encourageants)} (more than three million trees planted, a 70\% survival rate, villages green again, children back at school) \emph{but} also insists on the \cle{remaining challenges (défis restants)} (climate change accelerating, funds arriving slowly, the need for the whole world to reduce emissions: ``We cannot win this battle alone... But we have the duty to try.'').
\item \textbf{(Synthesis / Synthèse)} Suitable title for paragraph 4: \emph{``Encouraging Results of the Initiative''} or \emph{``Visible Impact on Villages and Schools''}. (Any short title mentioning positive results is acceptable.)
\end{enumerate}
\emph{Points de vigilance (Probatoire) :} answer in \textbf{complete sentences (phrases complètes)}; for literal questions, the answer must come from the text (no personal opinion); for the attitude question, the \cle{justification with a textual element (justification par un élément du texte)} is compulsory to get full marks.
\end{corrige}

\begin{corrige}[Exercice 12 --- Type Bac : Writing (marking guide and model outlines)]
\textbf{Barème indicatif sur 20 :}
\begin{center}
\begin{tabularx}{\linewidth}{|l|X|c|}\hline
\rowcolor{popblueL} \textbf{Critère} & \textbf{Attendus} & \textbf{Points} \\ \hline
Content & Relevance to the topic; all required points covered (Topic A: problems + concrete solutions; Topic B: substances + consequences + prevention). & 6 \\ \hline
Organisation & Clear \cle{paragraphing (paragraphage)} (introduction, body, conclusion); \cle{linking words (mots de liaison)} (\emph{furthermore, however, consequently, in conclusion}). & 4 \\ \hline
Grammar & Correct use of the \cle{target structures (structures ciblées)}; \cle{tense consistency (cohérence des temps)}; varied sentences. & 5 \\ \hline
Vocabulary & Range and accuracy of unit vocabulary (\emph{landfill, biodegradable, compost, substance abuse, rehabilitation, sustainability...}); spelling. & 3 \\ \hline
Mechanics & Punctuation, capitalisation, word count respected (150--180 words). & 2 \\ \hline
\end{tabularx}
\end{center}

\textbf{Topic A --- model outline:}
\begin{itemize}
\item \emph{Introduction:} waste is a serious problem in our community (open dumps, blocked drains, floods).
\item \emph{Body:} solutions --- \cle{sorting at home (tri à la maison)}, \cle{composting organic waste (compostage des déchets organiques)}, \cle{recycling plastics (recyclage des plastiques)}; sensitisation campaigns. Required structures, e.g.: \emph{``If every family sorts its waste, our neighbourhood will be cleaner''} (\cle{first conditional}); \emph{``If I were the mayor, I would build a modern landfill''} (\cle{second conditional}); \emph{``The bins which are placed in every quarter must be emptied regularly''} (\cle{relative clause}).
\item \emph{Conclusion:} waste management is everyone's responsibility.
\end{itemize}

\textbf{Topic B --- model outline:}
\begin{itemize}
\item \emph{Introduction:} date, place and subject of the talk (health expert on substance abuse).
\item \emph{Body -- reported speech:} \emph{The expert explained that tobacco, alcohol and drugs were the most abused substances among youths. He warned that drug abuse could destroy the brain and lead to school dropout. He advised students to refuse cigarettes and alcohol, and he told them to report dealers to the authorities. He added that rehabilitation centres helped addicted young people to recover.}
\item \emph{Conclusion:} recommendation to organise more sensitisation talks in the school.
\end{itemize}

\emph{Points de vigilance (Bac) :} award full grammar marks only if the mandatory structures are used \textbf{correctly and naturally} (Topic A: one first conditional, one second conditional, one relative clause; Topic B: at least three \cle{reported-speech verbs (verbes de discours rapporté)} with correct \cle{backshift (report du temps)}). Penalise word counts outside 150--180 words and any structure copied without control (e.g. \emph{If I was you}, \emph{he said that he will}).
\end{corrige}

\newpage

\coursec{Fiche bilan}

\begin{popfigure}
\centering
\begin{tikzpicture}[
    mindmap,
    concept color=popblue,
    level 1 concept/.append style={level distance=4.5cm, sibling angle=360/7},
    level 2 concept/.append style={level distance=3cm},
    every node/.style={font=\small, text=white},
    concept/.append style={text width=2.8cm, align=center, inner sep=4pt},
    branch1/.style={concept color=poporange},
    branch2/.style={concept color=popgreen},
    branch3/.style={concept color=poppurple},
    branch4/.style={concept color=poppink},
    branch5/.style={concept color=popgold!80!popdark},
    branch6/.style={concept color=popblue!70!popdark},
    branch7/.style={concept color=popmuted}
]
\node[concept, text width=3.5cm] {Environment,\\ Well-being\\ \& Health}
    child[branch1] { node[concept, align=center]{1. Substance Abuse\\(Drugs, Tobacco,\\Alcohol)\\
    \textit{Gram: Relative Clauses}} }
    child[branch2] { node[concept, align=center]{2. Communicable\\ Diseases\\
    Prevention\\(School, Work,\\Community)\\
    \textit{Gram: Reported Speech}} }
    child[branch3] { node[concept, align=center]{3. Environmental\\ Issues\\
    Climate Change\&\ Global Warming\\
    \textit{Reading/Vocab}} }
    child[branch4] { node[concept, align=center]{4. Waste\\ Management\\
    Collection, Disposal,\\Recycling\\
    \textit{Vocab/Writing}} }
    child[branch5] { node[concept, align=center]{5. Conditionals\\
    1st (Real Future)\\
    2nd (Unreal Present)\\
    \textit{If + Present, Will / If + Past, Would}} }
    child[branch6] { node[concept, align=center]{6. Adjective\\ Structures\\
    \textit{How + Adj}\\
    \textit{Too... To}\\
    \textit{Superlatives}} }
    child[branch7] { node[concept, align=center]{7. Writing\&\ Integration\\
    Compositions\\
    Projects\&\ Remediation\\
    \textit{Action-Oriented}} };
\end{tikzpicture}
\legende{Mindmap of Lesson 03: Environment, Well-being \& Health (Terminale Technique)}
\end{popfigure}

\begin{bilan}

\courssub{Key Definitions \& Grammar Rules}

\begin{definition}[Relative Clauses (Propositions relatives)]
A subordinate clause that modifies a noun (antecedent). Introduced by a \cle{relative pronoun}.
\begin{itemize}
    \item \textbf{Who / That} : People (Sujet/Objet). \textit{Ex: The doctor \textbf{who} treats patients...}
    \item \textbf{Which / That} : Things/Animals (Sujet/Objet). \textit{Ex: The cigarette \textbf{which} causes cancer...}
    \item \textbf{Whose} : Possession (People/Things). \textit{Ex: The man \textbf{whose} son is sick...}
    \item \textbf{Where} : Place. \textit{Ex: The hospital \textbf{where} I was born...}
    \item \textbf{When} : Time. \textit{Ex: The day \textbf{when} we quit smoking...}
\end{itemize}
\textbf{Defining (Restrictive)}: No commas, essential info. \textbf{Non-defining}: Commas, extra info (cannot use \textit{that}).
\end{definition}

\begin{definition}[Direct \& Indirect (Reported) Speech]
\begin{itemize}
    \item \textbf{Direct}: Exact words. \textit{He said, "I am sick."}
    \item \textbf{Indirect}: Reported meaning, tense shift (\cle{backshift}), pronoun/time/place changes.
\end{itemize}
\end{definition}

\begin{center}
\begin{tabularx}{\linewidth}{|l|X|X|}\hline
\rowcolor{popblueL}
\textbf{Direct Tense} & \textbf{Reported Tense} & \textbf{Example} \\ \hline
Present Simple & Past Simple & "I \textbf{smoke}" $\rightarrow$ He said he \textbf{smoked}. \\ \hline
Present Continuous & Past Continuous & "I \textbf{am quitting}" $\rightarrow$ He said he \textbf{was quitting}. \\ \hline
Present Perfect & Past Perfect & "I \textbf{have stopped}" $\rightarrow$ He said he \textbf{had stopped}. \\ \hline
Past Simple & Past Perfect & "I \textbf{quit}" $\rightarrow$ He said he \textbf{had quit}. \\ \hline
Will & Would & "I \textbf{will try}" $\rightarrow$ He said he \textbf{would try}. \\ \hline
Can & Could & "I \textbf{can} resist" $\rightarrow$ He said he \textbf{could} resist. \\ \hline
Must & Had to & "I \textbf{must} go" $\rightarrow$ He said he \textbf{had to} go. \\ \hline
\end{tabularx}
\end{center}
\textbf{Modals not changing}: Could, Would, Should, Might, Ought to.
\textbf{Time/Place shifts}: Now$\rightarrow$Then, Today$\rightarrow$That day, Here$\rightarrow$There, Tomorrow$\rightarrow$The next day, Yesterday$\rightarrow$The previous day.

\begin{definition}[First \& Second Conditionals]
\begin{itemize}
    \item \textbf{First Conditional (Real Future)}: \textit{If + Present Simple, ... Will + Base Verb}. Possible/likely future.
    \item \textbf{Second Conditional (Unreal Present/Future)}: \textit{If + Past Simple, ... Would + Base Verb}. Hypothetical, unlikely, imaginary. \textit{If I were you, I would stop.}
\end{itemize}
\end{definition}

\begin{definition}[Advanced Adjective Structures]
\begin{itemize}
    \item \textbf{How + Adjective (+ Subject + Verb!)}: Exclamation. \textit{\textbf{How dangerous} smoking is!}
    \item \textbf{Too + Adjective + To + Verb}: Excess $\rightarrow$ Negative result (Impossible). \textit{He is \textbf{too weak to work}.}
    \item \textbf{Adjective + Enough + To + Verb}: Sufficiency. \textit{He is \textbf{strong enough to resist}.}
    \item \textbf{Superlative}: \textit{The most + long adj / The + short adj + -est}. \textit{The \textbf{worst} habit.}
\end{itemize}
\end{definition}

\courssub{Essential Vocabulary Clusters (Mots-clés)}

\begin{center}
\begin{tabularx}{\linewidth}{|l|X|}\hline
\rowcolor{popblueL}
\textbf{Theme} & \textbf{Key Terms (English / Français)} \\ \hline
Substance Abuse & Addiction (dépendance), Overdose (surdose), Withdrawal symptoms (symptômes de sevrage), Peer pressure (pression des pairs), Rehabilitation (rééducation/desintox), Sobriety (sobriété), Craving (envie irrésistible). \\ \hline
Communicable Diseases & Pathogen (pathogène), Transmission (transmission), Contagious (contagieux), Immunization/Vaccination (vaccination), Hygiene (hygiène), Sanitation (assainissement), Outbreak (épidémie), Quarantine (quarantaine), Vector (vecteur). \\ \hline
Environment / Climate & Greenhouse effect (effet de serre), Global warming (réchauffement climatique), Carbon footprint (empreinte carbone), Fossil fuels (combustibles fossiles), Renewable energy (énergies renouvelables), Deforestation (déforestation), Biodiversity loss (perte de biodiversité). \\ \hline
Waste Management & Landfill (décharge), Incineration (incinération), Recycling (recyclage), Composting (compostage), Biodegradable (biodégradable), Hazardous waste (déchets dangereux), Segregation/Sorting (tri sélectif), Circular economy (économie circulaire). \\ \hline
\end{tabularx}
\end{center}

\courssub{Express Methods (Méthodes Express)}

\begin{methode}[Writing a Relative Clause]
1. Identify the \textbf{antecedent} (noun to define). 2. Choose the correct \textbf{relative pronoun} (Who/Which/Whose/Where/When). 3. Place clause \textbf{immediately after} the antecedent. 4. Check: Defining (no comma) vs Non-defining (commas). 5. Ensure verb agrees with antecedent.
\end{methode}

\begin{methode}[Reporting Speech (Backshift)]
1. Change \textbf{reporting verb} tense (Say $\rightarrow$ Said, Tell $\rightarrow$ Told). 2. Apply \textbf{tense backshift} (Present $\rightarrow$ Past, Past $\rightarrow$ Past Perfect, Will $\rightarrow$ Would). 3. Change \textbf{pronouns} (I $\rightarrow$ He/She, My $\rightarrow$ His/Her). 4. Adjust \textbf{time/place markers} (Now $\rightarrow$ Then, Here $\rightarrow$ There). 5. Remove quotation marks, add \textit{that} (optional).
\end{methode}

\begin{methode}[Choosing Conditional Type]
\textbf{Step 1}: Is the situation \textbf{Real/Possible} (Future)? $\rightarrow$ \textbf{1st Conditional} (If + Present, Will).\\
\textbf{Step 2}: Is the situation \textbf{Imaginary/Unlikely/Advice} (Present/Future)? $\rightarrow$ \textbf{2nd Conditional} (If + Past, Would).\\
\textbf{Step 3}: \textit{If I were you} (Advice) always uses 2nd Conditional.
\end{methode}

\begin{methode}[Writing a Composition (Bac Standard)]
1. \textbf{Analyse} the topic (Type: Argumentative, Descriptive, Narrative, Letter). 2. \textbf{Brainstorm} vocab \& arguments (Mindmap). 3. \textbf{Outline}: Introduction (Hook, Thesis), Body (2-3 paragraphs: Topic sentence, Evidence, Link), Conclusion (Restate thesis, Call to action). 4. \textbf{Draft} using connectors (Furthermore, However, Consequently). 5. \textbf{Proofread}: Grammar (Tenses, Conditionals, Relatives), Spelling, Word count ($\pm$10\%).
\end{methode}

\end{bilan}

\begin{aretenir}[Checklist avant le Bac \textendash{} Lesson 03]
\begin{itemize}
    \item $\square$ I can define \textbf{Relative Clauses} and choose the correct pronoun (Who/Which/Whose/Where/When).
    \item $\square$ I distinguish \textbf{Defining} vs \textbf{Non-defining} clauses (comma usage).
    \item $\square$ I master \textbf{Backshift rules} for Reported Speech (Tenses, Modals, Time/Place).
    \item $\square$ I correctly form \textbf{First Conditional} (Real future: If + Present, Will).
    \item $\square$ I correctly form \textbf{Second Conditional} (Unreal: If + Past, Would) + \textit{If I were you}.
    \item $\square$ I use \textbf{How + Adjective} for exclamations (\textit{How harmful...!}).
    \item $\square$ I distinguish \textbf{Too...To} (Impossible) vs \textbf{Adj+Enough+To} (Possible).
    \item $\square$ I form \textbf{Superlatives} correctly (Short/Long adjectives, Irregulars: Bad$\rightarrow$Worst).
    \item $\square$ I know key vocabulary for \textbf{Substance Abuse}, \textbf{Communicable Diseases}, \textbf{Climate Change}, \textbf{Waste Management}.
    \item $\square$ I can write a structured \textbf{Composition} (Intro, Body, Conclusion) on environment/health topics.
    \item $\square$ I can extract specific info from a \textbf{Listening} text on disease prevention.
    \item $\square$ I can participate in a \textbf{Speaking} exchange on drug/alcohol abuse risks.
\end{itemize}
\end{aretenir}

\findecours

\end{document}
